\documentclass[11pt]{article}
\usepackage[utf8]{inputenc}
\usepackage[T1]{fontenc}
\usepackage[italian]{babel}
\usepackage{geometry,amsmath,amssymb,mathtools,amsthm,bm}
\usepackage{graphicx,caption,subcaption,booktabs}
\usepackage{hyperref}
\usepackage{tikz}
\usepackage{pgfplots}
\pgfplotsset{compat=1.18}
\usepackage{siunitx,enumitem,listings}
\usepackage{longtable,array}
\geometry{margin=2.3cm}
\lstset{basicstyle=\ttfamily\small,breaklines=true,columns=fullflexible,keepspaces=true}
\title{Spiegazione estesa delle slide di Basi di Dati 2\\(Lezioni 00, 01, 02)}
\author{Fonte: Lezione 00 - Presentazione del corso.pdf\\Fonte: Lezione 01 - Dipendenze funzionali e Normalizzazione di DB Relazionali.pdf\\Fonte: Lezione 02 - Algoritmi per la progettazione di DB Relazionali e ulteriori dipendenze.pdf}
\date{14 marzo 2026}
\begin{document}
\begin{titlepage}\maketitle\end{titlepage}
\begin{abstract}
Questo documento fornisce una spiegazione estesa delle slide delle Lezioni 00, 01 e 02 del corso di Basi di Dati 2. Sono riportate trascrizioni per ciascuna slide e integrazioni didattiche marcate [AGGIUNTA]. Le parti non leggibili o ambigue sono marcate [INCERTO]. I temi coprono presentazione del corso, qualità degli schemi relazionali, dipendenze funzionali, normalizzazione, chiusure e algoritmi di progettazione. Sono inclusi grafici vettoriali TikZ/PGFPlots, tabelle con dati sintetici dichiarati, mappa slide->sezione, elenco simboli e riepilogo contenuti aggiunti. L'obiettivo è un supporto di studio trasparente e compilabile con pdflatex.
\end{abstract}
\tableofcontents
\listoffigures
\listoftables
\section{Introduzione}
[AGGIUNTA] Questo testo organizza e collega i contenuti presenti nelle slide originali per facilitare lo studio progressivo, mantenendo separata la trascrizione dalla spiegazione.
\section{Grafici esplicativi aggiunti}
[AGGIUNTA] I grafici seguenti sono stati aggiunti per visualizzare concetti qualitativi presenti nelle slide. Quando i dati numerici non erano forniti, sono stati usati dati sintetici coerenti e dichiarati.
\begin{table}[ht]
\centering
\caption{Dati sintetici per la curva di crescita combinatoria (numero di sottoinsiemi $2^n$)}
\label{tab:dati-sintetici-combinatoria}
\begin{tabular}{@{}rrrrrr@{}}
\toprule
$n$ & 2 & 4 & 6 & 8 & 10 \\
\midrule
$2^n$ & 4 & 16 & 64 & 256 & 1024 \\
\bottomrule
\end{tabular}
\end{table}
\begin{figure}[ht]
\centering
\begin{tikzpicture}
\begin{axis}[
  xlabel={$n$ attributi}, ylabel={Numero di sottoinsiemi $2^n$},
  grid=major, width=0.8\textwidth, height=6cm,
  legend pos=outer north east
]
\addplot[smooth, blue, thick, mark=*] table[row sep=crcr] {
  2 4
  4 16
  6 64
  8 256
  10 1024
};
\addlegendentry{Crescita combinatoria}
\end{axis}
\end{tikzpicture}
\caption{[AGGIUNTA] Crescita combinatoria dello spazio di ricerca in problemi di progettazione. Dati sintetici coerenti con $2^n$.}
\label{fig:crescita-combinatoria}
\end{figure}
\begin{figure}[ht]
\centering
\begin{tikzpicture}
\begin{axis}[
  ybar, bar width=18pt,
  symbolic x coords={Top-Down,Bottom-Up,Ibrido},
  xtick=data,
  ylabel={Punteggio qualitativo (0--10)},
  width=0.8\textwidth,height=6cm,
  ymin=0,ymax=10,
  legend pos=outer north east,
  grid=major
]
\addplot[fill=blue!30] coordinates {(Top-Down,8) (Bottom-Up,6) (Ibrido,9)};
\addplot[fill=red!30] coordinates {(Top-Down,5) (Bottom-Up,8) (Ibrido,7)};
\legend{Controllo semantico,Flessibilità incrementale}
\end{axis}
\end{tikzpicture}
\caption{[AGGIUNTA] Confronto qualitativo tra strategie progettuali. Dati sintetici per visualizzare il trade-off discusso nelle slide.}
\label{fig:confronto-strategie}
\end{figure}
\begin{figure}[ht]
\centering
\begin{tikzpicture}
\begin{axis}[
  xlabel={Schema candidato}, ylabel={Indice qualitativo normalizzato},
  grid=major, width=0.8\textwidth, height=6cm,
  legend pos=outer north east,
  ymin=0,ymax=1
]
\addplot[smooth, green!60!black, thick, mark=triangle*] table[row sep=crcr] {
  1 0.40
  2 0.58
  3 0.72
  4 0.83
  5 0.91
};
\addlegendentry{Miglioramento dopo normalizzazione}
\end{axis}
\end{tikzpicture}
\caption{[AGGIUNTA] Esempio qualitativo di miglioramento della qualità dello schema attraverso normalizzazione progressiva. Dati sintetici.}
\label{fig:miglioramento-normalizzazione}
\end{figure}
\section{CORSO DI LAUREA MAGISTRALE --- Slide 1}
\subsection*{Estratto originale dalla slide 1}
\begin{lstlisting}
CORSO DI LAUREA MAGISTRALE

  BASI DI DATI 2


  PRESENTAZIONE DEL CORSO



                       Prof.ssa G. Tortora
\end{lstlisting}
\subsection*{Spiegazione estesa}
[AGGIUNTA] Questa spiegazione collega il contenuto testuale della slide al contesto del corso. La slide enfatizza i concetti mostrati nell'estratto e il loro ruolo nella progettazione di basi di dati relazionali, con particolare attenzione a correttezza semantica, qualità dello schema e compromessi progettuali.
[AGGIUNTA] Fonte integrazione: riformulazione didattica del testo presente nella slide 1 del file Lezione 00 - Presentazione del corso.pdf.
\subsection*{Derivazioni}
[AGGIUNTA] In questa slide non sono state rilevate equazioni esplicite nel testo estratto; non è quindi richiesta una derivazione algebrica puntuale.
\subsection*{Esempio numerico / esercizio}
[AGGIUNTA] Esempio: dato uno schema R(A,B,C) con dipendenze funzionali A->B e B->C, mostrare la chiusura A^+ e discutere se A è superchiave. Soluzione sintetica: A^+={A,B,C}, quindi A determina tutti gli attributi di R ed è superchiave.
\subsection*{Note su assunzioni o integrazioni}
[ASSUNZIONE] L'estrazione OCR può perdere formattazione o simboli della slide originale.
[AGGIUNTA] Le spiegazioni sono estensioni didattiche per rendere espliciti passaggi impliciti nel materiale proiettato.
\paragraph{Riferimento fonte} File: Lezione 00 - Presentazione del corso.pdf. Numero slide: 1.
\section{Caratteristiche generali --- Slide 2}
\subsection*{Estratto originale dalla slide 2}
\begin{lstlisting}
Caratteristiche generali

-    Nome dell'insegnamento: BASI DI DATI 2 (BD2)

-  Crediti formativi: 9
-  Durata: 72 ore

    48 ore di lezione

    24 ore di esercitazioni in laboratorio

-  Periodo: 27 febbraio - 01 giugno
\end{lstlisting}
\subsection*{Spiegazione estesa}
[AGGIUNTA] Questa spiegazione collega il contenuto testuale della slide al contesto del corso. La slide enfatizza i concetti mostrati nell'estratto e il loro ruolo nella progettazione di basi di dati relazionali, con particolare attenzione a correttezza semantica, qualità dello schema e compromessi progettuali.
[AGGIUNTA] Fonte integrazione: riformulazione didattica del testo presente nella slide 2 del file Lezione 00 - Presentazione del corso.pdf.
\subsection*{Derivazioni}
[AGGIUNTA] In questa slide non sono state rilevate equazioni esplicite nel testo estratto; non è quindi richiesta una derivazione algebrica puntuale.
\subsection*{Esempio numerico / esercizio}
[AGGIUNTA] Esempio: dato uno schema R(A,B,C) con dipendenze funzionali A->B e B->C, mostrare la chiusura A^+ e discutere se A è superchiave. Soluzione sintetica: A^+={A,B,C}, quindi A determina tutti gli attributi di R ed è superchiave.
\subsection*{Note su assunzioni o integrazioni}
[ASSUNZIONE] L'estrazione OCR può perdere formattazione o simboli della slide originale.
[AGGIUNTA] Le spiegazioni sono estensioni didattiche per rendere espliciti passaggi impliciti nel materiale proiettato.
\paragraph{Riferimento fonte} File: Lezione 00 - Presentazione del corso.pdf. Numero slide: 2.
\section{Calendario del corso --- Slide 3}
\subsection*{Estratto originale dalla slide 3}
\begin{lstlisting}
Calendario del corso

-  Martedì 14.30 - 17.30 (F4)
-  Giovedì 09.00 - 12.00 (F4)



-    Orario ricevimento:
      Martedì 11.00 - 14.00



-    Studio:
      Edificio F2. 1 gradi  piano, stanza 81
\end{lstlisting}
\subsection*{Spiegazione estesa}
[AGGIUNTA] Questa spiegazione collega il contenuto testuale della slide al contesto del corso. La slide enfatizza i concetti mostrati nell'estratto e il loro ruolo nella progettazione di basi di dati relazionali, con particolare attenzione a correttezza semantica, qualità dello schema e compromessi progettuali.
[AGGIUNTA] Fonte integrazione: riformulazione didattica del testo presente nella slide 3 del file Lezione 00 - Presentazione del corso.pdf.
\subsection*{Derivazioni}
[AGGIUNTA] In questa slide non sono state rilevate equazioni esplicite nel testo estratto; non è quindi richiesta una derivazione algebrica puntuale.
\subsection*{Esempio numerico / esercizio}
[AGGIUNTA] Esempio: dato uno schema R(A,B,C) con dipendenze funzionali A->B e B->C, mostrare la chiusura A^+ e discutere se A è superchiave. Soluzione sintetica: A^+={A,B,C}, quindi A determina tutti gli attributi di R ed è superchiave.
\subsection*{Note su assunzioni o integrazioni}
[ASSUNZIONE] L'estrazione OCR può perdere formattazione o simboli della slide originale.
[AGGIUNTA] Le spiegazioni sono estensioni didattiche per rendere espliciti passaggi impliciti nel materiale proiettato.
\paragraph{Riferimento fonte} File: Lezione 00 - Presentazione del corso.pdf. Numero slide: 3.
\section{Laboratorio --- Slide 4}
\subsection*{Estratto originale dalla slide 4}
\begin{lstlisting}
Laboratorio

-    Il laboratorio è parte integrante del corso!




           -    Frequentare le esercitazioni guidate è
               molto importante:
                 Alcune problematiche si capiscono molto
                 meglio mettendo in pratica i concetti teorici
                 appresi a lezione
                 Acquisire manualità nell'uso degli strumenti
                 software è fondamentale
\end{lstlisting}
\subsection*{Spiegazione estesa}
[AGGIUNTA] Questa spiegazione collega il contenuto testuale della slide al contesto del corso. La slide enfatizza i concetti mostrati nell'estratto e il loro ruolo nella progettazione di basi di dati relazionali, con particolare attenzione a correttezza semantica, qualità dello schema e compromessi progettuali.
[AGGIUNTA] Fonte integrazione: riformulazione didattica del testo presente nella slide 4 del file Lezione 00 - Presentazione del corso.pdf.
\subsection*{Derivazioni}
[AGGIUNTA] In questa slide non sono state rilevate equazioni esplicite nel testo estratto; non è quindi richiesta una derivazione algebrica puntuale.
\subsection*{Esempio numerico / esercizio}
[AGGIUNTA] Esempio: dato uno schema R(A,B,C) con dipendenze funzionali A->B e B->C, mostrare la chiusura A^+ e discutere se A è superchiave. Soluzione sintetica: A^+={A,B,C}, quindi A determina tutti gli attributi di R ed è superchiave.
\subsection*{Note su assunzioni o integrazioni}
[ASSUNZIONE] L'estrazione OCR può perdere formattazione o simboli della slide originale.
[AGGIUNTA] Le spiegazioni sono estensioni didattiche per rendere espliciti passaggi impliciti nel materiale proiettato.
\paragraph{Riferimento fonte} File: Lezione 00 - Presentazione del corso.pdf. Numero slide: 4.
\section{Descrizione del corso --- Slide 5}
\subsection*{Estratto originale dalla slide 5}
\begin{lstlisting}
Descrizione del corso
-    Il corso intende approfondire alcuni aspetti teorici ed affrontare lo
    studio di argomenti avanzati inerenti le Basi di Dati e le tecnologie di
    supporto.
-    In particolare il corso tratterà:
      Approfondimenti sulla progettazione logica: regole di inferenza e
      forme normali superiori alla terza;
      Design fisico dei database: organizzazione dei file, indici primari e
      secondari;
      Gestione delle transazioni, tecniche di controllo della concorrenza,
      meccanismi di gestione dei lock, tecniche di recovery;
      Architetture distribuite: architetture Client Server, basi di dati e World
      Wide Web, Basi di dati distribuite, Basi di Dati semi-strutturate e per
      XML;
      Basi di dati semantiche, SPARQL e Linked Open Data;
      Gestione della conoscenza e Information Retrieval;
      Tecnologie emergenti: Data Warehouse, Sistemi di Business
      Intelligence, Basi di Dati multimediali, Big Data, database NoSQL.
\end{lstlisting}
\subsection*{Spiegazione estesa}
[AGGIUNTA] Questa spiegazione collega il contenuto testuale della slide al contesto del corso. La slide enfatizza i concetti mostrati nell'estratto e il loro ruolo nella progettazione di basi di dati relazionali, con particolare attenzione a correttezza semantica, qualità dello schema e compromessi progettuali.
[AGGIUNTA] Fonte integrazione: riformulazione didattica del testo presente nella slide 5 del file Lezione 00 - Presentazione del corso.pdf.
\subsection*{Derivazioni}
[AGGIUNTA] In questa slide non sono state rilevate equazioni esplicite nel testo estratto; non è quindi richiesta una derivazione algebrica puntuale.
\subsection*{Esempio numerico / esercizio}
[AGGIUNTA] Esempio: dato uno schema R(A,B,C) con dipendenze funzionali A->B e B->C, mostrare la chiusura A^+ e discutere se A è superchiave. Soluzione sintetica: A^+={A,B,C}, quindi A determina tutti gli attributi di R ed è superchiave.
\subsection*{Note su assunzioni o integrazioni}
[ASSUNZIONE] L'estrazione OCR può perdere formattazione o simboli della slide originale.
[AGGIUNTA] Le spiegazioni sono estensioni didattiche per rendere espliciti passaggi impliciti nel materiale proiettato.
\paragraph{Riferimento fonte} File: Lezione 00 - Presentazione del corso.pdf. Numero slide: 5.
\section{Obiettivi --- Slide 6}
\subsection*{Estratto originale dalla slide 6}
\begin{lstlisting}
Obiettivi

-    Il corso mira a sviluppare nello studente le
    seguenti capacità:
      Saper progettare Sistemi Informativi basati sulle
      tecnologie trattate.
      Saper progettare Data Warehouse e Sistemi di
      Business Intelligence.
      Capacità di analisi critica dei modelli e delle
      tecnologie di database avanzati.
      Saper progettare e realizzare tecnologie di
      supporto alla gestione di database avanzati e
      Data Analytics
\end{lstlisting}
\subsection*{Spiegazione estesa}
[AGGIUNTA] Questa spiegazione collega il contenuto testuale della slide al contesto del corso. La slide enfatizza i concetti mostrati nell'estratto e il loro ruolo nella progettazione di basi di dati relazionali, con particolare attenzione a correttezza semantica, qualità dello schema e compromessi progettuali.
[AGGIUNTA] Fonte integrazione: riformulazione didattica del testo presente nella slide 6 del file Lezione 00 - Presentazione del corso.pdf.
\subsection*{Derivazioni}
[AGGIUNTA] In questa slide non sono state rilevate equazioni esplicite nel testo estratto; non è quindi richiesta una derivazione algebrica puntuale.
\subsection*{Esempio numerico / esercizio}
[AGGIUNTA] Esempio: dato uno schema R(A,B,C) con dipendenze funzionali A->B e B->C, mostrare la chiusura A^+ e discutere se A è superchiave. Soluzione sintetica: A^+={A,B,C}, quindi A determina tutti gli attributi di R ed è superchiave.
\subsection*{Note su assunzioni o integrazioni}
[ASSUNZIONE] L'estrazione OCR può perdere formattazione o simboli della slide originale.
[AGGIUNTA] Le spiegazioni sono estensioni didattiche per rendere espliciti passaggi impliciti nel materiale proiettato.
\paragraph{Riferimento fonte} File: Lezione 00 - Presentazione del corso.pdf. Numero slide: 6.
\section{Prerequisiti --- Slide 7}
\subsection*{Estratto originale dalla slide 7}
\begin{lstlisting}
Prerequisiti

-  Conoscenze di base sui metodi e tecniche di
  progettazione di una base di dati;
-  Conoscenze di SQL;

-  Conoscenze di XML;

-  Fondamenti di sistemi distribuiti;

-  Paradigma di programmazione ad oggetti;

-  Il linguaggio di programmazione JAVA.
\end{lstlisting}
\subsection*{Spiegazione estesa}
[AGGIUNTA] Questa spiegazione collega il contenuto testuale della slide al contesto del corso. La slide enfatizza i concetti mostrati nell'estratto e il loro ruolo nella progettazione di basi di dati relazionali, con particolare attenzione a correttezza semantica, qualità dello schema e compromessi progettuali.
[AGGIUNTA] Fonte integrazione: riformulazione didattica del testo presente nella slide 7 del file Lezione 00 - Presentazione del corso.pdf.
\subsection*{Derivazioni}
[AGGIUNTA] In questa slide non sono state rilevate equazioni esplicite nel testo estratto; non è quindi richiesta una derivazione algebrica puntuale.
\subsection*{Esempio numerico / esercizio}
[AGGIUNTA] Esempio: dato uno schema R(A,B,C) con dipendenze funzionali A->B e B->C, mostrare la chiusura A^+ e discutere se A è superchiave. Soluzione sintetica: A^+={A,B,C}, quindi A determina tutti gli attributi di R ed è superchiave.
\subsection*{Note su assunzioni o integrazioni}
[ASSUNZIONE] L'estrazione OCR può perdere formattazione o simboli della slide originale.
[AGGIUNTA] Le spiegazioni sono estensioni didattiche per rendere espliciti passaggi impliciti nel materiale proiettato.
\paragraph{Riferimento fonte} File: Lezione 00 - Presentazione del corso.pdf. Numero slide: 7.
\section{Link --- Slide 8}
\subsection*{Estratto originale dalla slide 8}
\begin{lstlisting}
Link

-    Homepage del corso sulla piattaforma e-learning

-    Contiene:
      Copia delle slide in formato pdf
      Link utili
      Esercitazioni
      Risorse
      Prove d'esame
      Risultati
      Documentazione progetti
      Avvisi
      ...
\end{lstlisting}
\subsection*{Spiegazione estesa}
[AGGIUNTA] Questa spiegazione collega il contenuto testuale della slide al contesto del corso. La slide enfatizza i concetti mostrati nell'estratto e il loro ruolo nella progettazione di basi di dati relazionali, con particolare attenzione a correttezza semantica, qualità dello schema e compromessi progettuali.
[AGGIUNTA] Fonte integrazione: riformulazione didattica del testo presente nella slide 8 del file Lezione 00 - Presentazione del corso.pdf.
\subsection*{Derivazioni}
[AGGIUNTA] In questa slide non sono state rilevate equazioni esplicite nel testo estratto; non è quindi richiesta una derivazione algebrica puntuale.
\subsection*{Esempio numerico / esercizio}
[AGGIUNTA] Esempio: dato uno schema R(A,B,C) con dipendenze funzionali A->B e B->C, mostrare la chiusura A^+ e discutere se A è superchiave. Soluzione sintetica: A^+={A,B,C}, quindi A determina tutti gli attributi di R ed è superchiave.
\subsection*{Note su assunzioni o integrazioni}
[ASSUNZIONE] L'estrazione OCR può perdere formattazione o simboli della slide originale.
[AGGIUNTA] Le spiegazioni sono estensioni didattiche per rendere espliciti passaggi impliciti nel materiale proiettato.
\paragraph{Riferimento fonte} File: Lezione 00 - Presentazione del corso.pdf. Numero slide: 8.
\section{Modalità d'esame --- Slide 9}
\subsection*{Estratto originale dalla slide 9}
\begin{lstlisting}
Modalità d'esame

-    Una prova intercorso [inizio aprile] (prova scritta)

-    Progetto (gruppo composto da 2-3 componenti)

-    L'esame consiste di una prova scritta e di un colloquio
    orale (in cui si discuterà il progetto)

-    Il superamento delle prove intercorso dispensa dalla prova
    scritta
-    Una volta superata la prova intercorso o la prova scritta è
    possibile sostenere l'orale (con discussione del progetto) in
    qualsiasi appello
\end{lstlisting}
\subsection*{Spiegazione estesa}
[AGGIUNTA] Questa spiegazione collega il contenuto testuale della slide al contesto del corso. La slide enfatizza i concetti mostrati nell'estratto e il loro ruolo nella progettazione di basi di dati relazionali, con particolare attenzione a correttezza semantica, qualità dello schema e compromessi progettuali.
[AGGIUNTA] Fonte integrazione: riformulazione didattica del testo presente nella slide 9 del file Lezione 00 - Presentazione del corso.pdf.
\subsection*{Derivazioni}
[AGGIUNTA] In questa slide non sono state rilevate equazioni esplicite nel testo estratto; non è quindi richiesta una derivazione algebrica puntuale.
\subsection*{Esempio numerico / esercizio}
[AGGIUNTA] Esempio: dato uno schema R(A,B,C) con dipendenze funzionali A->B e B->C, mostrare la chiusura A^+ e discutere se A è superchiave. Soluzione sintetica: A^+={A,B,C}, quindi A determina tutti gli attributi di R ed è superchiave.
\subsection*{Note su assunzioni o integrazioni}
[ASSUNZIONE] L'estrazione OCR può perdere formattazione o simboli della slide originale.
[AGGIUNTA] Le spiegazioni sono estensioni didattiche per rendere espliciti passaggi impliciti nel materiale proiettato.
\paragraph{Riferimento fonte} File: Lezione 00 - Presentazione del corso.pdf. Numero slide: 9.
\section{Testi consigliati --- Slide 10}
\subsection*{Estratto originale dalla slide 10}
\begin{lstlisting}
Testi consigliati
-    Materiale a supporto fornito durante il corso: slide, articoli scientifici, link.

-    R.A Elmasri, S.B. Navathe,
    "Sistemi di Basi di Dati - Fondamenti",
    6a edizione, Addison Wesley, 2011.

-    R.A Elmasri, S.B. Navathe,
    "Sistemi di Basi di Dati - Complementi",
    4a edizione, Pearson - Addison Wesley, 2005.

-    M. Golfarelli, S. Rizzi,
    "Data Warehouse - Teoria e pratica della progettazione",
    The McGraw-Hill Companies, 2006.

-    P. Atzeni, S. Ceri, S. Paraboschi, R. Torlone,
    "Basi di dati - Architetture e Linee di Evoluzione",
    2a edizione, The McGraw-Hill Companies, 2007.
\end{lstlisting}
\subsection*{Spiegazione estesa}
[AGGIUNTA] Questa spiegazione collega il contenuto testuale della slide al contesto del corso. La slide enfatizza i concetti mostrati nell'estratto e il loro ruolo nella progettazione di basi di dati relazionali, con particolare attenzione a correttezza semantica, qualità dello schema e compromessi progettuali.
[AGGIUNTA] Fonte integrazione: riformulazione didattica del testo presente nella slide 10 del file Lezione 00 - Presentazione del corso.pdf.
\subsection*{Derivazioni}
[AGGIUNTA] In questa slide non sono state rilevate equazioni esplicite nel testo estratto; non è quindi richiesta una derivazione algebrica puntuale.
\subsection*{Esempio numerico / esercizio}
[AGGIUNTA] Esempio: dato uno schema R(A,B,C) con dipendenze funzionali A->B e B->C, mostrare la chiusura A^+ e discutere se A è superchiave. Soluzione sintetica: A^+={A,B,C}, quindi A determina tutti gli attributi di R ed è superchiave.
\subsection*{Note su assunzioni o integrazioni}
[ASSUNZIONE] L'estrazione OCR può perdere formattazione o simboli della slide originale.
[AGGIUNTA] Le spiegazioni sono estensioni didattiche per rendere espliciti passaggi impliciti nel materiale proiettato.
\paragraph{Riferimento fonte} File: Lezione 00 - Presentazione del corso.pdf. Numero slide: 10.
\section{BASI DI DATI 2 --- Slide 1}
\subsection*{Estratto originale dalla slide 1}
\begin{lstlisting}
BASI DI DATI 2


DIPENDENZE FUNZIONALI E
NORMALIZZAZIONE DI DB
RELAZIONALI

               Prof.ssa G. Tortora
\end{lstlisting}
\subsection*{Spiegazione estesa}
[AGGIUNTA] Questa spiegazione collega il contenuto testuale della slide al contesto del corso. La slide enfatizza i concetti mostrati nell'estratto e il loro ruolo nella progettazione di basi di dati relazionali, con particolare attenzione a correttezza semantica, qualità dello schema e compromessi progettuali.
[AGGIUNTA] Fonte integrazione: riformulazione didattica del testo presente nella slide 1 del file Lezione 01 - Dipendenze funzionali e Normalizzazione di DB Relazionali.pdf.
\subsection*{Derivazioni}
[AGGIUNTA] In questa slide non sono state rilevate equazioni esplicite nel testo estratto; non è quindi richiesta una derivazione algebrica puntuale.
\subsection*{Esempio numerico / esercizio}
[AGGIUNTA] Esempio: dato uno schema R(A,B,C) con dipendenze funzionali A->B e B->C, mostrare la chiusura A^+ e discutere se A è superchiave. Soluzione sintetica: A^+={A,B,C}, quindi A determina tutti gli attributi di R ed è superchiave.
\subsection*{Note su assunzioni o integrazioni}
[ASSUNZIONE] L'estrazione OCR può perdere formattazione o simboli della slide originale.
[AGGIUNTA] Le spiegazioni sono estensioni didattiche per rendere espliciti passaggi impliciti nel materiale proiettato.
\paragraph{Riferimento fonte} File: Lezione 01 - Dipendenze funzionali e Normalizzazione di DB Relazionali.pdf. Numero slide: 1.
\section{Le fasi di progettazione di un Database --- Slide 2}
\subsection*{Estratto originale dalla slide 2}
\begin{lstlisting}
Le fasi di progettazione di un Database
   Indipendente dal DBMS                                      Miniworld



                                                                                                   1
                                                   RACCOLTA ED ANALISI DEI REQUISITI
                                                     RACCOLTA ED ANALISI
                                                     DEI REQUISITI


                         Requisiti Funzionali                                                   Requisiti sui Dati


                                                                                                                         2
                                                                                        DISEGNO CONCETTUALE
                        ANALISI FUNZIONALE                                                DISEGNO CONCETTUALE

                                                                                         Schema Concettuale
                  Specifia transazioni di alto livello                             (Usando Data Model di alto livello)

                                                                                       (DATA MODEL
                                                                                       DISEGNO LOGICO
                                                                                                   MAPPING)


   Specifico del DBMS
                                                                                             DISEGNO LOGICO
                                                                                          (DATA MODEL MAPPING)           3
                     DISEGNO PROGRAMMI                                                 Schema Logico (Concettuale)
                     APPLICATIVI                                                       Usando il data model del DBMS


                                                                                                                     4
                                                                                             DISEGNO FISICO
                        IMPLEMENTAZIONE                                                        DISEGNO FISICO
                        TRANSAZIONI

                             Programmi                                                          Schema Interno
                             Applicativi
\end{lstlisting}
\subsection*{Spiegazione estesa}
[AGGIUNTA] Questa spiegazione collega il contenuto testuale della slide al contesto del corso. La slide enfatizza i concetti mostrati nell'estratto e il loro ruolo nella progettazione di basi di dati relazionali, con particolare attenzione a correttezza semantica, qualità dello schema e compromessi progettuali.
[AGGIUNTA] Fonte integrazione: riformulazione didattica del testo presente nella slide 2 del file Lezione 01 - Dipendenze funzionali e Normalizzazione di DB Relazionali.pdf.
\subsection*{Derivazioni}
[AGGIUNTA] In questa slide non sono state rilevate equazioni esplicite nel testo estratto; non è quindi richiesta una derivazione algebrica puntuale.
\subsection*{Esempio numerico / esercizio}
[AGGIUNTA] Esempio: dato uno schema R(A,B,C) con dipendenze funzionali A->B e B->C, mostrare la chiusura A^+ e discutere se A è superchiave. Soluzione sintetica: A^+={A,B,C}, quindi A determina tutti gli attributi di R ed è superchiave.
\subsection*{Note su assunzioni o integrazioni}
[ASSUNZIONE] L'estrazione OCR può perdere formattazione o simboli della slide originale.
[AGGIUNTA] Le spiegazioni sono estensioni didattiche per rendere espliciti passaggi impliciti nel materiale proiettato.
\paragraph{Riferimento fonte} File: Lezione 01 - Dipendenze funzionali e Normalizzazione di DB Relazionali.pdf. Numero slide: 2.
\section{Qualità di schemi relazionali --- Slide 3}
\subsection*{Estratto originale dalla slide 3}
\begin{lstlisting}
Qualità di schemi relazionali
-    Non sempre dalla fase di progettazione si ottiene
    uno schema privo di difetti:
      Possono sorgere dei problemi nella fase di mapping da
      ER a Relazionale.
      Si può ottenere uno schema non ottimale progettando
      direttamente uno schema logico saltando la fase di
      analisi concettuale.


-    Tali difetti possono portare ad anomalie nella
    base di dati.
\end{lstlisting}
\subsection*{Spiegazione estesa}
[AGGIUNTA] Questa spiegazione collega il contenuto testuale della slide al contesto del corso. La slide enfatizza i concetti mostrati nell'estratto e il loro ruolo nella progettazione di basi di dati relazionali, con particolare attenzione a correttezza semantica, qualità dello schema e compromessi progettuali.
[AGGIUNTA] Fonte integrazione: riformulazione didattica del testo presente nella slide 3 del file Lezione 01 - Dipendenze funzionali e Normalizzazione di DB Relazionali.pdf.
\subsection*{Derivazioni}
[AGGIUNTA] In questa slide non sono state rilevate equazioni esplicite nel testo estratto; non è quindi richiesta una derivazione algebrica puntuale.
\subsection*{Esempio numerico / esercizio}
[AGGIUNTA] Esempio: dato uno schema R(A,B,C) con dipendenze funzionali A->B e B->C, mostrare la chiusura A^+ e discutere se A è superchiave. Soluzione sintetica: A^+={A,B,C}, quindi A determina tutti gli attributi di R ed è superchiave.
\subsection*{Note su assunzioni o integrazioni}
[ASSUNZIONE] L'estrazione OCR può perdere formattazione o simboli della slide originale.
[AGGIUNTA] Le spiegazioni sono estensioni didattiche per rendere espliciti passaggi impliciti nel materiale proiettato.
\paragraph{Riferimento fonte} File: Lezione 01 - Dipendenze funzionali e Normalizzazione di DB Relazionali.pdf. Numero slide: 3.
\section{Qualità di schemi relazionali (2) --- Slide 4}
\subsection*{Estratto originale dalla slide 4}
\begin{lstlisting}
Qualità di schemi relazionali (2)
-  Come misurare la "bontà" di uno schema di
  relazione?
-  La discussione è a due livelli:
      livello logico
        riferito al significato dello schema di relazione e degli attributi;
        la bontà, a questo livello, aiuta l'utente a capire chiaramente il
        significato dei dati ed a formulare query correttamente;

      livello di memorizzazione e manipolazione
        Riguarda come le tuple sono memorizzate e aggiornate (solo
        per le relazioni base);
\end{lstlisting}
\subsection*{Spiegazione estesa}
[AGGIUNTA] Questa spiegazione collega il contenuto testuale della slide al contesto del corso. La slide enfatizza i concetti mostrati nell'estratto e il loro ruolo nella progettazione di basi di dati relazionali, con particolare attenzione a correttezza semantica, qualità dello schema e compromessi progettuali.
[AGGIUNTA] Fonte integrazione: riformulazione didattica del testo presente nella slide 4 del file Lezione 01 - Dipendenze funzionali e Normalizzazione di DB Relazionali.pdf.
\subsection*{Derivazioni}
[AGGIUNTA] In questa slide non sono state rilevate equazioni esplicite nel testo estratto; non è quindi richiesta una derivazione algebrica puntuale.
\subsection*{Esempio numerico / esercizio}
[AGGIUNTA] Esempio: dato uno schema R(A,B,C) con dipendenze funzionali A->B e B->C, mostrare la chiusura A^+ e discutere se A è superchiave. Soluzione sintetica: A^+={A,B,C}, quindi A determina tutti gli attributi di R ed è superchiave.
\subsection*{Note su assunzioni o integrazioni}
[ASSUNZIONE] L'estrazione OCR può perdere formattazione o simboli della slide originale.
[AGGIUNTA] Le spiegazioni sono estensioni didattiche per rendere espliciti passaggi impliciti nel materiale proiettato.
\paragraph{Riferimento fonte} File: Lezione 01 - Dipendenze funzionali e Normalizzazione di DB Relazionali.pdf. Numero slide: 4.
\section{Misure informali di qualità --- Slide 5}
\subsection*{Estratto originale dalla slide 5}
\begin{lstlisting}
Misure informali di qualità
-    Esistono delle misure informali di qualità per il
    disegno di schemi di relazione:
     1. Semantica degli attributi.
     2. Riduzione dei valori ridondanti nelle tuple.

     3. Riduzione dei valori null nelle tuple.

     4. Non consentire tuple spurie.



-    Tali misure non sempre sono indipendenti.
\end{lstlisting}
\subsection*{Spiegazione estesa}
[AGGIUNTA] Questa spiegazione collega il contenuto testuale della slide al contesto del corso. La slide enfatizza i concetti mostrati nell'estratto e il loro ruolo nella progettazione di basi di dati relazionali, con particolare attenzione a correttezza semantica, qualità dello schema e compromessi progettuali.
[AGGIUNTA] Fonte integrazione: riformulazione didattica del testo presente nella slide 5 del file Lezione 01 - Dipendenze funzionali e Normalizzazione di DB Relazionali.pdf.
\subsection*{Derivazioni}
[AGGIUNTA] In questa slide non sono state rilevate equazioni esplicite nel testo estratto; non è quindi richiesta una derivazione algebrica puntuale.
\subsection*{Esempio numerico / esercizio}
[AGGIUNTA] Esempio: dato uno schema R(A,B,C) con dipendenze funzionali A->B e B->C, mostrare la chiusura A^+ e discutere se A è superchiave. Soluzione sintetica: A^+={A,B,C}, quindi A determina tutti gli attributi di R ed è superchiave.
\subsection*{Note su assunzioni o integrazioni}
[ASSUNZIONE] L'estrazione OCR può perdere formattazione o simboli della slide originale.
[AGGIUNTA] Le spiegazioni sono estensioni didattiche per rendere espliciti passaggi impliciti nel materiale proiettato.
\paragraph{Riferimento fonte} File: Lezione 01 - Dipendenze funzionali e Normalizzazione di DB Relazionali.pdf. Numero slide: 5.
\section{Modello ER Company --- Slide 6}
\subsection*{Estratto originale dalla slide 6}
\begin{lstlisting}
Modello ER Company
        Minit                                             N                           1                     Number
                                                                  WORKS_FOR
Fname           Lname                                                                                                  Locations
                                                                                                     Name
                            Address
    Name                                                    Start_date
                                                                                                             DEPARTMENT
                      Sex          Salary
                                                              1                   1
                                                                      MANAGES
    Ssn
                        EMPLOYEE                                                                                         1
                                                                                      Number_of_employees
                                                                          Hours                                CONTROLS
   Bdate
                                                                      M                   N                          N
         supervisor                    supervisee                           WORKS_ON
                 1      SUPERVISION     N                         1                                          PROJECT


                                                                                                   Name
                                                                                                                         Location
                                                                      N                                     Number
                                                          DEPENDENT



                                       Name         Sex           Birth_date              Relationship
\end{lstlisting}
\subsection*{Spiegazione estesa}
[AGGIUNTA] Questa spiegazione collega il contenuto testuale della slide al contesto del corso. La slide enfatizza i concetti mostrati nell'estratto e il loro ruolo nella progettazione di basi di dati relazionali, con particolare attenzione a correttezza semantica, qualità dello schema e compromessi progettuali.
[AGGIUNTA] Fonte integrazione: riformulazione didattica del testo presente nella slide 6 del file Lezione 01 - Dipendenze funzionali e Normalizzazione di DB Relazionali.pdf.
\subsection*{Derivazioni}
[AGGIUNTA] In questa slide non sono state rilevate equazioni esplicite nel testo estratto; non è quindi richiesta una derivazione algebrica puntuale.
\subsection*{Esempio numerico / esercizio}
[AGGIUNTA] Esempio: dato uno schema R(A,B,C) con dipendenze funzionali A->B e B->C, mostrare la chiusura A^+ e discutere se A è superchiave. Soluzione sintetica: A^+={A,B,C}, quindi A determina tutti gli attributi di R ed è superchiave.
\subsection*{Note su assunzioni o integrazioni}
[ASSUNZIONE] L'estrazione OCR può perdere formattazione o simboli della slide originale.
[AGGIUNTA] Le spiegazioni sono estensioni didattiche per rendere espliciti passaggi impliciti nel materiale proiettato.
\paragraph{Riferimento fonte} File: Lezione 01 - Dipendenze funzionali e Normalizzazione di DB Relazionali.pdf. Numero slide: 6.
\section{Outline --- Slide 7}
\subsection*{Estratto originale dalla slide 7}
\begin{lstlisting}
Outline
  1. Semantica degli attributi.
  2. Riduzione dei valori ridondanti nelle tuple.

  3. Riduzione dei valori null nelle tuple.

  4. Non consentire tuple spurie.
\end{lstlisting}
\subsection*{Spiegazione estesa}
[AGGIUNTA] Questa spiegazione collega il contenuto testuale della slide al contesto del corso. La slide enfatizza i concetti mostrati nell'estratto e il loro ruolo nella progettazione di basi di dati relazionali, con particolare attenzione a correttezza semantica, qualità dello schema e compromessi progettuali.
[AGGIUNTA] Fonte integrazione: riformulazione didattica del testo presente nella slide 7 del file Lezione 01 - Dipendenze funzionali e Normalizzazione di DB Relazionali.pdf.
\subsection*{Derivazioni}
[AGGIUNTA] In questa slide non sono state rilevate equazioni esplicite nel testo estratto; non è quindi richiesta una derivazione algebrica puntuale.
\subsection*{Esempio numerico / esercizio}
[AGGIUNTA] Esempio: dato uno schema R(A,B,C) con dipendenze funzionali A->B e B->C, mostrare la chiusura A^+ e discutere se A è superchiave. Soluzione sintetica: A^+={A,B,C}, quindi A determina tutti gli attributi di R ed è superchiave.
\subsection*{Note su assunzioni o integrazioni}
[ASSUNZIONE] L'estrazione OCR può perdere formattazione o simboli della slide originale.
[AGGIUNTA] Le spiegazioni sono estensioni didattiche per rendere espliciti passaggi impliciti nel materiale proiettato.
\paragraph{Riferimento fonte} File: Lezione 01 - Dipendenze funzionali e Normalizzazione di DB Relazionali.pdf. Numero slide: 7.
\section{1) Semantica degli attributi di una --- Slide 8}
\subsection*{Estratto originale dalla slide 8}
\begin{lstlisting}
1)      Semantica degli attributi di una
        relazione
-    Quando si raggruppano gli attributi in uno
    schema di relazione, assumiamo che essi
    abbiano associato un significato.

-    Il significato, o semantica, specifica come
    interpretare i valori degli attributi di una relazione:
      Più è semplice spiegare la semantica della relazione,
      migliore è il disegno dello schema di relazione.
\end{lstlisting}
\subsection*{Spiegazione estesa}
[AGGIUNTA] Questa spiegazione collega il contenuto testuale della slide al contesto del corso. La slide enfatizza i concetti mostrati nell'estratto e il loro ruolo nella progettazione di basi di dati relazionali, con particolare attenzione a correttezza semantica, qualità dello schema e compromessi progettuali.
[AGGIUNTA] Fonte integrazione: riformulazione didattica del testo presente nella slide 8 del file Lezione 01 - Dipendenze funzionali e Normalizzazione di DB Relazionali.pdf.
\subsection*{Derivazioni}
[AGGIUNTA] In questa slide non sono state rilevate equazioni esplicite nel testo estratto; non è quindi richiesta una derivazione algebrica puntuale.
\subsection*{Esempio numerico / esercizio}
[AGGIUNTA] Esempio: dato uno schema R(A,B,C) con dipendenze funzionali A->B e B->C, mostrare la chiusura A^+ e discutere se A è superchiave. Soluzione sintetica: A^+={A,B,C}, quindi A determina tutti gli attributi di R ed è superchiave.
\subsection*{Note su assunzioni o integrazioni}
[ASSUNZIONE] L'estrazione OCR può perdere formattazione o simboli della slide originale.
[AGGIUNTA] Le spiegazioni sono estensioni didattiche per rendere espliciti passaggi impliciti nel materiale proiettato.
\paragraph{Riferimento fonte} File: Lezione 01 - Dipendenze funzionali e Normalizzazione di DB Relazionali.pdf. Numero slide: 8.
\section{1) Semantica degli attributi di una relazione: --- Slide 9}
\subsection*{Estratto originale dalla slide 9}
\begin{lstlisting}
1) Semantica degli attributi di una relazione:
Esempio




                                Versione semplificata
                                dello schema del
                                database Company
\end{lstlisting}
\subsection*{Spiegazione estesa}
[AGGIUNTA] Questa spiegazione collega il contenuto testuale della slide al contesto del corso. La slide enfatizza i concetti mostrati nell'estratto e il loro ruolo nella progettazione di basi di dati relazionali, con particolare attenzione a correttezza semantica, qualità dello schema e compromessi progettuali.
[AGGIUNTA] Fonte integrazione: riformulazione didattica del testo presente nella slide 9 del file Lezione 01 - Dipendenze funzionali e Normalizzazione di DB Relazionali.pdf.
\subsection*{Derivazioni}
[AGGIUNTA] In questa slide non sono state rilevate equazioni esplicite nel testo estratto; non è quindi richiesta una derivazione algebrica puntuale.
\subsection*{Esempio numerico / esercizio}
[AGGIUNTA] Esempio: dato uno schema R(A,B,C) con dipendenze funzionali A->B e B->C, mostrare la chiusura A^+ e discutere se A è superchiave. Soluzione sintetica: A^+={A,B,C}, quindi A determina tutti gli attributi di R ed è superchiave.
\subsection*{Note su assunzioni o integrazioni}
[ASSUNZIONE] L'estrazione OCR può perdere formattazione o simboli della slide originale.
[AGGIUNTA] Le spiegazioni sono estensioni didattiche per rendere espliciti passaggi impliciti nel materiale proiettato.
\paragraph{Riferimento fonte} File: Lezione 01 - Dipendenze funzionali e Normalizzazione di DB Relazionali.pdf. Numero slide: 9.
\section{1) Semantica degli attributi di una relazione: --- Slide 10}
\subsection*{Estratto originale dalla slide 10}
\begin{lstlisting}
1) Semantica degli attributi di una relazione:
Esempio (2)




                                 Esempio di relazioni
                                 dello schema del db
                                 Company
\end{lstlisting}
\subsection*{Spiegazione estesa}
[AGGIUNTA] Questa spiegazione collega il contenuto testuale della slide al contesto del corso. La slide enfatizza i concetti mostrati nell'estratto e il loro ruolo nella progettazione di basi di dati relazionali, con particolare attenzione a correttezza semantica, qualità dello schema e compromessi progettuali.
[AGGIUNTA] Fonte integrazione: riformulazione didattica del testo presente nella slide 10 del file Lezione 01 - Dipendenze funzionali e Normalizzazione di DB Relazionali.pdf.
\subsection*{Derivazioni}
[AGGIUNTA] In questa slide non sono state rilevate equazioni esplicite nel testo estratto; non è quindi richiesta una derivazione algebrica puntuale.
\subsection*{Esempio numerico / esercizio}
[AGGIUNTA] Esempio: dato uno schema R(A,B,C) con dipendenze funzionali A->B e B->C, mostrare la chiusura A^+ e discutere se A è superchiave. Soluzione sintetica: A^+={A,B,C}, quindi A determina tutti gli attributi di R ed è superchiave.
\subsection*{Note su assunzioni o integrazioni}
[ASSUNZIONE] L'estrazione OCR può perdere formattazione o simboli della slide originale.
[AGGIUNTA] Le spiegazioni sono estensioni didattiche per rendere espliciti passaggi impliciti nel materiale proiettato.
\paragraph{Riferimento fonte} File: Lezione 01 - Dipendenze funzionali e Normalizzazione di DB Relazionali.pdf. Numero slide: 10.
\section{1) Semantica degli attributi di una relazione: --- Slide 11}
\subsection*{Estratto originale dalla slide 11}
\begin{lstlisting}
1) Semantica degli attributi di una relazione:
Esempio (3)
-    Significato dello schema del database:

      ENAME, SSN, BDATE, ADDRESS -> dati
      dell'impiegato.
      DNUMBER -> dipartimento per cui lavora.
      DNUMBER (foreign key - f.k.) -> relazione implicita tra
      EMPLOYEE e DEPARTMENT.
      DEPARTMENT -> un'entità dipartimento.
      PROJECT -> una entità progetto.
      f.k. DMGRSSN (di DEPARTMENT) -> correla il
      dipartimento all'impiegato che è suo manager.
\end{lstlisting}
\subsection*{Spiegazione estesa}
[AGGIUNTA] Questa spiegazione collega il contenuto testuale della slide al contesto del corso. La slide enfatizza i concetti mostrati nell'estratto e il loro ruolo nella progettazione di basi di dati relazionali, con particolare attenzione a correttezza semantica, qualità dello schema e compromessi progettuali.
[AGGIUNTA] Fonte integrazione: riformulazione didattica del testo presente nella slide 11 del file Lezione 01 - Dipendenze funzionali e Normalizzazione di DB Relazionali.pdf.
\subsection*{Derivazioni}
[AGGIUNTA] Nella slide compaiono espressioni simboliche/relazioni. Derivazione didattica generale:
\begin{align*}
X^+ &= X \\
X^+ &\leftarrow X^+ \cup Y \quad \text{se } (U \to V) \in F \text{ e } U \subseteq X^+ \text{ con } Y=V \\
\text{Stop} &\text{ quando non si aggiungono nuovi attributi in } X^+
\end{align*}
[AGGIUNTA] Questa procedura esplicita il metodo di chiusura usato nelle slide su dipendenze e chiavi.
\begin{itemize}[leftmargin=*]
\item Forma osservata nella slide: \texttt{ENAME, SSN, BDATE, ADDRESS -> dati}
\item Forma osservata nella slide: \texttt{DNUMBER -> dipartimento per cui lavora.}
\item Forma osservata nella slide: \texttt{DNUMBER (foreign key - f.k.) -> relazione implicita tra}
\end{itemize}
\subsection*{Esempio numerico / esercizio}
[AGGIUNTA] Esempio: dato uno schema R(A,B,C) con dipendenze funzionali A->B e B->C, mostrare la chiusura A^+ e discutere se A è superchiave. Soluzione sintetica: A^+={A,B,C}, quindi A determina tutti gli attributi di R ed è superchiave.
\subsection*{Note su assunzioni o integrazioni}
[ASSUNZIONE] L'estrazione OCR può perdere formattazione o simboli della slide originale.
[AGGIUNTA] Le spiegazioni sono estensioni didattiche per rendere espliciti passaggi impliciti nel materiale proiettato.
\paragraph{Riferimento fonte} File: Lezione 01 - Dipendenze funzionali e Normalizzazione di DB Relazionali.pdf. Numero slide: 11.
\section{1) Semantica degli attributi di una relazione: --- Slide 12}
\subsection*{Estratto originale dalla slide 12}
\begin{lstlisting}
1) Semantica degli attributi di una relazione:
Esempio (4)
    f.k. DNUM (di PROJECT) -> correla un progetto al suo
   dipartimento di controllo.

    Ogni tupla in DEPT_LOCATIONS contiene un numero
   di dipartimento (DNUMBER) e una delle locazioni del
   dipartimento (DLOCATIONS).

    Ogni tupla in WORKS_ON contiene l'SSN di un
   impiegato, il numero di progetto PNUMBER di uno dei
   progetti su cui lavora l'impiegato, e il numero di ore
   settimanali HOURS.
\end{lstlisting}
\subsection*{Spiegazione estesa}
[AGGIUNTA] Questa spiegazione collega il contenuto testuale della slide al contesto del corso. La slide enfatizza i concetti mostrati nell'estratto e il loro ruolo nella progettazione di basi di dati relazionali, con particolare attenzione a correttezza semantica, qualità dello schema e compromessi progettuali.
[AGGIUNTA] Fonte integrazione: riformulazione didattica del testo presente nella slide 12 del file Lezione 01 - Dipendenze funzionali e Normalizzazione di DB Relazionali.pdf.
\subsection*{Derivazioni}
[AGGIUNTA] Nella slide compaiono espressioni simboliche/relazioni. Derivazione didattica generale:
\begin{align*}
X^+ &= X \\
X^+ &\leftarrow X^+ \cup Y \quad \text{se } (U \to V) \in F \text{ e } U \subseteq X^+ \text{ con } Y=V \\
\text{Stop} &\text{ quando non si aggiungono nuovi attributi in } X^+
\end{align*}
[AGGIUNTA] Questa procedura esplicita il metodo di chiusura usato nelle slide su dipendenze e chiavi.
\begin{itemize}[leftmargin=*]
\item Forma osservata nella slide: \texttt{f.k. DNUM (di PROJECT) -> correla un progetto al suo}
\end{itemize}
\subsection*{Esempio numerico / esercizio}
[AGGIUNTA] Esempio: dato uno schema R(A,B,C) con dipendenze funzionali A->B e B->C, mostrare la chiusura A^+ e discutere se A è superchiave. Soluzione sintetica: A^+={A,B,C}, quindi A determina tutti gli attributi di R ed è superchiave.
\subsection*{Note su assunzioni o integrazioni}
[ASSUNZIONE] L'estrazione OCR può perdere formattazione o simboli della slide originale.
[AGGIUNTA] Le spiegazioni sono estensioni didattiche per rendere espliciti passaggi impliciti nel materiale proiettato.
\paragraph{Riferimento fonte} File: Lezione 01 - Dipendenze funzionali e Normalizzazione di DB Relazionali.pdf. Numero slide: 12.
\section{1) Semantica degli attributi di una relazione: --- Slide 13}
\subsection*{Estratto originale dalla slide 13}
\begin{lstlisting}
1) Semantica degli attributi di una relazione:
Esempio (5)
      DEPT_LOCATIONS -> rappresenta un attributo
      multivalued di DEPARTMENT.
      WORKS_ON -> rappresenta una relazione di
      cardinalità M:N tra EMPLOYEE e PROJECT.


-    Tutte le relazioni quindi possono essere
    considerate "buone" dal punto di vista di avere
    una semantica chiara.
\end{lstlisting}
\subsection*{Spiegazione estesa}
[AGGIUNTA] Questa spiegazione collega il contenuto testuale della slide al contesto del corso. La slide enfatizza i concetti mostrati nell'estratto e il loro ruolo nella progettazione di basi di dati relazionali, con particolare attenzione a correttezza semantica, qualità dello schema e compromessi progettuali.
[AGGIUNTA] Fonte integrazione: riformulazione didattica del testo presente nella slide 13 del file Lezione 01 - Dipendenze funzionali e Normalizzazione di DB Relazionali.pdf.
\subsection*{Derivazioni}
[AGGIUNTA] Nella slide compaiono espressioni simboliche/relazioni. Derivazione didattica generale:
\begin{align*}
X^+ &= X \\
X^+ &\leftarrow X^+ \cup Y \quad \text{se } (U \to V) \in F \text{ e } U \subseteq X^+ \text{ con } Y=V \\
\text{Stop} &\text{ quando non si aggiungono nuovi attributi in } X^+
\end{align*}
[AGGIUNTA] Questa procedura esplicita il metodo di chiusura usato nelle slide su dipendenze e chiavi.
\begin{itemize}[leftmargin=*]
\item Forma osservata nella slide: \texttt{DEPT\_LOCATIONS -> rappresenta un attributo}
\item Forma osservata nella slide: \texttt{WORKS\_ON -> rappresenta una relazione di}
\end{itemize}
\subsection*{Esempio numerico / esercizio}
[AGGIUNTA] Esempio: dato uno schema R(A,B,C) con dipendenze funzionali A->B e B->C, mostrare la chiusura A^+ e discutere se A è superchiave. Soluzione sintetica: A^+={A,B,C}, quindi A determina tutti gli attributi di R ed è superchiave.
\subsection*{Note su assunzioni o integrazioni}
[ASSUNZIONE] L'estrazione OCR può perdere formattazione o simboli della slide originale.
[AGGIUNTA] Le spiegazioni sono estensioni didattiche per rendere espliciti passaggi impliciti nel materiale proiettato.
\paragraph{Riferimento fonte} File: Lezione 01 - Dipendenze funzionali e Normalizzazione di DB Relazionali.pdf. Numero slide: 13.
\section{Linea guida 1 --- Slide 14}
\subsection*{Estratto originale dalla slide 14}
\begin{lstlisting}
Linea guida 1
-    Disegnare uno schema di relazione del quale sia
    facile spiegarne il significato.

-    Non combinare attributi da tipi di entità e tipi di
    relazioni differenti in una singola relazione.
      Intuitivamente: se uno schema di relazione corrisponde
      ad un tipo di entità o ad un tipo di relazione, il
      significato tende ad essere chiaro.
\end{lstlisting}
\subsection*{Spiegazione estesa}
[AGGIUNTA] Questa spiegazione collega il contenuto testuale della slide al contesto del corso. La slide enfatizza i concetti mostrati nell'estratto e il loro ruolo nella progettazione di basi di dati relazionali, con particolare attenzione a correttezza semantica, qualità dello schema e compromessi progettuali.
[AGGIUNTA] Fonte integrazione: riformulazione didattica del testo presente nella slide 14 del file Lezione 01 - Dipendenze funzionali e Normalizzazione di DB Relazionali.pdf.
\subsection*{Derivazioni}
[AGGIUNTA] In questa slide non sono state rilevate equazioni esplicite nel testo estratto; non è quindi richiesta una derivazione algebrica puntuale.
\subsection*{Esempio numerico / esercizio}
[AGGIUNTA] Esempio: dato uno schema R(A,B,C) con dipendenze funzionali A->B e B->C, mostrare la chiusura A^+ e discutere se A è superchiave. Soluzione sintetica: A^+={A,B,C}, quindi A determina tutti gli attributi di R ed è superchiave.
\subsection*{Note su assunzioni o integrazioni}
[ASSUNZIONE] L'estrazione OCR può perdere formattazione o simboli della slide originale.
[AGGIUNTA] Le spiegazioni sono estensioni didattiche per rendere espliciti passaggi impliciti nel materiale proiettato.
\paragraph{Riferimento fonte} File: Lezione 01 - Dipendenze funzionali e Normalizzazione di DB Relazionali.pdf. Numero slide: 14.
\section{Linea guida 1: Esempio --- Slide 15}
\subsection*{Estratto originale dalla slide 15}
\begin{lstlisting}
Linea guida 1: Esempio
-    Considerando le due relazioni:




       EMPT_DEPT -> ogni entità rappresenta un impiegato ma include
       informazioni sul dipartimento in cui lavora.
       EMP_PROJ -> ogni tupla correla un impiegato a un progetto ma
       richiede anche il nome dell'impiegato ENAME e il nome e la
       locazione del progetto, PNAME e PLOCATION.

-    Tali schemi di relazione hanno anch'essi una semantica
    chiara, però entrambe contravvengono la linea guida 1,
    contenendo attributi di entità distinte.
\end{lstlisting}
\subsection*{Spiegazione estesa}
[AGGIUNTA] Questa spiegazione collega il contenuto testuale della slide al contesto del corso. La slide enfatizza i concetti mostrati nell'estratto e il loro ruolo nella progettazione di basi di dati relazionali, con particolare attenzione a correttezza semantica, qualità dello schema e compromessi progettuali.
[AGGIUNTA] Fonte integrazione: riformulazione didattica del testo presente nella slide 15 del file Lezione 01 - Dipendenze funzionali e Normalizzazione di DB Relazionali.pdf.
\subsection*{Derivazioni}
[AGGIUNTA] Nella slide compaiono espressioni simboliche/relazioni. Derivazione didattica generale:
\begin{align*}
X^+ &= X \\
X^+ &\leftarrow X^+ \cup Y \quad \text{se } (U \to V) \in F \text{ e } U \subseteq X^+ \text{ con } Y=V \\
\text{Stop} &\text{ quando non si aggiungono nuovi attributi in } X^+
\end{align*}
[AGGIUNTA] Questa procedura esplicita il metodo di chiusura usato nelle slide su dipendenze e chiavi.
\begin{itemize}[leftmargin=*]
\item Forma osservata nella slide: \texttt{EMPT\_DEPT -> ogni entità rappresenta un impiegato ma include}
\item Forma osservata nella slide: \texttt{EMP\_PROJ -> ogni tupla correla un impiegato a un progetto ma}
\end{itemize}
\subsection*{Esempio numerico / esercizio}
[AGGIUNTA] Esempio: dato uno schema R(A,B,C) con dipendenze funzionali A->B e B->C, mostrare la chiusura A^+ e discutere se A è superchiave. Soluzione sintetica: A^+={A,B,C}, quindi A determina tutti gli attributi di R ed è superchiave.
\subsection*{Note su assunzioni o integrazioni}
[ASSUNZIONE] L'estrazione OCR può perdere formattazione o simboli della slide originale.
[AGGIUNTA] Le spiegazioni sono estensioni didattiche per rendere espliciti passaggi impliciti nel materiale proiettato.
\paragraph{Riferimento fonte} File: Lezione 01 - Dipendenze funzionali e Normalizzazione di DB Relazionali.pdf. Numero slide: 15.
\section{Outline --- Slide 16}
\subsection*{Estratto originale dalla slide 16}
\begin{lstlisting}
Outline
  1. Semantica degli attributi.
  2. Riduzione dei valori ridondanti nelle tuple.

  3. Riduzione dei valori null nelle tuple.

  4. Non consentire tuple spurie.
\end{lstlisting}
\subsection*{Spiegazione estesa}
[AGGIUNTA] Questa spiegazione collega il contenuto testuale della slide al contesto del corso. La slide enfatizza i concetti mostrati nell'estratto e il loro ruolo nella progettazione di basi di dati relazionali, con particolare attenzione a correttezza semantica, qualità dello schema e compromessi progettuali.
[AGGIUNTA] Fonte integrazione: riformulazione didattica del testo presente nella slide 16 del file Lezione 01 - Dipendenze funzionali e Normalizzazione di DB Relazionali.pdf.
\subsection*{Derivazioni}
[AGGIUNTA] In questa slide non sono state rilevate equazioni esplicite nel testo estratto; non è quindi richiesta una derivazione algebrica puntuale.
\subsection*{Esempio numerico / esercizio}
[AGGIUNTA] Esempio: dato uno schema R(A,B,C) con dipendenze funzionali A->B e B->C, mostrare la chiusura A^+ e discutere se A è superchiave. Soluzione sintetica: A^+={A,B,C}, quindi A determina tutti gli attributi di R ed è superchiave.
\subsection*{Note su assunzioni o integrazioni}
[ASSUNZIONE] L'estrazione OCR può perdere formattazione o simboli della slide originale.
[AGGIUNTA] Le spiegazioni sono estensioni didattiche per rendere espliciti passaggi impliciti nel materiale proiettato.
\paragraph{Riferimento fonte} File: Lezione 01 - Dipendenze funzionali e Normalizzazione di DB Relazionali.pdf. Numero slide: 16.
\section{2) Riduzione dei valori ridondanti nelle tuple --- Slide 17}
\subsection*{Estratto originale dalla slide 17}
\begin{lstlisting}
2) Riduzione dei valori ridondanti nelle tuple

-    Un obiettivo del disegno dello schema è di
    minimizzare lo spazio di memoria occupato dalle
    relazioni base.
\end{lstlisting}
\subsection*{Spiegazione estesa}
[AGGIUNTA] Questa spiegazione collega il contenuto testuale della slide al contesto del corso. La slide enfatizza i concetti mostrati nell'estratto e il loro ruolo nella progettazione di basi di dati relazionali, con particolare attenzione a correttezza semantica, qualità dello schema e compromessi progettuali.
[AGGIUNTA] Fonte integrazione: riformulazione didattica del testo presente nella slide 17 del file Lezione 01 - Dipendenze funzionali e Normalizzazione di DB Relazionali.pdf.
\subsection*{Derivazioni}
[AGGIUNTA] In questa slide non sono state rilevate equazioni esplicite nel testo estratto; non è quindi richiesta una derivazione algebrica puntuale.
\subsection*{Esempio numerico / esercizio}
[AGGIUNTA] Esempio: dato uno schema R(A,B,C) con dipendenze funzionali A->B e B->C, mostrare la chiusura A^+ e discutere se A è superchiave. Soluzione sintetica: A^+={A,B,C}, quindi A determina tutti gli attributi di R ed è superchiave.
\subsection*{Note su assunzioni o integrazioni}
[ASSUNZIONE] L'estrazione OCR può perdere formattazione o simboli della slide originale.
[AGGIUNTA] Le spiegazioni sono estensioni didattiche per rendere espliciti passaggi impliciti nel materiale proiettato.
\paragraph{Riferimento fonte} File: Lezione 01 - Dipendenze funzionali e Normalizzazione di DB Relazionali.pdf. Numero slide: 17.
\section{2) Riduzione dei valori ridondanti nelle --- Slide 18}
\subsection*{Estratto originale dalla slide 18}
\begin{lstlisting}
2) Riduzione dei valori ridondanti nelle
tuple: Esempio




                                  Schema corretto
\end{lstlisting}
\subsection*{Spiegazione estesa}
[AGGIUNTA] Questa spiegazione collega il contenuto testuale della slide al contesto del corso. La slide enfatizza i concetti mostrati nell'estratto e il loro ruolo nella progettazione di basi di dati relazionali, con particolare attenzione a correttezza semantica, qualità dello schema e compromessi progettuali.
[AGGIUNTA] Fonte integrazione: riformulazione didattica del testo presente nella slide 18 del file Lezione 01 - Dipendenze funzionali e Normalizzazione di DB Relazionali.pdf.
\subsection*{Derivazioni}
[AGGIUNTA] In questa slide non sono state rilevate equazioni esplicite nel testo estratto; non è quindi richiesta una derivazione algebrica puntuale.
\subsection*{Esempio numerico / esercizio}
[AGGIUNTA] Esempio: dato uno schema R(A,B,C) con dipendenze funzionali A->B e B->C, mostrare la chiusura A^+ e discutere se A è superchiave. Soluzione sintetica: A^+={A,B,C}, quindi A determina tutti gli attributi di R ed è superchiave.
\subsection*{Note su assunzioni o integrazioni}
[ASSUNZIONE] L'estrazione OCR può perdere formattazione o simboli della slide originale.
[AGGIUNTA] Le spiegazioni sono estensioni didattiche per rendere espliciti passaggi impliciti nel materiale proiettato.
\paragraph{Riferimento fonte} File: Lezione 01 - Dipendenze funzionali e Normalizzazione di DB Relazionali.pdf. Numero slide: 18.
\section{2) Riduzione dei valori ridondanti nelle --- Slide 19}
\subsection*{Estratto originale dalla slide 19}
\begin{lstlisting}
2) Riduzione dei valori ridondanti nelle
tuple: Esempio (2)




                                   Schema errato
\end{lstlisting}
\subsection*{Spiegazione estesa}
[AGGIUNTA] Questa spiegazione collega il contenuto testuale della slide al contesto del corso. La slide enfatizza i concetti mostrati nell'estratto e il loro ruolo nella progettazione di basi di dati relazionali, con particolare attenzione a correttezza semantica, qualità dello schema e compromessi progettuali.
[AGGIUNTA] Fonte integrazione: riformulazione didattica del testo presente nella slide 19 del file Lezione 01 - Dipendenze funzionali e Normalizzazione di DB Relazionali.pdf.
\subsection*{Derivazioni}
[AGGIUNTA] In questa slide non sono state rilevate equazioni esplicite nel testo estratto; non è quindi richiesta una derivazione algebrica puntuale.
\subsection*{Esempio numerico / esercizio}
[AGGIUNTA] Esempio: dato uno schema R(A,B,C) con dipendenze funzionali A->B e B->C, mostrare la chiusura A^+ e discutere se A è superchiave. Soluzione sintetica: A^+={A,B,C}, quindi A determina tutti gli attributi di R ed è superchiave.
\subsection*{Note su assunzioni o integrazioni}
[ASSUNZIONE] L'estrazione OCR può perdere formattazione o simboli della slide originale.
[AGGIUNTA] Le spiegazioni sono estensioni didattiche per rendere espliciti passaggi impliciti nel materiale proiettato.
\paragraph{Riferimento fonte} File: Lezione 01 - Dipendenze funzionali e Normalizzazione di DB Relazionali.pdf. Numero slide: 19.
\section{2) Riduzione dei valori ridondanti nelle --- Slide 20}
\subsection*{Estratto originale dalla slide 20}
\begin{lstlisting}
2) Riduzione dei valori ridondanti nelle
tuple: Esempio (3)
-    Oltre all'occupazione di più spazio, le relazioni
    EMP_DEPT e EMP_PROJ usate come relazioni
    base, presentano anche il problema delle
    "anomalie di aggiornamento" (update anomalies)...
\end{lstlisting}
\subsection*{Spiegazione estesa}
[AGGIUNTA] Questa spiegazione collega il contenuto testuale della slide al contesto del corso. La slide enfatizza i concetti mostrati nell'estratto e il loro ruolo nella progettazione di basi di dati relazionali, con particolare attenzione a correttezza semantica, qualità dello schema e compromessi progettuali.
[AGGIUNTA] Fonte integrazione: riformulazione didattica del testo presente nella slide 20 del file Lezione 01 - Dipendenze funzionali e Normalizzazione di DB Relazionali.pdf.
\subsection*{Derivazioni}
[AGGIUNTA] In questa slide non sono state rilevate equazioni esplicite nel testo estratto; non è quindi richiesta una derivazione algebrica puntuale.
\subsection*{Esempio numerico / esercizio}
[AGGIUNTA] Esempio: dato uno schema R(A,B,C) con dipendenze funzionali A->B e B->C, mostrare la chiusura A^+ e discutere se A è superchiave. Soluzione sintetica: A^+={A,B,C}, quindi A determina tutti gli attributi di R ed è superchiave.
\subsection*{Note su assunzioni o integrazioni}
[ASSUNZIONE] L'estrazione OCR può perdere formattazione o simboli della slide originale.
[AGGIUNTA] Le spiegazioni sono estensioni didattiche per rendere espliciti passaggi impliciti nel materiale proiettato.
\paragraph{Riferimento fonte} File: Lezione 01 - Dipendenze funzionali e Normalizzazione di DB Relazionali.pdf. Numero slide: 20.
\section{Update Anomalies --- Slide 21}
\subsection*{Estratto originale dalla slide 21}
\begin{lstlisting}
Update Anomalies
-    Le "update anomalies" sono delle anomalie che
    possono sorgere nel gestire i dati di un database
    relazionale non correttamente progettato.

-    Si dividono in:
      Insertion Anomalies
      Deletion Anomalies
      Modification Anomalies
\end{lstlisting}
\subsection*{Spiegazione estesa}
[AGGIUNTA] Questa spiegazione collega il contenuto testuale della slide al contesto del corso. La slide enfatizza i concetti mostrati nell'estratto e il loro ruolo nella progettazione di basi di dati relazionali, con particolare attenzione a correttezza semantica, qualità dello schema e compromessi progettuali.
[AGGIUNTA] Fonte integrazione: riformulazione didattica del testo presente nella slide 21 del file Lezione 01 - Dipendenze funzionali e Normalizzazione di DB Relazionali.pdf.
\subsection*{Derivazioni}
[AGGIUNTA] In questa slide non sono state rilevate equazioni esplicite nel testo estratto; non è quindi richiesta una derivazione algebrica puntuale.
\subsection*{Esempio numerico / esercizio}
[AGGIUNTA] Esempio: dato uno schema R(A,B,C) con dipendenze funzionali A->B e B->C, mostrare la chiusura A^+ e discutere se A è superchiave. Soluzione sintetica: A^+={A,B,C}, quindi A determina tutti gli attributi di R ed è superchiave.
\subsection*{Note su assunzioni o integrazioni}
[ASSUNZIONE] L'estrazione OCR può perdere formattazione o simboli della slide originale.
[AGGIUNTA] Le spiegazioni sono estensioni didattiche per rendere espliciti passaggi impliciti nel materiale proiettato.
\paragraph{Riferimento fonte} File: Lezione 01 - Dipendenze funzionali e Normalizzazione di DB Relazionali.pdf. Numero slide: 21.
\section{Insertion Anomalies --- Slide 22}
\subsection*{Estratto originale dalla slide 22}
\begin{lstlisting}
Insertion Anomalies
-    Possono essere di due tipi, mostrati riferendosi
    allo schema EMP_DEPT:



     1.   Per inserire una nuova tupla impiegato, dobbiamo
          anche richiedere i valori degli attributi del
          dipartimento per cui lavora.
          Inoltre tali valori devono essere consistenti.

          Esempio: inserendo un nuovo impiegato che lavora nel
          dipartimento 5, dobbiamo inserire i valori degli attributi DNAME
          e DMGRSSN consistenti con tali attributi di altre tuple in
          EMP_DEPT per il dipartimento 5.
\end{lstlisting}
\subsection*{Spiegazione estesa}
[AGGIUNTA] Questa spiegazione collega il contenuto testuale della slide al contesto del corso. La slide enfatizza i concetti mostrati nell'estratto e il loro ruolo nella progettazione di basi di dati relazionali, con particolare attenzione a correttezza semantica, qualità dello schema e compromessi progettuali.
[AGGIUNTA] Fonte integrazione: riformulazione didattica del testo presente nella slide 22 del file Lezione 01 - Dipendenze funzionali e Normalizzazione di DB Relazionali.pdf.
\subsection*{Derivazioni}
[AGGIUNTA] In questa slide non sono state rilevate equazioni esplicite nel testo estratto; non è quindi richiesta una derivazione algebrica puntuale.
\subsection*{Esempio numerico / esercizio}
[AGGIUNTA] Esempio: dato uno schema R(A,B,C) con dipendenze funzionali A->B e B->C, mostrare la chiusura A^+ e discutere se A è superchiave. Soluzione sintetica: A^+={A,B,C}, quindi A determina tutti gli attributi di R ed è superchiave.
\subsection*{Note su assunzioni o integrazioni}
[ASSUNZIONE] L'estrazione OCR può perdere formattazione o simboli della slide originale.
[AGGIUNTA] Le spiegazioni sono estensioni didattiche per rendere espliciti passaggi impliciti nel materiale proiettato.
\paragraph{Riferimento fonte} File: Lezione 01 - Dipendenze funzionali e Normalizzazione di DB Relazionali.pdf. Numero slide: 22.
\section{Insertion Anomalies (2) --- Slide 23}
\subsection*{Estratto originale dalla slide 23}
\begin{lstlisting}
Insertion Anomalies (2)



  2.   È difficile inserire un nuovo dipartimento che non ha
       ancora impiegati nella relazione EMP_DEPT.
       L'unico modo è di inserirlo ponendo a null gli attributi
       per l'impiegato, ma ciò crea un problema poiché SSN
       è la chiave primaria.
       Inoltre inserendo il primo impiegato per quel
       dipartimento non servirebbe più la tupla con valori
       null.
\end{lstlisting}
\subsection*{Spiegazione estesa}
[AGGIUNTA] Questa spiegazione collega il contenuto testuale della slide al contesto del corso. La slide enfatizza i concetti mostrati nell'estratto e il loro ruolo nella progettazione di basi di dati relazionali, con particolare attenzione a correttezza semantica, qualità dello schema e compromessi progettuali.
[AGGIUNTA] Fonte integrazione: riformulazione didattica del testo presente nella slide 23 del file Lezione 01 - Dipendenze funzionali e Normalizzazione di DB Relazionali.pdf.
\subsection*{Derivazioni}
[AGGIUNTA] In questa slide non sono state rilevate equazioni esplicite nel testo estratto; non è quindi richiesta una derivazione algebrica puntuale.
\subsection*{Esempio numerico / esercizio}
[AGGIUNTA] Esempio: dato uno schema R(A,B,C) con dipendenze funzionali A->B e B->C, mostrare la chiusura A^+ e discutere se A è superchiave. Soluzione sintetica: A^+={A,B,C}, quindi A determina tutti gli attributi di R ed è superchiave.
\subsection*{Note su assunzioni o integrazioni}
[ASSUNZIONE] L'estrazione OCR può perdere formattazione o simboli della slide originale.
[AGGIUNTA] Le spiegazioni sono estensioni didattiche per rendere espliciti passaggi impliciti nel materiale proiettato.
\paragraph{Riferimento fonte} File: Lezione 01 - Dipendenze funzionali e Normalizzazione di DB Relazionali.pdf. Numero slide: 23.
\section{Deletion Anomalies --- Slide 24}
\subsection*{Estratto originale dalla slide 24}
\begin{lstlisting}
Deletion Anomalies
-    Questo problema è correlato al secondo
    problema delle insertion anomalies.

-    Se cancelliamo da EMP_DEPT una tupla
    impiegato che rappresenta l'ultimo impiegato che
    lavora per un particolare dipartimento,
    l'informazione sul dipartimento è persa dal
    database.
\end{lstlisting}
\subsection*{Spiegazione estesa}
[AGGIUNTA] Questa spiegazione collega il contenuto testuale della slide al contesto del corso. La slide enfatizza i concetti mostrati nell'estratto e il loro ruolo nella progettazione di basi di dati relazionali, con particolare attenzione a correttezza semantica, qualità dello schema e compromessi progettuali.
[AGGIUNTA] Fonte integrazione: riformulazione didattica del testo presente nella slide 24 del file Lezione 01 - Dipendenze funzionali e Normalizzazione di DB Relazionali.pdf.
\subsection*{Derivazioni}
[AGGIUNTA] In questa slide non sono state rilevate equazioni esplicite nel testo estratto; non è quindi richiesta una derivazione algebrica puntuale.
\subsection*{Esempio numerico / esercizio}
[AGGIUNTA] Esempio: dato uno schema R(A,B,C) con dipendenze funzionali A->B e B->C, mostrare la chiusura A^+ e discutere se A è superchiave. Soluzione sintetica: A^+={A,B,C}, quindi A determina tutti gli attributi di R ed è superchiave.
\subsection*{Note su assunzioni o integrazioni}
[ASSUNZIONE] L'estrazione OCR può perdere formattazione o simboli della slide originale.
[AGGIUNTA] Le spiegazioni sono estensioni didattiche per rendere espliciti passaggi impliciti nel materiale proiettato.
\paragraph{Riferimento fonte} File: Lezione 01 - Dipendenze funzionali e Normalizzazione di DB Relazionali.pdf. Numero slide: 24.
\section{Modification Anomalies --- Slide 25}
\subsection*{Estratto originale dalla slide 25}
\begin{lstlisting}
Modification Anomalies
-    Se cambiamo il valore di uno degli attributi di un
    particolare dipartimento, dobbiamo aggiornare
    tutte le tuple di tutti gli impiegati che lavorano per
    quel dipartimento.
    Altrimenti il database diventa inconsistente.
\end{lstlisting}
\subsection*{Spiegazione estesa}
[AGGIUNTA] Questa spiegazione collega il contenuto testuale della slide al contesto del corso. La slide enfatizza i concetti mostrati nell'estratto e il loro ruolo nella progettazione di basi di dati relazionali, con particolare attenzione a correttezza semantica, qualità dello schema e compromessi progettuali.
[AGGIUNTA] Fonte integrazione: riformulazione didattica del testo presente nella slide 25 del file Lezione 01 - Dipendenze funzionali e Normalizzazione di DB Relazionali.pdf.
\subsection*{Derivazioni}
[AGGIUNTA] In questa slide non sono state rilevate equazioni esplicite nel testo estratto; non è quindi richiesta una derivazione algebrica puntuale.
\subsection*{Esempio numerico / esercizio}
[AGGIUNTA] Esempio: dato uno schema R(A,B,C) con dipendenze funzionali A->B e B->C, mostrare la chiusura A^+ e discutere se A è superchiave. Soluzione sintetica: A^+={A,B,C}, quindi A determina tutti gli attributi di R ed è superchiave.
\subsection*{Note su assunzioni o integrazioni}
[ASSUNZIONE] L'estrazione OCR può perdere formattazione o simboli della slide originale.
[AGGIUNTA] Le spiegazioni sono estensioni didattiche per rendere espliciti passaggi impliciti nel materiale proiettato.
\paragraph{Riferimento fonte} File: Lezione 01 - Dipendenze funzionali e Normalizzazione di DB Relazionali.pdf. Numero slide: 25.
\section{Linea guida 2 --- Slide 26}
\subsection*{Estratto originale dalla slide 26}
\begin{lstlisting}
Linea guida 2
-    Disegnare gli schemi di relazione di base in modo
    che non possano accadere insertion, deletion o
    modification anomalies.
      Se qualcuna è presente, annotare ciò chiaramente in
      modo che i programmi che aggiornano il database
      operino correttamente.

-    Talvolta le linee guida possono essere violate per
    scopi di efficienza. Una soluzione migliore
    potrebbe comunque essere quella di definire
    delle viste.
\end{lstlisting}
\subsection*{Spiegazione estesa}
[AGGIUNTA] Questa spiegazione collega il contenuto testuale della slide al contesto del corso. La slide enfatizza i concetti mostrati nell'estratto e il loro ruolo nella progettazione di basi di dati relazionali, con particolare attenzione a correttezza semantica, qualità dello schema e compromessi progettuali.
[AGGIUNTA] Fonte integrazione: riformulazione didattica del testo presente nella slide 26 del file Lezione 01 - Dipendenze funzionali e Normalizzazione di DB Relazionali.pdf.
\subsection*{Derivazioni}
[AGGIUNTA] In questa slide non sono state rilevate equazioni esplicite nel testo estratto; non è quindi richiesta una derivazione algebrica puntuale.
\subsection*{Esempio numerico / esercizio}
[AGGIUNTA] Esempio: dato uno schema R(A,B,C) con dipendenze funzionali A->B e B->C, mostrare la chiusura A^+ e discutere se A è superchiave. Soluzione sintetica: A^+={A,B,C}, quindi A determina tutti gli attributi di R ed è superchiave.
\subsection*{Note su assunzioni o integrazioni}
[ASSUNZIONE] L'estrazione OCR può perdere formattazione o simboli della slide originale.
[AGGIUNTA] Le spiegazioni sono estensioni didattiche per rendere espliciti passaggi impliciti nel materiale proiettato.
\paragraph{Riferimento fonte} File: Lezione 01 - Dipendenze funzionali e Normalizzazione di DB Relazionali.pdf. Numero slide: 26.
\section{Outline --- Slide 27}
\subsection*{Estratto originale dalla slide 27}
\begin{lstlisting}
Outline
  1. Semantica degli attributi.
  2. Riduzione dei valori ridondanti nelle tuple.

  3. Riduzione dei valori null nelle tuple.

  4. Non consentire tuple spurie.
\end{lstlisting}
\subsection*{Spiegazione estesa}
[AGGIUNTA] Questa spiegazione collega il contenuto testuale della slide al contesto del corso. La slide enfatizza i concetti mostrati nell'estratto e il loro ruolo nella progettazione di basi di dati relazionali, con particolare attenzione a correttezza semantica, qualità dello schema e compromessi progettuali.
[AGGIUNTA] Fonte integrazione: riformulazione didattica del testo presente nella slide 27 del file Lezione 01 - Dipendenze funzionali e Normalizzazione di DB Relazionali.pdf.
\subsection*{Derivazioni}
[AGGIUNTA] In questa slide non sono state rilevate equazioni esplicite nel testo estratto; non è quindi richiesta una derivazione algebrica puntuale.
\subsection*{Esempio numerico / esercizio}
[AGGIUNTA] Esempio: dato uno schema R(A,B,C) con dipendenze funzionali A->B e B->C, mostrare la chiusura A^+ e discutere se A è superchiave. Soluzione sintetica: A^+={A,B,C}, quindi A determina tutti gli attributi di R ed è superchiave.
\subsection*{Note su assunzioni o integrazioni}
[ASSUNZIONE] L'estrazione OCR può perdere formattazione o simboli della slide originale.
[AGGIUNTA] Le spiegazioni sono estensioni didattiche per rendere espliciti passaggi impliciti nel materiale proiettato.
\paragraph{Riferimento fonte} File: Lezione 01 - Dipendenze funzionali e Normalizzazione di DB Relazionali.pdf. Numero slide: 27.
\section{3) Riduzione dei valori null nelle tuple --- Slide 28}
\subsection*{Estratto originale dalla slide 28}
\begin{lstlisting}
3) Riduzione dei valori null nelle tuple

-    In alcuni disegni di schemi possiamo raggruppare
    molti attributi in una relazione grossa. Se molti
    degli attributi non applicano a tutte le tuple
    possiamo avere molti valori null.

-    Oltre allo spreco di spazio, ciò crea problemi con
    le funzioni di aggregazione COUNT e SUM.
\end{lstlisting}
\subsection*{Spiegazione estesa}
[AGGIUNTA] Questa spiegazione collega il contenuto testuale della slide al contesto del corso. La slide enfatizza i concetti mostrati nell'estratto e il loro ruolo nella progettazione di basi di dati relazionali, con particolare attenzione a correttezza semantica, qualità dello schema e compromessi progettuali.
[AGGIUNTA] Fonte integrazione: riformulazione didattica del testo presente nella slide 28 del file Lezione 01 - Dipendenze funzionali e Normalizzazione di DB Relazionali.pdf.
\subsection*{Derivazioni}
[AGGIUNTA] In questa slide non sono state rilevate equazioni esplicite nel testo estratto; non è quindi richiesta una derivazione algebrica puntuale.
\subsection*{Esempio numerico / esercizio}
[AGGIUNTA] Esempio: dato uno schema R(A,B,C) con dipendenze funzionali A->B e B->C, mostrare la chiusura A^+ e discutere se A è superchiave. Soluzione sintetica: A^+={A,B,C}, quindi A determina tutti gli attributi di R ed è superchiave.
\subsection*{Note su assunzioni o integrazioni}
[ASSUNZIONE] L'estrazione OCR può perdere formattazione o simboli della slide originale.
[AGGIUNTA] Le spiegazioni sono estensioni didattiche per rendere espliciti passaggi impliciti nel materiale proiettato.
\paragraph{Riferimento fonte} File: Lezione 01 - Dipendenze funzionali e Normalizzazione di DB Relazionali.pdf. Numero slide: 28.
\section{Interpretazioni di null --- Slide 29}
\subsection*{Estratto originale dalla slide 29}
\begin{lstlisting}
Interpretazioni di null
-    Il valore null può avere diverse interpretazioni:
      L'attributo non si applica a questa tupla.

      Il valore dell'attributo per questa tupla non è noto.

      Il valore dell'attributo è noto ma assente, cioè non è
      stato ancora registrato.
\end{lstlisting}
\subsection*{Spiegazione estesa}
[AGGIUNTA] Questa spiegazione collega il contenuto testuale della slide al contesto del corso. La slide enfatizza i concetti mostrati nell'estratto e il loro ruolo nella progettazione di basi di dati relazionali, con particolare attenzione a correttezza semantica, qualità dello schema e compromessi progettuali.
[AGGIUNTA] Fonte integrazione: riformulazione didattica del testo presente nella slide 29 del file Lezione 01 - Dipendenze funzionali e Normalizzazione di DB Relazionali.pdf.
\subsection*{Derivazioni}
[AGGIUNTA] In questa slide non sono state rilevate equazioni esplicite nel testo estratto; non è quindi richiesta una derivazione algebrica puntuale.
\subsection*{Esempio numerico / esercizio}
[AGGIUNTA] Esempio: dato uno schema R(A,B,C) con dipendenze funzionali A->B e B->C, mostrare la chiusura A^+ e discutere se A è superchiave. Soluzione sintetica: A^+={A,B,C}, quindi A determina tutti gli attributi di R ed è superchiave.
\subsection*{Note su assunzioni o integrazioni}
[ASSUNZIONE] L'estrazione OCR può perdere formattazione o simboli della slide originale.
[AGGIUNTA] Le spiegazioni sono estensioni didattiche per rendere espliciti passaggi impliciti nel materiale proiettato.
\paragraph{Riferimento fonte} File: Lezione 01 - Dipendenze funzionali e Normalizzazione di DB Relazionali.pdf. Numero slide: 29.
\section{Linea guida 3 --- Slide 30}
\subsection*{Estratto originale dalla slide 30}
\begin{lstlisting}
Linea guida 3
-    Evitare (per quanto possibile) di porre attributi in
    una relazione base i cui valori possono essere
    null.

-    Se i valori null sono inevitabili, assicurarsi che
    essi ricorrano in casi eccezionali e non per la
    maggioranza delle tuple.
\end{lstlisting}
\subsection*{Spiegazione estesa}
[AGGIUNTA] Questa spiegazione collega il contenuto testuale della slide al contesto del corso. La slide enfatizza i concetti mostrati nell'estratto e il loro ruolo nella progettazione di basi di dati relazionali, con particolare attenzione a correttezza semantica, qualità dello schema e compromessi progettuali.
[AGGIUNTA] Fonte integrazione: riformulazione didattica del testo presente nella slide 30 del file Lezione 01 - Dipendenze funzionali e Normalizzazione di DB Relazionali.pdf.
\subsection*{Derivazioni}
[AGGIUNTA] In questa slide non sono state rilevate equazioni esplicite nel testo estratto; non è quindi richiesta una derivazione algebrica puntuale.
\subsection*{Esempio numerico / esercizio}
[AGGIUNTA] Esempio: dato uno schema R(A,B,C) con dipendenze funzionali A->B e B->C, mostrare la chiusura A^+ e discutere se A è superchiave. Soluzione sintetica: A^+={A,B,C}, quindi A determina tutti gli attributi di R ed è superchiave.
\subsection*{Note su assunzioni o integrazioni}
[ASSUNZIONE] L'estrazione OCR può perdere formattazione o simboli della slide originale.
[AGGIUNTA] Le spiegazioni sono estensioni didattiche per rendere espliciti passaggi impliciti nel materiale proiettato.
\paragraph{Riferimento fonte} File: Lezione 01 - Dipendenze funzionali e Normalizzazione di DB Relazionali.pdf. Numero slide: 30.
\section{Linea guida 3: Esempio --- Slide 31}
\subsection*{Estratto originale dalla slide 31}
\begin{lstlisting}
Linea guida 3: Esempio
-    Se solo il 10% degli impiegati ha un ufficio
    individuale, è poco giustificato include l'attributo
    OFFICE_NUMBER nella relazione EMPLOYEE;

    piuttosto creare la relazione:
    EMP_OFFICES (ESSN, OFFICE_NUMBER)
\end{lstlisting}
\subsection*{Spiegazione estesa}
[AGGIUNTA] Questa spiegazione collega il contenuto testuale della slide al contesto del corso. La slide enfatizza i concetti mostrati nell'estratto e il loro ruolo nella progettazione di basi di dati relazionali, con particolare attenzione a correttezza semantica, qualità dello schema e compromessi progettuali.
[AGGIUNTA] Fonte integrazione: riformulazione didattica del testo presente nella slide 31 del file Lezione 01 - Dipendenze funzionali e Normalizzazione di DB Relazionali.pdf.
\subsection*{Derivazioni}
[AGGIUNTA] In questa slide non sono state rilevate equazioni esplicite nel testo estratto; non è quindi richiesta una derivazione algebrica puntuale.
\subsection*{Esempio numerico / esercizio}
[AGGIUNTA] Esempio: dato uno schema R(A,B,C) con dipendenze funzionali A->B e B->C, mostrare la chiusura A^+ e discutere se A è superchiave. Soluzione sintetica: A^+={A,B,C}, quindi A determina tutti gli attributi di R ed è superchiave.
\subsection*{Note su assunzioni o integrazioni}
[ASSUNZIONE] L'estrazione OCR può perdere formattazione o simboli della slide originale.
[AGGIUNTA] Le spiegazioni sono estensioni didattiche per rendere espliciti passaggi impliciti nel materiale proiettato.
\paragraph{Riferimento fonte} File: Lezione 01 - Dipendenze funzionali e Normalizzazione di DB Relazionali.pdf. Numero slide: 31.
\section{Outline --- Slide 32}
\subsection*{Estratto originale dalla slide 32}
\begin{lstlisting}
Outline
  1. Semantica degli attributi.
  2. Riduzione dei valori ridondanti nelle tuple.

  3. Riduzione dei valori null nelle tuple.

  4. Non consentire tuple spurie.
\end{lstlisting}
\subsection*{Spiegazione estesa}
[AGGIUNTA] Questa spiegazione collega il contenuto testuale della slide al contesto del corso. La slide enfatizza i concetti mostrati nell'estratto e il loro ruolo nella progettazione di basi di dati relazionali, con particolare attenzione a correttezza semantica, qualità dello schema e compromessi progettuali.
[AGGIUNTA] Fonte integrazione: riformulazione didattica del testo presente nella slide 32 del file Lezione 01 - Dipendenze funzionali e Normalizzazione di DB Relazionali.pdf.
\subsection*{Derivazioni}
[AGGIUNTA] In questa slide non sono state rilevate equazioni esplicite nel testo estratto; non è quindi richiesta una derivazione algebrica puntuale.
\subsection*{Esempio numerico / esercizio}
[AGGIUNTA] Esempio: dato uno schema R(A,B,C) con dipendenze funzionali A->B e B->C, mostrare la chiusura A^+ e discutere se A è superchiave. Soluzione sintetica: A^+={A,B,C}, quindi A determina tutti gli attributi di R ed è superchiave.
\subsection*{Note su assunzioni o integrazioni}
[ASSUNZIONE] L'estrazione OCR può perdere formattazione o simboli della slide originale.
[AGGIUNTA] Le spiegazioni sono estensioni didattiche per rendere espliciti passaggi impliciti nel materiale proiettato.
\paragraph{Riferimento fonte} File: Lezione 01 - Dipendenze funzionali e Normalizzazione di DB Relazionali.pdf. Numero slide: 32.
\section{Tuple spurie --- Slide 33}
\subsection*{Estratto originale dalla slide 33}
\begin{lstlisting}
Tuple spurie
-    Consideriamo i due schemi, alternativi a EMP_PROJ:




-    Le tuple rappresentano:
      EMP_LOCS: l'impiegato di nome ENAME lavora per qualche
      progetto la cui locazione è PLOCATION.
      EMP_PROJ1: l'impiegato con codice SSN lavora HOURS ore alla
      settimana sul progetto avente numero PNUMBER, nome PNAME
      e locazione PLOCATION.
\end{lstlisting}
\subsection*{Spiegazione estesa}
[AGGIUNTA] Questa spiegazione collega il contenuto testuale della slide al contesto del corso. La slide enfatizza i concetti mostrati nell'estratto e il loro ruolo nella progettazione di basi di dati relazionali, con particolare attenzione a correttezza semantica, qualità dello schema e compromessi progettuali.
[AGGIUNTA] Fonte integrazione: riformulazione didattica del testo presente nella slide 33 del file Lezione 01 - Dipendenze funzionali e Normalizzazione di DB Relazionali.pdf.
\subsection*{Derivazioni}
[AGGIUNTA] In questa slide non sono state rilevate equazioni esplicite nel testo estratto; non è quindi richiesta una derivazione algebrica puntuale.
\subsection*{Esempio numerico / esercizio}
[AGGIUNTA] Esempio: dato uno schema R(A,B,C) con dipendenze funzionali A->B e B->C, mostrare la chiusura A^+ e discutere se A è superchiave. Soluzione sintetica: A^+={A,B,C}, quindi A determina tutti gli attributi di R ed è superchiave.
\subsection*{Note su assunzioni o integrazioni}
[ASSUNZIONE] L'estrazione OCR può perdere formattazione o simboli della slide originale.
[AGGIUNTA] Le spiegazioni sono estensioni didattiche per rendere espliciti passaggi impliciti nel materiale proiettato.
\paragraph{Riferimento fonte} File: Lezione 01 - Dipendenze funzionali e Normalizzazione di DB Relazionali.pdf. Numero slide: 33.
\section{Tuple spurie (2) --- Slide 34}
\subsection*{Estratto originale dalla slide 34}
\begin{lstlisting}
Tuple spurie (2)
-    La suddivisione appena vista non corrisponde ad
    un buon disegno di db in quanto, volendo
    riottenere le informazioni di EMP_PROJ mediante
    un natural join sui due schemi, si ottengono tuple
    spurie, rappresentanti informazioni errate.

EMP_LOCS * EMP_PROJ1
\end{lstlisting}
\subsection*{Spiegazione estesa}
[AGGIUNTA] Questa spiegazione collega il contenuto testuale della slide al contesto del corso. La slide enfatizza i concetti mostrati nell'estratto e il loro ruolo nella progettazione di basi di dati relazionali, con particolare attenzione a correttezza semantica, qualità dello schema e compromessi progettuali.
[AGGIUNTA] Fonte integrazione: riformulazione didattica del testo presente nella slide 34 del file Lezione 01 - Dipendenze funzionali e Normalizzazione di DB Relazionali.pdf.
\subsection*{Derivazioni}
[AGGIUNTA] In questa slide non sono state rilevate equazioni esplicite nel testo estratto; non è quindi richiesta una derivazione algebrica puntuale.
\subsection*{Esempio numerico / esercizio}
[AGGIUNTA] Esempio: dato uno schema R(A,B,C) con dipendenze funzionali A->B e B->C, mostrare la chiusura A^+ e discutere se A è superchiave. Soluzione sintetica: A^+={A,B,C}, quindi A determina tutti gli attributi di R ed è superchiave.
\subsection*{Note su assunzioni o integrazioni}
[ASSUNZIONE] L'estrazione OCR può perdere formattazione o simboli della slide originale.
[AGGIUNTA] Le spiegazioni sono estensioni didattiche per rendere espliciti passaggi impliciti nel materiale proiettato.
\paragraph{Riferimento fonte} File: Lezione 01 - Dipendenze funzionali e Normalizzazione di DB Relazionali.pdf. Numero slide: 34.
\section{Tuple spurie: Esempio --- Slide 35}
\subsection*{Estratto originale dalla slide 35}
\begin{lstlisting}
Tuple spurie: Esempio




             I valori presenti nelle due relazioni
\end{lstlisting}
\subsection*{Spiegazione estesa}
[AGGIUNTA] Questa spiegazione collega il contenuto testuale della slide al contesto del corso. La slide enfatizza i concetti mostrati nell'estratto e il loro ruolo nella progettazione di basi di dati relazionali, con particolare attenzione a correttezza semantica, qualità dello schema e compromessi progettuali.
[AGGIUNTA] Fonte integrazione: riformulazione didattica del testo presente nella slide 35 del file Lezione 01 - Dipendenze funzionali e Normalizzazione di DB Relazionali.pdf.
\subsection*{Derivazioni}
[AGGIUNTA] In questa slide non sono state rilevate equazioni esplicite nel testo estratto; non è quindi richiesta una derivazione algebrica puntuale.
\subsection*{Esempio numerico / esercizio}
[AGGIUNTA] Esempio: dato uno schema R(A,B,C) con dipendenze funzionali A->B e B->C, mostrare la chiusura A^+ e discutere se A è superchiave. Soluzione sintetica: A^+={A,B,C}, quindi A determina tutti gli attributi di R ed è superchiave.
\subsection*{Note su assunzioni o integrazioni}
[ASSUNZIONE] L'estrazione OCR può perdere formattazione o simboli della slide originale.
[AGGIUNTA] Le spiegazioni sono estensioni didattiche per rendere espliciti passaggi impliciti nel materiale proiettato.
\paragraph{Riferimento fonte} File: Lezione 01 - Dipendenze funzionali e Normalizzazione di DB Relazionali.pdf. Numero slide: 35.
\section{Tuple spurie: Esempio (2) --- Slide 36}
\subsection*{Estratto originale dalla slide 36}
\begin{lstlisting}
Tuple spurie: Esempio (2)




      Risultato del natural join:
      le tuple spurie sono contrassegnate da un '*'
\end{lstlisting}
\subsection*{Spiegazione estesa}
[AGGIUNTA] Questa spiegazione collega il contenuto testuale della slide al contesto del corso. La slide enfatizza i concetti mostrati nell'estratto e il loro ruolo nella progettazione di basi di dati relazionali, con particolare attenzione a correttezza semantica, qualità dello schema e compromessi progettuali.
[AGGIUNTA] Fonte integrazione: riformulazione didattica del testo presente nella slide 36 del file Lezione 01 - Dipendenze funzionali e Normalizzazione di DB Relazionali.pdf.
\subsection*{Derivazioni}
[AGGIUNTA] In questa slide non sono state rilevate equazioni esplicite nel testo estratto; non è quindi richiesta una derivazione algebrica puntuale.
\subsection*{Esempio numerico / esercizio}
[AGGIUNTA] Esempio: dato uno schema R(A,B,C) con dipendenze funzionali A->B e B->C, mostrare la chiusura A^+ e discutere se A è superchiave. Soluzione sintetica: A^+={A,B,C}, quindi A determina tutti gli attributi di R ed è superchiave.
\subsection*{Note su assunzioni o integrazioni}
[ASSUNZIONE] L'estrazione OCR può perdere formattazione o simboli della slide originale.
[AGGIUNTA] Le spiegazioni sono estensioni didattiche per rendere espliciti passaggi impliciti nel materiale proiettato.
\paragraph{Riferimento fonte} File: Lezione 01 - Dipendenze funzionali e Normalizzazione di DB Relazionali.pdf. Numero slide: 36.
\section{Motivo delle tuple spurie --- Slide 37}
\subsection*{Estratto originale dalla slide 37}
\begin{lstlisting}
Motivo delle tuple spurie
-    Tale problema sorge perché PLOCATION è
    l'attributo che correla EMP_LOCS e
    EMP_PROJ1, ma non è né chiave primaria, né
    chiave esterna in nessuna delle due relazioni.
\end{lstlisting}
\subsection*{Spiegazione estesa}
[AGGIUNTA] Questa spiegazione collega il contenuto testuale della slide al contesto del corso. La slide enfatizza i concetti mostrati nell'estratto e il loro ruolo nella progettazione di basi di dati relazionali, con particolare attenzione a correttezza semantica, qualità dello schema e compromessi progettuali.
[AGGIUNTA] Fonte integrazione: riformulazione didattica del testo presente nella slide 37 del file Lezione 01 - Dipendenze funzionali e Normalizzazione di DB Relazionali.pdf.
\subsection*{Derivazioni}
[AGGIUNTA] In questa slide non sono state rilevate equazioni esplicite nel testo estratto; non è quindi richiesta una derivazione algebrica puntuale.
\subsection*{Esempio numerico / esercizio}
[AGGIUNTA] Esempio: dato uno schema R(A,B,C) con dipendenze funzionali A->B e B->C, mostrare la chiusura A^+ e discutere se A è superchiave. Soluzione sintetica: A^+={A,B,C}, quindi A determina tutti gli attributi di R ed è superchiave.
\subsection*{Note su assunzioni o integrazioni}
[ASSUNZIONE] L'estrazione OCR può perdere formattazione o simboli della slide originale.
[AGGIUNTA] Le spiegazioni sono estensioni didattiche per rendere espliciti passaggi impliciti nel materiale proiettato.
\paragraph{Riferimento fonte} File: Lezione 01 - Dipendenze funzionali e Normalizzazione di DB Relazionali.pdf. Numero slide: 37.
\section{Linea guida 4 --- Slide 38}
\subsection*{Estratto originale dalla slide 38}
\begin{lstlisting}
Linea guida 4
-    Progettare gli schemi di relazione in modo da
    poter effettuare JOIN con condizioni di
    uguaglianza su attributi che sono o chiave
    primaria o chiave esterna, in modo da non
    generare tuple spurie.
\end{lstlisting}
\subsection*{Spiegazione estesa}
[AGGIUNTA] Questa spiegazione collega il contenuto testuale della slide al contesto del corso. La slide enfatizza i concetti mostrati nell'estratto e il loro ruolo nella progettazione di basi di dati relazionali, con particolare attenzione a correttezza semantica, qualità dello schema e compromessi progettuali.
[AGGIUNTA] Fonte integrazione: riformulazione didattica del testo presente nella slide 38 del file Lezione 01 - Dipendenze funzionali e Normalizzazione di DB Relazionali.pdf.
\subsection*{Derivazioni}
[AGGIUNTA] In questa slide non sono state rilevate equazioni esplicite nel testo estratto; non è quindi richiesta una derivazione algebrica puntuale.
\subsection*{Esempio numerico / esercizio}
[AGGIUNTA] Esempio: dato uno schema R(A,B,C) con dipendenze funzionali A->B e B->C, mostrare la chiusura A^+ e discutere se A è superchiave. Soluzione sintetica: A^+={A,B,C}, quindi A determina tutti gli attributi di R ed è superchiave.
\subsection*{Note su assunzioni o integrazioni}
[ASSUNZIONE] L'estrazione OCR può perdere formattazione o simboli della slide originale.
[AGGIUNTA] Le spiegazioni sono estensioni didattiche per rendere espliciti passaggi impliciti nel materiale proiettato.
\paragraph{Riferimento fonte} File: Lezione 01 - Dipendenze funzionali e Normalizzazione di DB Relazionali.pdf. Numero slide: 38.
\section{Sommario delle linee guida --- Slide 39}
\subsection*{Estratto originale dalla slide 39}
\begin{lstlisting}
Sommario delle linee guida
-    Abbiamo mostrato in modo informale che una
    cattiva progettazione dello schema di un db può
    portare ad una serie di problemi:
      Anomalie che implicano un maggiore sforzo
      nell'inserimento e nella modifica di una relazione, e che
      possono portare a perdite accidentali di dati durante la
      cancellazione.
      Spreco di spazio a causa dei valori null.
      Generazione di tuple spurie o non valide durante
      operazioni di join.
\end{lstlisting}
\subsection*{Spiegazione estesa}
[AGGIUNTA] Questa spiegazione collega il contenuto testuale della slide al contesto del corso. La slide enfatizza i concetti mostrati nell'estratto e il loro ruolo nella progettazione di basi di dati relazionali, con particolare attenzione a correttezza semantica, qualità dello schema e compromessi progettuali.
[AGGIUNTA] Fonte integrazione: riformulazione didattica del testo presente nella slide 39 del file Lezione 01 - Dipendenze funzionali e Normalizzazione di DB Relazionali.pdf.
\subsection*{Derivazioni}
[AGGIUNTA] In questa slide non sono state rilevate equazioni esplicite nel testo estratto; non è quindi richiesta una derivazione algebrica puntuale.
\subsection*{Esempio numerico / esercizio}
[AGGIUNTA] Esempio: dato uno schema R(A,B,C) con dipendenze funzionali A->B e B->C, mostrare la chiusura A^+ e discutere se A è superchiave. Soluzione sintetica: A^+={A,B,C}, quindi A determina tutti gli attributi di R ed è superchiave.
\subsection*{Note su assunzioni o integrazioni}
[ASSUNZIONE] L'estrazione OCR può perdere formattazione o simboli della slide originale.
[AGGIUNTA] Le spiegazioni sono estensioni didattiche per rendere espliciti passaggi impliciti nel materiale proiettato.
\paragraph{Riferimento fonte} File: Lezione 01 - Dipendenze funzionali e Normalizzazione di DB Relazionali.pdf. Numero slide: 39.
\section{Dipendenze Funzionali --- Slide 40}
\subsection*{Estratto originale dalla slide 40}
\begin{lstlisting}
Dipendenze Funzionali
\end{lstlisting}
\subsection*{Spiegazione estesa}
[AGGIUNTA] Questa spiegazione collega il contenuto testuale della slide al contesto del corso. La slide enfatizza i concetti mostrati nell'estratto e il loro ruolo nella progettazione di basi di dati relazionali, con particolare attenzione a correttezza semantica, qualità dello schema e compromessi progettuali.
[AGGIUNTA] Fonte integrazione: riformulazione didattica del testo presente nella slide 40 del file Lezione 01 - Dipendenze funzionali e Normalizzazione di DB Relazionali.pdf.
\subsection*{Derivazioni}
[AGGIUNTA] In questa slide non sono state rilevate equazioni esplicite nel testo estratto; non è quindi richiesta una derivazione algebrica puntuale.
\subsection*{Esempio numerico / esercizio}
[AGGIUNTA] Esempio: dato uno schema R(A,B,C) con dipendenze funzionali A->B e B->C, mostrare la chiusura A^+ e discutere se A è superchiave. Soluzione sintetica: A^+={A,B,C}, quindi A determina tutti gli attributi di R ed è superchiave.
\subsection*{Note su assunzioni o integrazioni}
[ASSUNZIONE] L'estrazione OCR può perdere formattazione o simboli della slide originale.
[AGGIUNTA] Le spiegazioni sono estensioni didattiche per rendere espliciti passaggi impliciti nel materiale proiettato.
\paragraph{Riferimento fonte} File: Lezione 01 - Dipendenze funzionali e Normalizzazione di DB Relazionali.pdf. Numero slide: 40.
\section{Qualità di schemi relazionali --- Slide 41}
\subsection*{Estratto originale dalla slide 41}
\begin{lstlisting}
Qualità di schemi relazionali
-    Dopo una visione informale, vediamo delle teorie
    e dei concetti formali per descrivere la "bontà" dei
    singoli schemi di relazione con maggiore
    precisione, usando le "dipendenze funzionali" e le
    "forme normali".
\end{lstlisting}
\subsection*{Spiegazione estesa}
[AGGIUNTA] Questa spiegazione collega il contenuto testuale della slide al contesto del corso. La slide enfatizza i concetti mostrati nell'estratto e il loro ruolo nella progettazione di basi di dati relazionali, con particolare attenzione a correttezza semantica, qualità dello schema e compromessi progettuali.
[AGGIUNTA] Fonte integrazione: riformulazione didattica del testo presente nella slide 41 del file Lezione 01 - Dipendenze funzionali e Normalizzazione di DB Relazionali.pdf.
\subsection*{Derivazioni}
[AGGIUNTA] In questa slide non sono state rilevate equazioni esplicite nel testo estratto; non è quindi richiesta una derivazione algebrica puntuale.
\subsection*{Esempio numerico / esercizio}
[AGGIUNTA] Esempio: dato uno schema R(A,B,C) con dipendenze funzionali A->B e B->C, mostrare la chiusura A^+ e discutere se A è superchiave. Soluzione sintetica: A^+={A,B,C}, quindi A determina tutti gli attributi di R ed è superchiave.
\subsection*{Note su assunzioni o integrazioni}
[ASSUNZIONE] L'estrazione OCR può perdere formattazione o simboli della slide originale.
[AGGIUNTA] Le spiegazioni sono estensioni didattiche per rendere espliciti passaggi impliciti nel materiale proiettato.
\paragraph{Riferimento fonte} File: Lezione 01 - Dipendenze funzionali e Normalizzazione di DB Relazionali.pdf. Numero slide: 41.
\section{Dipendenza Funzionale --- Slide 42}
\subsection*{Estratto originale dalla slide 42}
\begin{lstlisting}
Dipendenza Funzionale
-    Una dipendenza funzionale (FD) è un vincolo tra due
    insiemi di attributi del database.

-    Supponiamo che lo schema di db relazionale abbia n
    attributi A1, A2, ..., An e che l'intero database sia descritto
    da uno schema di relazione universale
                R = {A1, A2, ..., An}.

-    Una dipendenza funzionale, denotata da X->Y, tra due
    insiemi di attributi X e Y che sono sottoinsiemi di R,
    specifica un vincolo sulle possibili tuple che possono
    formare una istanza di relazione r di R.
\end{lstlisting}
\subsection*{Spiegazione estesa}
[AGGIUNTA] Questa spiegazione collega il contenuto testuale della slide al contesto del corso. La slide enfatizza i concetti mostrati nell'estratto e il loro ruolo nella progettazione di basi di dati relazionali, con particolare attenzione a correttezza semantica, qualità dello schema e compromessi progettuali.
[AGGIUNTA] Fonte integrazione: riformulazione didattica del testo presente nella slide 42 del file Lezione 01 - Dipendenze funzionali e Normalizzazione di DB Relazionali.pdf.
\subsection*{Derivazioni}
[AGGIUNTA] Nella slide compaiono espressioni simboliche/relazioni. Derivazione didattica generale:
\begin{align*}
X^+ &= X \\
X^+ &\leftarrow X^+ \cup Y \quad \text{se } (U \to V) \in F \text{ e } U \subseteq X^+ \text{ con } Y=V \\
\text{Stop} &\text{ quando non si aggiungono nuovi attributi in } X^+
\end{align*}
[AGGIUNTA] Questa procedura esplicita il metodo di chiusura usato nelle slide su dipendenze e chiavi.
\begin{itemize}[leftmargin=*]
\item Forma osservata nella slide: \texttt{R = (A1, A2, ..., An).}
\item Forma osservata nella slide: \texttt{-    Una dipendenza funzionale, denotata da X->Y, tra due}
\end{itemize}
\subsection*{Esempio numerico / esercizio}
[AGGIUNTA] Esempio: dato uno schema R(A,B,C) con dipendenze funzionali A->B e B->C, mostrare la chiusura A^+ e discutere se A è superchiave. Soluzione sintetica: A^+={A,B,C}, quindi A determina tutti gli attributi di R ed è superchiave.
\subsection*{Note su assunzioni o integrazioni}
[ASSUNZIONE] L'estrazione OCR può perdere formattazione o simboli della slide originale.
[AGGIUNTA] Le spiegazioni sono estensioni didattiche per rendere espliciti passaggi impliciti nel materiale proiettato.
\paragraph{Riferimento fonte} File: Lezione 01 - Dipendenze funzionali e Normalizzazione di DB Relazionali.pdf. Numero slide: 42.
\section{Dipendenza Funzionale (2) --- Slide 43}
\subsection*{Estratto originale dalla slide 43}
\begin{lstlisting}
Dipendenza Funzionale (2)
-    Il vincolo stabilisce che se X -> Y, allora
      t1, t2 in r tali che
                t1[X] = t2[X],
    deve valere
                t1[Y] = t2[Y].

      Ciò significa che i valori della componente Y di una
      tupla di r dipendono da (o sono determinati da) i valori
      della componente X.
\end{lstlisting}
\subsection*{Spiegazione estesa}
[AGGIUNTA] Questa spiegazione collega il contenuto testuale della slide al contesto del corso. La slide enfatizza i concetti mostrati nell'estratto e il loro ruolo nella progettazione di basi di dati relazionali, con particolare attenzione a correttezza semantica, qualità dello schema e compromessi progettuali.
[AGGIUNTA] Fonte integrazione: riformulazione didattica del testo presente nella slide 43 del file Lezione 01 - Dipendenze funzionali e Normalizzazione di DB Relazionali.pdf.
\subsection*{Derivazioni}
[AGGIUNTA] Nella slide compaiono espressioni simboliche/relazioni. Derivazione didattica generale:
\begin{align*}
X^+ &= X \\
X^+ &\leftarrow X^+ \cup Y \quad \text{se } (U \to V) \in F \text{ e } U \subseteq X^+ \text{ con } Y=V \\
\text{Stop} &\text{ quando non si aggiungono nuovi attributi in } X^+
\end{align*}
[AGGIUNTA] Questa procedura esplicita il metodo di chiusura usato nelle slide su dipendenze e chiavi.
\begin{itemize}[leftmargin=*]
\item Forma osservata nella slide: \texttt{-    Il vincolo stabilisce che se X -> Y, allora}
\item Forma osservata nella slide: \texttt{t1[X] = t2[X],}
\item Forma osservata nella slide: \texttt{t1[Y] = t2[Y].}
\end{itemize}
\subsection*{Esempio numerico / esercizio}
[AGGIUNTA] Esempio: dato uno schema R(A,B,C) con dipendenze funzionali A->B e B->C, mostrare la chiusura A^+ e discutere se A è superchiave. Soluzione sintetica: A^+={A,B,C}, quindi A determina tutti gli attributi di R ed è superchiave.
\subsection*{Note su assunzioni o integrazioni}
[ASSUNZIONE] L'estrazione OCR può perdere formattazione o simboli della slide originale.
[AGGIUNTA] Le spiegazioni sono estensioni didattiche per rendere espliciti passaggi impliciti nel materiale proiettato.
\paragraph{Riferimento fonte} File: Lezione 01 - Dipendenze funzionali e Normalizzazione di DB Relazionali.pdf. Numero slide: 43.
\section{Dipendenza Funzionale (3) --- Slide 44}
\subsection*{Estratto originale dalla slide 44}
\begin{lstlisting}
Dipendenza Funzionale (3)
-    Alternativamente, i valori della componente X di
    una tupla determinano univocamente (o
    funzionalmente) i valori della componente Y.

-    Cioè esiste una dipendenza funzionale da X a Y:
      Y è funzionalmente dipendente da X

      X è la parte sinistra della FD

      Y è la parte destra della FD
\end{lstlisting}
\subsection*{Spiegazione estesa}
[AGGIUNTA] Questa spiegazione collega il contenuto testuale della slide al contesto del corso. La slide enfatizza i concetti mostrati nell'estratto e il loro ruolo nella progettazione di basi di dati relazionali, con particolare attenzione a correttezza semantica, qualità dello schema e compromessi progettuali.
[AGGIUNTA] Fonte integrazione: riformulazione didattica del testo presente nella slide 44 del file Lezione 01 - Dipendenze funzionali e Normalizzazione di DB Relazionali.pdf.
\subsection*{Derivazioni}
[AGGIUNTA] In questa slide non sono state rilevate equazioni esplicite nel testo estratto; non è quindi richiesta una derivazione algebrica puntuale.
\subsection*{Esempio numerico / esercizio}
[AGGIUNTA] Esempio: dato uno schema R(A,B,C) con dipendenze funzionali A->B e B->C, mostrare la chiusura A^+ e discutere se A è superchiave. Soluzione sintetica: A^+={A,B,C}, quindi A determina tutti gli attributi di R ed è superchiave.
\subsection*{Note su assunzioni o integrazioni}
[ASSUNZIONE] L'estrazione OCR può perdere formattazione o simboli della slide originale.
[AGGIUNTA] Le spiegazioni sono estensioni didattiche per rendere espliciti passaggi impliciti nel materiale proiettato.
\paragraph{Riferimento fonte} File: Lezione 01 - Dipendenze funzionali e Normalizzazione di DB Relazionali.pdf. Numero slide: 44.
\section{Dipendenza Funzionale (4) --- Slide 45}
\subsection*{Estratto originale dalla slide 45}
\begin{lstlisting}
Dipendenza Funzionale (4)
-    Si noti che:
      Se un vincolo su R stabilisce che non può esistere più di
      una tupla con un dato valore per X in ogni istanza di
      relazione r(R): "cioè X è una chiave candidata di R"; ciò
      implica che X -> Y per ogni sottoinsieme Y di attributi di R.

      Il fatto che X -> Y in R non implica nulla circa Y -> X in R.
\end{lstlisting}
\subsection*{Spiegazione estesa}
[AGGIUNTA] Questa spiegazione collega il contenuto testuale della slide al contesto del corso. La slide enfatizza i concetti mostrati nell'estratto e il loro ruolo nella progettazione di basi di dati relazionali, con particolare attenzione a correttezza semantica, qualità dello schema e compromessi progettuali.
[AGGIUNTA] Fonte integrazione: riformulazione didattica del testo presente nella slide 45 del file Lezione 01 - Dipendenze funzionali e Normalizzazione di DB Relazionali.pdf.
\subsection*{Derivazioni}
[AGGIUNTA] Nella slide compaiono espressioni simboliche/relazioni. Derivazione didattica generale:
\begin{align*}
X^+ &= X \\
X^+ &\leftarrow X^+ \cup Y \quad \text{se } (U \to V) \in F \text{ e } U \subseteq X^+ \text{ con } Y=V \\
\text{Stop} &\text{ quando non si aggiungono nuovi attributi in } X^+
\end{align*}
[AGGIUNTA] Questa procedura esplicita il metodo di chiusura usato nelle slide su dipendenze e chiavi.
\begin{itemize}[leftmargin=*]
\item Forma osservata nella slide: \texttt{implica che X -> Y per ogni sottoinsieme Y di attributi di R.}
\item Forma osservata nella slide: \texttt{Il fatto che X -> Y in R non implica nulla circa Y -> X in R.}
\end{itemize}
\subsection*{Esempio numerico / esercizio}
[AGGIUNTA] Esempio: dato uno schema R(A,B,C) con dipendenze funzionali A->B e B->C, mostrare la chiusura A^+ e discutere se A è superchiave. Soluzione sintetica: A^+={A,B,C}, quindi A determina tutti gli attributi di R ed è superchiave.
\subsection*{Note su assunzioni o integrazioni}
[ASSUNZIONE] L'estrazione OCR può perdere formattazione o simboli della slide originale.
[AGGIUNTA] Le spiegazioni sono estensioni didattiche per rendere espliciti passaggi impliciti nel materiale proiettato.
\paragraph{Riferimento fonte} File: Lezione 01 - Dipendenze funzionali e Normalizzazione di DB Relazionali.pdf. Numero slide: 45.
\section{Dipendenza Funzionale: Esempio --- Slide 46}
\subsection*{Estratto originale dalla slide 46}
\begin{lstlisting}
Dipendenza Funzionale: Esempio
-    La dipendenza funzionale è una proprietà della
    semantica degli attributi.




      FD3: PNUMBER -> {PNAME, PLOCATION}

      FD2: SSN -> ENAME
      FD1: {SSN, PNUMBER} -> HOURS
\end{lstlisting}
\subsection*{Spiegazione estesa}
[AGGIUNTA] Questa spiegazione collega il contenuto testuale della slide al contesto del corso. La slide enfatizza i concetti mostrati nell'estratto e il loro ruolo nella progettazione di basi di dati relazionali, con particolare attenzione a correttezza semantica, qualità dello schema e compromessi progettuali.
[AGGIUNTA] Fonte integrazione: riformulazione didattica del testo presente nella slide 46 del file Lezione 01 - Dipendenze funzionali e Normalizzazione di DB Relazionali.pdf.
\subsection*{Derivazioni}
[AGGIUNTA] Nella slide compaiono espressioni simboliche/relazioni. Derivazione didattica generale:
\begin{align*}
X^+ &= X \\
X^+ &\leftarrow X^+ \cup Y \quad \text{se } (U \to V) \in F \text{ e } U \subseteq X^+ \text{ con } Y=V \\
\text{Stop} &\text{ quando non si aggiungono nuovi attributi in } X^+
\end{align*}
[AGGIUNTA] Questa procedura esplicita il metodo di chiusura usato nelle slide su dipendenze e chiavi.
\begin{itemize}[leftmargin=*]
\item Forma osservata nella slide: \texttt{FD3: PNUMBER -> (PNAME, PLOCATION)}
\item Forma osservata nella slide: \texttt{FD2: SSN -> ENAME}
\item Forma osservata nella slide: \texttt{FD1: (SSN, PNUMBER) -> HOURS}
\end{itemize}
\subsection*{Esempio numerico / esercizio}
[AGGIUNTA] Esempio: dato uno schema R(A,B,C) con dipendenze funzionali A->B e B->C, mostrare la chiusura A^+ e discutere se A è superchiave. Soluzione sintetica: A^+={A,B,C}, quindi A determina tutti gli attributi di R ed è superchiave.
\subsection*{Note su assunzioni o integrazioni}
[ASSUNZIONE] L'estrazione OCR può perdere formattazione o simboli della slide originale.
[AGGIUNTA] Le spiegazioni sono estensioni didattiche per rendere espliciti passaggi impliciti nel materiale proiettato.
\paragraph{Riferimento fonte} File: Lezione 01 - Dipendenze funzionali e Normalizzazione di DB Relazionali.pdf. Numero slide: 46.
\section{Dipendenza funzionale dalla semantica --- Slide 47}
\subsection*{Estratto originale dalla slide 47}
\begin{lstlisting}
Dipendenza funzionale dalla semantica
degli attributi
-    Una dipendenza funzionale è una proprietà dello
    schema di relazione (intenzione) R e non di un
    particolare stato di relazione (estensione r di R).
      Le estensioni valide o stati di relazione legali sono delle
      estensioni di relazione r(R) che soddisfano i vincoli di
      FD.


-    Una FD non può essere inferita automaticamente
    da un'estensione di relazione r, ma deve essere
    definita esplicitamente da chi conosce la
    semantica degli attributi.
\end{lstlisting}
\subsection*{Spiegazione estesa}
[AGGIUNTA] Questa spiegazione collega il contenuto testuale della slide al contesto del corso. La slide enfatizza i concetti mostrati nell'estratto e il loro ruolo nella progettazione di basi di dati relazionali, con particolare attenzione a correttezza semantica, qualità dello schema e compromessi progettuali.
[AGGIUNTA] Fonte integrazione: riformulazione didattica del testo presente nella slide 47 del file Lezione 01 - Dipendenze funzionali e Normalizzazione di DB Relazionali.pdf.
\subsection*{Derivazioni}
[AGGIUNTA] In questa slide non sono state rilevate equazioni esplicite nel testo estratto; non è quindi richiesta una derivazione algebrica puntuale.
\subsection*{Esempio numerico / esercizio}
[AGGIUNTA] Esempio: dato uno schema R(A,B,C) con dipendenze funzionali A->B e B->C, mostrare la chiusura A^+ e discutere se A è superchiave. Soluzione sintetica: A^+={A,B,C}, quindi A determina tutti gli attributi di R ed è superchiave.
\subsection*{Note su assunzioni o integrazioni}
[ASSUNZIONE] L'estrazione OCR può perdere formattazione o simboli della slide originale.
[AGGIUNTA] Le spiegazioni sono estensioni didattiche per rendere espliciti passaggi impliciti nel materiale proiettato.
\paragraph{Riferimento fonte} File: Lezione 01 - Dipendenze funzionali e Normalizzazione di DB Relazionali.pdf. Numero slide: 47.
\section{Dipendenza Funzionale: Esempio --- Slide 48}
\subsection*{Estratto originale dalla slide 48}
\begin{lstlisting}
Dipendenza Funzionale: Esempio




-    Si potrebbe pensare che TEACHER determina
    funzionalmente COURSE e che COURSE determina
    funzionalmente TEXT.
    Ci sono però delle tuple che non seguono queste FD, e
    quindi giungiamo alla conclusione che:
       TEACHER x-> COURSE
       COURSE x-> TEXT
\end{lstlisting}
\subsection*{Spiegazione estesa}
[AGGIUNTA] Questa spiegazione collega il contenuto testuale della slide al contesto del corso. La slide enfatizza i concetti mostrati nell'estratto e il loro ruolo nella progettazione di basi di dati relazionali, con particolare attenzione a correttezza semantica, qualità dello schema e compromessi progettuali.
[AGGIUNTA] Fonte integrazione: riformulazione didattica del testo presente nella slide 48 del file Lezione 01 - Dipendenze funzionali e Normalizzazione di DB Relazionali.pdf.
\subsection*{Derivazioni}
[AGGIUNTA] Nella slide compaiono espressioni simboliche/relazioni. Derivazione didattica generale:
\begin{align*}
X^+ &= X \\
X^+ &\leftarrow X^+ \cup Y \quad \text{se } (U \to V) \in F \text{ e } U \subseteq X^+ \text{ con } Y=V \\
\text{Stop} &\text{ quando non si aggiungono nuovi attributi in } X^+
\end{align*}
[AGGIUNTA] Questa procedura esplicita il metodo di chiusura usato nelle slide su dipendenze e chiavi.
\begin{itemize}[leftmargin=*]
\item Forma osservata nella slide: \texttt{TEACHER x-> COURSE}
\item Forma osservata nella slide: \texttt{COURSE x-> TEXT}
\end{itemize}
\subsection*{Esempio numerico / esercizio}
[AGGIUNTA] Esempio: dato uno schema R(A,B,C) con dipendenze funzionali A->B e B->C, mostrare la chiusura A^+ e discutere se A è superchiave. Soluzione sintetica: A^+={A,B,C}, quindi A determina tutti gli attributi di R ed è superchiave.
\subsection*{Note su assunzioni o integrazioni}
[ASSUNZIONE] L'estrazione OCR può perdere formattazione o simboli della slide originale.
[AGGIUNTA] Le spiegazioni sono estensioni didattiche per rendere espliciti passaggi impliciti nel materiale proiettato.
\paragraph{Riferimento fonte} File: Lezione 01 - Dipendenze funzionali e Normalizzazione di DB Relazionali.pdf. Numero slide: 48.
\section{Regole di inferenza per le FD --- Slide 49}
\subsection*{Estratto originale dalla slide 49}
\begin{lstlisting}
Regole di inferenza per le FD
-    Sia F l'insieme delle dipendenze funzionali di uno
    schema di relazione R. Il progettista dello schema
    specifica le dipendenze funzionali
    semanticamente ovvie F.

-    In tutte le istanze di relazione legali, però,
    valgono numerose altre dipendenze funzionali:
    queste possono essere inferite (o dedotte) da F.

-    L'insieme di tali dipendenze è detto chiusura di F,
    ed è denotato da F+.
\end{lstlisting}
\subsection*{Spiegazione estesa}
[AGGIUNTA] Questa spiegazione collega il contenuto testuale della slide al contesto del corso. La slide enfatizza i concetti mostrati nell'estratto e il loro ruolo nella progettazione di basi di dati relazionali, con particolare attenzione a correttezza semantica, qualità dello schema e compromessi progettuali.
[AGGIUNTA] Fonte integrazione: riformulazione didattica del testo presente nella slide 49 del file Lezione 01 - Dipendenze funzionali e Normalizzazione di DB Relazionali.pdf.
\subsection*{Derivazioni}
[AGGIUNTA] In questa slide non sono state rilevate equazioni esplicite nel testo estratto; non è quindi richiesta una derivazione algebrica puntuale.
\subsection*{Esempio numerico / esercizio}
[AGGIUNTA] Esempio: dato uno schema R(A,B,C) con dipendenze funzionali A->B e B->C, mostrare la chiusura A^+ e discutere se A è superchiave. Soluzione sintetica: A^+={A,B,C}, quindi A determina tutti gli attributi di R ed è superchiave.
\subsection*{Note su assunzioni o integrazioni}
[ASSUNZIONE] L'estrazione OCR può perdere formattazione o simboli della slide originale.
[AGGIUNTA] Le spiegazioni sono estensioni didattiche per rendere espliciti passaggi impliciti nel materiale proiettato.
\paragraph{Riferimento fonte} File: Lezione 01 - Dipendenze funzionali e Normalizzazione di DB Relazionali.pdf. Numero slide: 49.
\section{Regole di inferenza: Esempio --- Slide 50}
\subsection*{Estratto originale dalla slide 50}
\begin{lstlisting}
Regole di inferenza: Esempio



-    F={
      SSN -> { ENAME, BDATE, ADDRESS, DNUMBER},
      DNUMBER -> {DNAME,DMGRSSN}
      }

-    Possiamo inferire:
      SSN -> {DNAME,DMGRSSN}
      SSN -> SSN

      DNUMBER -> DNAME
      ...
\end{lstlisting}
\subsection*{Spiegazione estesa}
[AGGIUNTA] Questa spiegazione collega il contenuto testuale della slide al contesto del corso. La slide enfatizza i concetti mostrati nell'estratto e il loro ruolo nella progettazione di basi di dati relazionali, con particolare attenzione a correttezza semantica, qualità dello schema e compromessi progettuali.
[AGGIUNTA] Fonte integrazione: riformulazione didattica del testo presente nella slide 50 del file Lezione 01 - Dipendenze funzionali e Normalizzazione di DB Relazionali.pdf.
\subsection*{Derivazioni}
[AGGIUNTA] Nella slide compaiono espressioni simboliche/relazioni. Derivazione didattica generale:
\begin{align*}
X^+ &= X \\
X^+ &\leftarrow X^+ \cup Y \quad \text{se } (U \to V) \in F \text{ e } U \subseteq X^+ \text{ con } Y=V \\
\text{Stop} &\text{ quando non si aggiungono nuovi attributi in } X^+
\end{align*}
[AGGIUNTA] Questa procedura esplicita il metodo di chiusura usato nelle slide su dipendenze e chiavi.
\begin{itemize}[leftmargin=*]
\item Forma osservata nella slide: \texttt{-    F=(}
\item Forma osservata nella slide: \texttt{SSN -> ( ENAME, BDATE, ADDRESS, DNUMBER),}
\item Forma osservata nella slide: \texttt{DNUMBER -> (DNAME,DMGRSSN)}
\end{itemize}
\subsection*{Esempio numerico / esercizio}
[AGGIUNTA] Esempio: dato uno schema R(A,B,C) con dipendenze funzionali A->B e B->C, mostrare la chiusura A^+ e discutere se A è superchiave. Soluzione sintetica: A^+={A,B,C}, quindi A determina tutti gli attributi di R ed è superchiave.
\subsection*{Note su assunzioni o integrazioni}
[ASSUNZIONE] L'estrazione OCR può perdere formattazione o simboli della slide originale.
[AGGIUNTA] Le spiegazioni sono estensioni didattiche per rendere espliciti passaggi impliciti nel materiale proiettato.
\paragraph{Riferimento fonte} File: Lezione 01 - Dipendenze funzionali e Normalizzazione di DB Relazionali.pdf. Numero slide: 50.
\section{Regole di inferenza per le FD (2) --- Slide 51}
\subsection*{Estratto originale dalla slide 51}
\begin{lstlisting}
Regole di inferenza per le FD (2)
-    Una FD X->Y è inferita da un insieme di
    dipendenze F su R se X->Y vale in ogni stato di
    relazione r che è estensione legale di R.

-    F  X->Y denota che la FD X->Y è inferita da F.

-    Notazioni:
      {X,Y} -> Z è abbreviato in XY->Z
      {X,Y,Z} -> {U,V} è abbreviato in XYZ -> UV
\end{lstlisting}
\subsection*{Spiegazione estesa}
[AGGIUNTA] Questa spiegazione collega il contenuto testuale della slide al contesto del corso. La slide enfatizza i concetti mostrati nell'estratto e il loro ruolo nella progettazione di basi di dati relazionali, con particolare attenzione a correttezza semantica, qualità dello schema e compromessi progettuali.
[AGGIUNTA] Fonte integrazione: riformulazione didattica del testo presente nella slide 51 del file Lezione 01 - Dipendenze funzionali e Normalizzazione di DB Relazionali.pdf.
\subsection*{Derivazioni}
[AGGIUNTA] Nella slide compaiono espressioni simboliche/relazioni. Derivazione didattica generale:
\begin{align*}
X^+ &= X \\
X^+ &\leftarrow X^+ \cup Y \quad \text{se } (U \to V) \in F \text{ e } U \subseteq X^+ \text{ con } Y=V \\
\text{Stop} &\text{ quando non si aggiungono nuovi attributi in } X^+
\end{align*}
[AGGIUNTA] Questa procedura esplicita il metodo di chiusura usato nelle slide su dipendenze e chiavi.
\begin{itemize}[leftmargin=*]
\item Forma osservata nella slide: \texttt{-    Una FD X->Y è inferita da un insieme di}
\item Forma osservata nella slide: \texttt{dipendenze F su R se X->Y vale in ogni stato di}
\item Forma osservata nella slide: \texttt{-    F  X->Y denota che la FD X->Y è inferita da F.}
\end{itemize}
\subsection*{Esempio numerico / esercizio}
[AGGIUNTA] Esempio: dato uno schema R(A,B,C) con dipendenze funzionali A->B e B->C, mostrare la chiusura A^+ e discutere se A è superchiave. Soluzione sintetica: A^+={A,B,C}, quindi A determina tutti gli attributi di R ed è superchiave.
\subsection*{Note su assunzioni o integrazioni}
[ASSUNZIONE] L'estrazione OCR può perdere formattazione o simboli della slide originale.
[AGGIUNTA] Le spiegazioni sono estensioni didattiche per rendere espliciti passaggi impliciti nel materiale proiettato.
\paragraph{Riferimento fonte} File: Lezione 01 - Dipendenze funzionali e Normalizzazione di DB Relazionali.pdf. Numero slide: 51.
\section{Regole di inferenza fondamentali --- Slide 52}
\subsection*{Estratto originale dalla slide 52}
\begin{lstlisting}
Regole di inferenza fondamentali
-    Regole di inferenza per le dipendenze funzionali:
-    (IR1) (Reflexive rule)
      Se X Ê Y allora X -> Y


-    (IR2) (Augmentation rule)
      X -> Y  XZ -> YZ


-    (IR3) (Transitive rule)
      {X -> Y , Y -> Z}  X -> Z
\end{lstlisting}
\subsection*{Spiegazione estesa}
[AGGIUNTA] Questa spiegazione collega il contenuto testuale della slide al contesto del corso. La slide enfatizza i concetti mostrati nell'estratto e il loro ruolo nella progettazione di basi di dati relazionali, con particolare attenzione a correttezza semantica, qualità dello schema e compromessi progettuali.
[AGGIUNTA] Fonte integrazione: riformulazione didattica del testo presente nella slide 52 del file Lezione 01 - Dipendenze funzionali e Normalizzazione di DB Relazionali.pdf.
\subsection*{Derivazioni}
[AGGIUNTA] Nella slide compaiono espressioni simboliche/relazioni. Derivazione didattica generale:
\begin{align*}
X^+ &= X \\
X^+ &\leftarrow X^+ \cup Y \quad \text{se } (U \to V) \in F \text{ e } U \subseteq X^+ \text{ con } Y=V \\
\text{Stop} &\text{ quando non si aggiungono nuovi attributi in } X^+
\end{align*}
[AGGIUNTA] Questa procedura esplicita il metodo di chiusura usato nelle slide su dipendenze e chiavi.
\begin{itemize}[leftmargin=*]
\item Forma osservata nella slide: \texttt{Se X Ê Y allora X -> Y}
\item Forma osservata nella slide: \texttt{X -> Y  XZ -> YZ}
\item Forma osservata nella slide: \texttt{(X -> Y , Y -> Z)  X -> Z}
\end{itemize}
\subsection*{Esempio numerico / esercizio}
[AGGIUNTA] Esempio: dato uno schema R(A,B,C) con dipendenze funzionali A->B e B->C, mostrare la chiusura A^+ e discutere se A è superchiave. Soluzione sintetica: A^+={A,B,C}, quindi A determina tutti gli attributi di R ed è superchiave.
\subsection*{Note su assunzioni o integrazioni}
[ASSUNZIONE] L'estrazione OCR può perdere formattazione o simboli della slide originale.
[AGGIUNTA] Le spiegazioni sono estensioni didattiche per rendere espliciti passaggi impliciti nel materiale proiettato.
\paragraph{Riferimento fonte} File: Lezione 01 - Dipendenze funzionali e Normalizzazione di DB Relazionali.pdf. Numero slide: 52.
\section{Regole di inferenza fondamentali (2) --- Slide 53}
\subsection*{Estratto originale dalla slide 53}
\begin{lstlisting}
Regole di inferenza fondamentali (2)
-    (IR4) (Decomposition or Projective rule)
      {X -> YZ}  X -> Z


-    (IR5) (Union (or additive) rule)
      {X -> Y , X -> Z}  X -> YZ


-    (IR6) (Pseudotransitive rule)
      {X -> Y , WY -> Z}  WX -> Z
\end{lstlisting}
\subsection*{Spiegazione estesa}
[AGGIUNTA] Questa spiegazione collega il contenuto testuale della slide al contesto del corso. La slide enfatizza i concetti mostrati nell'estratto e il loro ruolo nella progettazione di basi di dati relazionali, con particolare attenzione a correttezza semantica, qualità dello schema e compromessi progettuali.
[AGGIUNTA] Fonte integrazione: riformulazione didattica del testo presente nella slide 53 del file Lezione 01 - Dipendenze funzionali e Normalizzazione di DB Relazionali.pdf.
\subsection*{Derivazioni}
[AGGIUNTA] Nella slide compaiono espressioni simboliche/relazioni. Derivazione didattica generale:
\begin{align*}
X^+ &= X \\
X^+ &\leftarrow X^+ \cup Y \quad \text{se } (U \to V) \in F \text{ e } U \subseteq X^+ \text{ con } Y=V \\
\text{Stop} &\text{ quando non si aggiungono nuovi attributi in } X^+
\end{align*}
[AGGIUNTA] Questa procedura esplicita il metodo di chiusura usato nelle slide su dipendenze e chiavi.
\begin{itemize}[leftmargin=*]
\item Forma osservata nella slide: \texttt{(X -> YZ)  X -> Z}
\item Forma osservata nella slide: \texttt{(X -> Y , X -> Z)  X -> YZ}
\item Forma osservata nella slide: \texttt{(X -> Y , WY -> Z)  WX -> Z}
\end{itemize}
\subsection*{Esempio numerico / esercizio}
[AGGIUNTA] Esempio: dato uno schema R(A,B,C) con dipendenze funzionali A->B e B->C, mostrare la chiusura A^+ e discutere se A è superchiave. Soluzione sintetica: A^+={A,B,C}, quindi A determina tutti gli attributi di R ed è superchiave.
\subsection*{Note su assunzioni o integrazioni}
[ASSUNZIONE] L'estrazione OCR può perdere formattazione o simboli della slide originale.
[AGGIUNTA] Le spiegazioni sono estensioni didattiche per rendere espliciti passaggi impliciti nel materiale proiettato.
\paragraph{Riferimento fonte} File: Lezione 01 - Dipendenze funzionali e Normalizzazione di DB Relazionali.pdf. Numero slide: 53.
\section{Dimostrazione di IR1 --- Slide 54}
\subsection*{Estratto originale dalla slide 54}
\begin{lstlisting}
Dimostrazione di IR1
              Se X ÊY allora X -> Y

-  La tesi X -> Y significa che
    t1, t2, t1[X] = t2[X]   t1[Y] = t2[Y]
-  Sia X Ê Y e supponiamo che esistano tuple t 1, t2

  in qualche istanza r di R tale che t1[X] = t2[X].
  Ma allora t1[Y] = t2[Y] poiché X Ê Y.
  Quindi X -> Y.
\end{lstlisting}
\subsection*{Spiegazione estesa}
[AGGIUNTA] Questa spiegazione collega il contenuto testuale della slide al contesto del corso. La slide enfatizza i concetti mostrati nell'estratto e il loro ruolo nella progettazione di basi di dati relazionali, con particolare attenzione a correttezza semantica, qualità dello schema e compromessi progettuali.
[AGGIUNTA] Fonte integrazione: riformulazione didattica del testo presente nella slide 54 del file Lezione 01 - Dipendenze funzionali e Normalizzazione di DB Relazionali.pdf.
\subsection*{Derivazioni}
[AGGIUNTA] Nella slide compaiono espressioni simboliche/relazioni. Derivazione didattica generale:
\begin{align*}
X^+ &= X \\
X^+ &\leftarrow X^+ \cup Y \quad \text{se } (U \to V) \in F \text{ e } U \subseteq X^+ \text{ con } Y=V \\
\text{Stop} &\text{ quando non si aggiungono nuovi attributi in } X^+
\end{align*}
[AGGIUNTA] Questa procedura esplicita il metodo di chiusura usato nelle slide su dipendenze e chiavi.
\begin{itemize}[leftmargin=*]
\item Forma osservata nella slide: \texttt{Se X ÊY allora X -> Y}
\item Forma osservata nella slide: \texttt{-  La tesi X -> Y significa che}
\item Forma osservata nella slide: \texttt{t1, t2, t1[X] = t2[X]   t1[Y] = t2[Y]}
\end{itemize}
\subsection*{Esempio numerico / esercizio}
[AGGIUNTA] Esempio: dato uno schema R(A,B,C) con dipendenze funzionali A->B e B->C, mostrare la chiusura A^+ e discutere se A è superchiave. Soluzione sintetica: A^+={A,B,C}, quindi A determina tutti gli attributi di R ed è superchiave.
\subsection*{Note su assunzioni o integrazioni}
[ASSUNZIONE] L'estrazione OCR può perdere formattazione o simboli della slide originale.
[AGGIUNTA] Le spiegazioni sono estensioni didattiche per rendere espliciti passaggi impliciti nel materiale proiettato.
\paragraph{Riferimento fonte} File: Lezione 01 - Dipendenze funzionali e Normalizzazione di DB Relazionali.pdf. Numero slide: 54.
\section{Dimostrazione di IR2 --- Slide 55}
\subsection*{Estratto originale dalla slide 55}
\begin{lstlisting}
Dimostrazione di IR2
                  {X -> Y}   XZ -> YZ
-  Per assurdo assumiamo che X -> Y vale ma

  XZ -> YZ non vale.
-  Allora devono esistere t 1 e t2 in r tali che

  (1) t1[X] = t2[X], (2) t1[Y] = t2[Y],
  (3) t1[XZ] = t2[XZ] e (4) t1[YZ]   t2[YZ].
  Ma ciò è falso perché:
  t1[X] = t2[X] e t1[XZ] = t2[XZ]   t1[Z] = t2[Z]
  Inoltre
  t1[Y] = t2[Y] e t1[Z] = t2[Z]   t1[YZ] = t2[YZ] che
  contraddice l'assunto.
\end{lstlisting}
\subsection*{Spiegazione estesa}
[AGGIUNTA] Questa spiegazione collega il contenuto testuale della slide al contesto del corso. La slide enfatizza i concetti mostrati nell'estratto e il loro ruolo nella progettazione di basi di dati relazionali, con particolare attenzione a correttezza semantica, qualità dello schema e compromessi progettuali.
[AGGIUNTA] Fonte integrazione: riformulazione didattica del testo presente nella slide 55 del file Lezione 01 - Dipendenze funzionali e Normalizzazione di DB Relazionali.pdf.
\subsection*{Derivazioni}
[AGGIUNTA] Nella slide compaiono espressioni simboliche/relazioni. Derivazione didattica generale:
\begin{align*}
X^+ &= X \\
X^+ &\leftarrow X^+ \cup Y \quad \text{se } (U \to V) \in F \text{ e } U \subseteq X^+ \text{ con } Y=V \\
\text{Stop} &\text{ quando non si aggiungono nuovi attributi in } X^+
\end{align*}
[AGGIUNTA] Questa procedura esplicita il metodo di chiusura usato nelle slide su dipendenze e chiavi.
\begin{itemize}[leftmargin=*]
\item Forma osservata nella slide: \texttt{(X -> Y)   XZ -> YZ}
\item Forma osservata nella slide: \texttt{-  Per assurdo assumiamo che X -> Y vale ma}
\item Forma osservata nella slide: \texttt{XZ -> YZ non vale.}
\end{itemize}
\subsection*{Esempio numerico / esercizio}
[AGGIUNTA] Esempio: dato uno schema R(A,B,C) con dipendenze funzionali A->B e B->C, mostrare la chiusura A^+ e discutere se A è superchiave. Soluzione sintetica: A^+={A,B,C}, quindi A determina tutti gli attributi di R ed è superchiave.
\subsection*{Note su assunzioni o integrazioni}
[ASSUNZIONE] L'estrazione OCR può perdere formattazione o simboli della slide originale.
[AGGIUNTA] Le spiegazioni sono estensioni didattiche per rendere espliciti passaggi impliciti nel materiale proiettato.
\paragraph{Riferimento fonte} File: Lezione 01 - Dipendenze funzionali e Normalizzazione di DB Relazionali.pdf. Numero slide: 55.
\section{Dimostrazione di IR3 --- Slide 56}
\subsection*{Estratto originale dalla slide 56}
\begin{lstlisting}
Dimostrazione di IR3
              {X -> Y , Y -> Z}   X -> Z

-  Assumiamo che X -> Y e Y -> Z valgono in una
  relazione r.
-  Allora   t1, t2 in r: t1[X] = t2[X]   t1[Y] = t2[Y] e

  quindi poiché Y -> Z, t1[Y] = t2[Y]   t1[Z] = t2[Z];
  quindi X -> Z vale in r.
\end{lstlisting}
\subsection*{Spiegazione estesa}
[AGGIUNTA] Questa spiegazione collega il contenuto testuale della slide al contesto del corso. La slide enfatizza i concetti mostrati nell'estratto e il loro ruolo nella progettazione di basi di dati relazionali, con particolare attenzione a correttezza semantica, qualità dello schema e compromessi progettuali.
[AGGIUNTA] Fonte integrazione: riformulazione didattica del testo presente nella slide 56 del file Lezione 01 - Dipendenze funzionali e Normalizzazione di DB Relazionali.pdf.
\subsection*{Derivazioni}
[AGGIUNTA] Nella slide compaiono espressioni simboliche/relazioni. Derivazione didattica generale:
\begin{align*}
X^+ &= X \\
X^+ &\leftarrow X^+ \cup Y \quad \text{se } (U \to V) \in F \text{ e } U \subseteq X^+ \text{ con } Y=V \\
\text{Stop} &\text{ quando non si aggiungono nuovi attributi in } X^+
\end{align*}
[AGGIUNTA] Questa procedura esplicita il metodo di chiusura usato nelle slide su dipendenze e chiavi.
\begin{itemize}[leftmargin=*]
\item Forma osservata nella slide: \texttt{(X -> Y , Y -> Z)   X -> Z}
\item Forma osservata nella slide: \texttt{-  Assumiamo che X -> Y e Y -> Z valgono in una}
\item Forma osservata nella slide: \texttt{-  Allora   t1, t2 in r: t1[X] = t2[X]   t1[Y] = t2[Y] e}
\end{itemize}
\subsection*{Esempio numerico / esercizio}
[AGGIUNTA] Esempio: dato uno schema R(A,B,C) con dipendenze funzionali A->B e B->C, mostrare la chiusura A^+ e discutere se A è superchiave. Soluzione sintetica: A^+={A,B,C}, quindi A determina tutti gli attributi di R ed è superchiave.
\subsection*{Note su assunzioni o integrazioni}
[ASSUNZIONE] L'estrazione OCR può perdere formattazione o simboli della slide originale.
[AGGIUNTA] Le spiegazioni sono estensioni didattiche per rendere espliciti passaggi impliciti nel materiale proiettato.
\paragraph{Riferimento fonte} File: Lezione 01 - Dipendenze funzionali e Normalizzazione di DB Relazionali.pdf. Numero slide: 56.
\section{Dimostrazione di IR4 --- Slide 57}
\subsection*{Estratto originale dalla slide 57}
\begin{lstlisting}
Dimostrazione di IR4
 (IR4) / (IR6) possono essere provate usando
 (IR1) / (IR3)

 (IR4):    {X -> YZ}  X -> Y
 1.   X -> YZ   (per ipotesi)
 2.   YZ -> Y   da IR1 e dato che YZ   ÊY
 3.   X->Y      da IR3, 1 e 2.
\end{lstlisting}
\subsection*{Spiegazione estesa}
[AGGIUNTA] Questa spiegazione collega il contenuto testuale della slide al contesto del corso. La slide enfatizza i concetti mostrati nell'estratto e il loro ruolo nella progettazione di basi di dati relazionali, con particolare attenzione a correttezza semantica, qualità dello schema e compromessi progettuali.
[AGGIUNTA] Fonte integrazione: riformulazione didattica del testo presente nella slide 57 del file Lezione 01 - Dipendenze funzionali e Normalizzazione di DB Relazionali.pdf.
\subsection*{Derivazioni}
[AGGIUNTA] Nella slide compaiono espressioni simboliche/relazioni. Derivazione didattica generale:
\begin{align*}
X^+ &= X \\
X^+ &\leftarrow X^+ \cup Y \quad \text{se } (U \to V) \in F \text{ e } U \subseteq X^+ \text{ con } Y=V \\
\text{Stop} &\text{ quando non si aggiungono nuovi attributi in } X^+
\end{align*}
[AGGIUNTA] Questa procedura esplicita il metodo di chiusura usato nelle slide su dipendenze e chiavi.
\begin{itemize}[leftmargin=*]
\item Forma osservata nella slide: \texttt{(IR4):    (X -> YZ)  X -> Y}
\item Forma osservata nella slide: \texttt{1.   X -> YZ   (per ipotesi)}
\item Forma osservata nella slide: \texttt{2.   YZ -> Y   da IR1 e dato che YZ   ÊY}
\end{itemize}
\subsection*{Esempio numerico / esercizio}
[AGGIUNTA] Esempio: dato uno schema R(A,B,C) con dipendenze funzionali A->B e B->C, mostrare la chiusura A^+ e discutere se A è superchiave. Soluzione sintetica: A^+={A,B,C}, quindi A determina tutti gli attributi di R ed è superchiave.
\subsection*{Note su assunzioni o integrazioni}
[ASSUNZIONE] L'estrazione OCR può perdere formattazione o simboli della slide originale.
[AGGIUNTA] Le spiegazioni sono estensioni didattiche per rendere espliciti passaggi impliciti nel materiale proiettato.
\paragraph{Riferimento fonte} File: Lezione 01 - Dipendenze funzionali e Normalizzazione di DB Relazionali.pdf. Numero slide: 57.
\section{Dimostrazione di IR5 --- Slide 58}
\subsection*{Estratto originale dalla slide 58}
\begin{lstlisting}
Dimostrazione di IR5
             { X -> Y , X -> Z}  X -> YZ

 1.   X->Y          (per ipotesi)
 2.   X->Z          (per ipotesi)
 3.   X -> XY
      Poiché X = XX allora X -> XX, e da 1. e usando IR2
      XX -> XY e IR3 X -> XY
 4.   XY -> YZ      da 2. e usando IR2 (aumentando con Y)
 5.   X -> YZ       da IR3, 3. e 4.
\end{lstlisting}
\subsection*{Spiegazione estesa}
[AGGIUNTA] Questa spiegazione collega il contenuto testuale della slide al contesto del corso. La slide enfatizza i concetti mostrati nell'estratto e il loro ruolo nella progettazione di basi di dati relazionali, con particolare attenzione a correttezza semantica, qualità dello schema e compromessi progettuali.
[AGGIUNTA] Fonte integrazione: riformulazione didattica del testo presente nella slide 58 del file Lezione 01 - Dipendenze funzionali e Normalizzazione di DB Relazionali.pdf.
\subsection*{Derivazioni}
[AGGIUNTA] Nella slide compaiono espressioni simboliche/relazioni. Derivazione didattica generale:
\begin{align*}
X^+ &= X \\
X^+ &\leftarrow X^+ \cup Y \quad \text{se } (U \to V) \in F \text{ e } U \subseteq X^+ \text{ con } Y=V \\
\text{Stop} &\text{ quando non si aggiungono nuovi attributi in } X^+
\end{align*}
[AGGIUNTA] Questa procedura esplicita il metodo di chiusura usato nelle slide su dipendenze e chiavi.
\begin{itemize}[leftmargin=*]
\item Forma osservata nella slide: \texttt{( X -> Y , X -> Z)  X -> YZ}
\item Forma osservata nella slide: \texttt{1.   X->Y          (per ipotesi)}
\item Forma osservata nella slide: \texttt{2.   X->Z          (per ipotesi)}
\end{itemize}
\subsection*{Esempio numerico / esercizio}
[AGGIUNTA] Esempio: dato uno schema R(A,B,C) con dipendenze funzionali A->B e B->C, mostrare la chiusura A^+ e discutere se A è superchiave. Soluzione sintetica: A^+={A,B,C}, quindi A determina tutti gli attributi di R ed è superchiave.
\subsection*{Note su assunzioni o integrazioni}
[ASSUNZIONE] L'estrazione OCR può perdere formattazione o simboli della slide originale.
[AGGIUNTA] Le spiegazioni sono estensioni didattiche per rendere espliciti passaggi impliciti nel materiale proiettato.
\paragraph{Riferimento fonte} File: Lezione 01 - Dipendenze funzionali e Normalizzazione di DB Relazionali.pdf. Numero slide: 58.
\section{Dimostrazione di IR6 --- Slide 59}
\subsection*{Estratto originale dalla slide 59}
\begin{lstlisting}
Dimostrazione di IR6
          { X -> Y , WY -> Z}  WX -> Z

 1.   X-> Y      (per ipotesi)
 2.   WY -> Z    (per ipotesi)
 3.   WX -> WY   da 1. e IR2
 4.   WX -> Z    da 3. e 2. con IR3.
\end{lstlisting}
\subsection*{Spiegazione estesa}
[AGGIUNTA] Questa spiegazione collega il contenuto testuale della slide al contesto del corso. La slide enfatizza i concetti mostrati nell'estratto e il loro ruolo nella progettazione di basi di dati relazionali, con particolare attenzione a correttezza semantica, qualità dello schema e compromessi progettuali.
[AGGIUNTA] Fonte integrazione: riformulazione didattica del testo presente nella slide 59 del file Lezione 01 - Dipendenze funzionali e Normalizzazione di DB Relazionali.pdf.
\subsection*{Derivazioni}
[AGGIUNTA] Nella slide compaiono espressioni simboliche/relazioni. Derivazione didattica generale:
\begin{align*}
X^+ &= X \\
X^+ &\leftarrow X^+ \cup Y \quad \text{se } (U \to V) \in F \text{ e } U \subseteq X^+ \text{ con } Y=V \\
\text{Stop} &\text{ quando non si aggiungono nuovi attributi in } X^+
\end{align*}
[AGGIUNTA] Questa procedura esplicita il metodo di chiusura usato nelle slide su dipendenze e chiavi.
\begin{itemize}[leftmargin=*]
\item Forma osservata nella slide: \texttt{( X -> Y , WY -> Z)  WX -> Z}
\item Forma osservata nella slide: \texttt{1.   X-> Y      (per ipotesi)}
\item Forma osservata nella slide: \texttt{2.   WY -> Z    (per ipotesi)}
\end{itemize}
\subsection*{Esempio numerico / esercizio}
[AGGIUNTA] Esempio: dato uno schema R(A,B,C) con dipendenze funzionali A->B e B->C, mostrare la chiusura A^+ e discutere se A è superchiave. Soluzione sintetica: A^+={A,B,C}, quindi A determina tutti gli attributi di R ed è superchiave.
\subsection*{Note su assunzioni o integrazioni}
[ASSUNZIONE] L'estrazione OCR può perdere formattazione o simboli della slide originale.
[AGGIUNTA] Le spiegazioni sono estensioni didattiche per rendere espliciti passaggi impliciti nel materiale proiettato.
\paragraph{Riferimento fonte} File: Lezione 01 - Dipendenze funzionali e Normalizzazione di DB Relazionali.pdf. Numero slide: 59.
\section{Regole di inferenza di Armstrong --- Slide 60}
\subsection*{Estratto originale dalla slide 60}
\begin{lstlisting}
Regole di inferenza di Armstrong
-    È stato provato che (IR1) / (IR3) sono regole di
    inferenza corrette e complete (Armstrong 1974):
      Corretto: Dato un insieme F di dipendenze funzionali su
      R ogni dipendenza inferita da F usando IR1 / IR3 è
      vera in ogni stato che soddisfa le dipendenze in F.

      Completo: la chiusura di F (F+, l'insieme di tutte le
      possibili dipendenze inferibili da F) può essere
      determinata usando solo IR1 / IR3 (regole di inferenza
      di Armstrong).
\end{lstlisting}
\subsection*{Spiegazione estesa}
[AGGIUNTA] Questa spiegazione collega il contenuto testuale della slide al contesto del corso. La slide enfatizza i concetti mostrati nell'estratto e il loro ruolo nella progettazione di basi di dati relazionali, con particolare attenzione a correttezza semantica, qualità dello schema e compromessi progettuali.
[AGGIUNTA] Fonte integrazione: riformulazione didattica del testo presente nella slide 60 del file Lezione 01 - Dipendenze funzionali e Normalizzazione di DB Relazionali.pdf.
\subsection*{Derivazioni}
[AGGIUNTA] In questa slide non sono state rilevate equazioni esplicite nel testo estratto; non è quindi richiesta una derivazione algebrica puntuale.
\subsection*{Esempio numerico / esercizio}
[AGGIUNTA] Esempio: dato uno schema R(A,B,C) con dipendenze funzionali A->B e B->C, mostrare la chiusura A^+ e discutere se A è superchiave. Soluzione sintetica: A^+={A,B,C}, quindi A determina tutti gli attributi di R ed è superchiave.
\subsection*{Note su assunzioni o integrazioni}
[ASSUNZIONE] L'estrazione OCR può perdere formattazione o simboli della slide originale.
[AGGIUNTA] Le spiegazioni sono estensioni didattiche per rendere espliciti passaggi impliciti nel materiale proiettato.
\paragraph{Riferimento fonte} File: Lezione 01 - Dipendenze funzionali e Normalizzazione di DB Relazionali.pdf. Numero slide: 60.
\section{Regole di inferenza nella progettazione dei db --- Slide 61}
\subsection*{Estratto originale dalla slide 61}
\begin{lstlisting}
Regole di inferenza nella progettazione dei db

-  Tipicamente, il progettista del db prima specifica
  le dipendenze funzionali F a partire dalla
  semantica degli attributi e poi usa le regole di
  inferenza di Armstrong per inferire dipendenze
  funzionali aggiuntive.
-  Metodo:
    1. determinare ogni insieme di attributi X che appare
       come parte sinistra di qualche dipendenza funzionale
       in F.
    2. usare le regole di Armstrong per determinare
       l'insieme di tutti gli attributi dipendenti da X.
\end{lstlisting}
\subsection*{Spiegazione estesa}
[AGGIUNTA] Questa spiegazione collega il contenuto testuale della slide al contesto del corso. La slide enfatizza i concetti mostrati nell'estratto e il loro ruolo nella progettazione di basi di dati relazionali, con particolare attenzione a correttezza semantica, qualità dello schema e compromessi progettuali.
[AGGIUNTA] Fonte integrazione: riformulazione didattica del testo presente nella slide 61 del file Lezione 01 - Dipendenze funzionali e Normalizzazione di DB Relazionali.pdf.
\subsection*{Derivazioni}
[AGGIUNTA] In questa slide non sono state rilevate equazioni esplicite nel testo estratto; non è quindi richiesta una derivazione algebrica puntuale.
\subsection*{Esempio numerico / esercizio}
[AGGIUNTA] Esempio: dato uno schema R(A,B,C) con dipendenze funzionali A->B e B->C, mostrare la chiusura A^+ e discutere se A è superchiave. Soluzione sintetica: A^+={A,B,C}, quindi A determina tutti gli attributi di R ed è superchiave.
\subsection*{Note su assunzioni o integrazioni}
[ASSUNZIONE] L'estrazione OCR può perdere formattazione o simboli della slide originale.
[AGGIUNTA] Le spiegazioni sono estensioni didattiche per rendere espliciti passaggi impliciti nel materiale proiettato.
\paragraph{Riferimento fonte} File: Lezione 01 - Dipendenze funzionali e Normalizzazione di DB Relazionali.pdf. Numero slide: 61.
\section{Algoritmo per il calcolo di X+ --- Slide 62}
\subsection*{Estratto originale dalla slide 62}
\begin{lstlisting}
Algoritmo per il calcolo di X+
X+:= X;
repeat
  oldX+ := X+;
  for each functional dependency Y -> Z in F do
       if X+ Ê Y
       then X+ := X+   Z;
until (oldX+ = X+);
\end{lstlisting}
\subsection*{Spiegazione estesa}
[AGGIUNTA] Questa spiegazione collega il contenuto testuale della slide al contesto del corso. La slide enfatizza i concetti mostrati nell'estratto e il loro ruolo nella progettazione di basi di dati relazionali, con particolare attenzione a correttezza semantica, qualità dello schema e compromessi progettuali.
[AGGIUNTA] Fonte integrazione: riformulazione didattica del testo presente nella slide 62 del file Lezione 01 - Dipendenze funzionali e Normalizzazione di DB Relazionali.pdf.
\subsection*{Derivazioni}
[AGGIUNTA] Nella slide compaiono espressioni simboliche/relazioni. Derivazione didattica generale:
\begin{align*}
X^+ &= X \\
X^+ &\leftarrow X^+ \cup Y \quad \text{se } (U \to V) \in F \text{ e } U \subseteq X^+ \text{ con } Y=V \\
\text{Stop} &\text{ quando non si aggiungono nuovi attributi in } X^+
\end{align*}
[AGGIUNTA] Questa procedura esplicita il metodo di chiusura usato nelle slide su dipendenze e chiavi.
\begin{itemize}[leftmargin=*]
\item Forma osservata nella slide: \texttt{X+:= X;}
\item Forma osservata nella slide: \texttt{oldX+ := X+;}
\item Forma osservata nella slide: \texttt{for each functional dependency Y -> Z in F do}
\end{itemize}
\subsection*{Esempio numerico / esercizio}
[AGGIUNTA] Esempio: dato uno schema R(A,B,C) con dipendenze funzionali A->B e B->C, mostrare la chiusura A^+ e discutere se A è superchiave. Soluzione sintetica: A^+={A,B,C}, quindi A determina tutti gli attributi di R ed è superchiave.
\subsection*{Note su assunzioni o integrazioni}
[ASSUNZIONE] L'estrazione OCR può perdere formattazione o simboli della slide originale.
[AGGIUNTA] Le spiegazioni sono estensioni didattiche per rendere espliciti passaggi impliciti nel materiale proiettato.
\paragraph{Riferimento fonte} File: Lezione 01 - Dipendenze funzionali e Normalizzazione di DB Relazionali.pdf. Numero slide: 62.
\section{Algoritmo per il calcolo di X+: Esempio --- Slide 63}
\subsection*{Estratto originale dalla slide 63}
\begin{lstlisting}
Algoritmo per il calcolo di X+: Esempio




-  Dalla semantica degli attributi si ha:
-  F={ {SSN, PNUMBER} -> HOURS,

       SSN -> ENAME,
       PNUMBER -> { PNAME, PLOCATION}
     }
-  Usando l'algoritmo calcoliamo:
      {SSN}+ = {SSN, ENAME}

      {SSN, PNUMBER}+ = {SSN, PNUMBER, ENAME, PNAME,
      PLOCATIONS, HOURS}
      {PNUMBER}+= {PNUMBER, PNAME, PLOCATION}
\end{lstlisting}
\subsection*{Spiegazione estesa}
[AGGIUNTA] Questa spiegazione collega il contenuto testuale della slide al contesto del corso. La slide enfatizza i concetti mostrati nell'estratto e il loro ruolo nella progettazione di basi di dati relazionali, con particolare attenzione a correttezza semantica, qualità dello schema e compromessi progettuali.
[AGGIUNTA] Fonte integrazione: riformulazione didattica del testo presente nella slide 63 del file Lezione 01 - Dipendenze funzionali e Normalizzazione di DB Relazionali.pdf.
\subsection*{Derivazioni}
[AGGIUNTA] Nella slide compaiono espressioni simboliche/relazioni. Derivazione didattica generale:
\begin{align*}
X^+ &= X \\
X^+ &\leftarrow X^+ \cup Y \quad \text{se } (U \to V) \in F \text{ e } U \subseteq X^+ \text{ con } Y=V \\
\text{Stop} &\text{ quando non si aggiungono nuovi attributi in } X^+
\end{align*}
[AGGIUNTA] Questa procedura esplicita il metodo di chiusura usato nelle slide su dipendenze e chiavi.
\begin{itemize}[leftmargin=*]
\item Forma osservata nella slide: \texttt{-  F=( (SSN, PNUMBER) -> HOURS,}
\item Forma osservata nella slide: \texttt{SSN -> ENAME,}
\item Forma osservata nella slide: \texttt{PNUMBER -> ( PNAME, PLOCATION)}
\end{itemize}
\subsection*{Esempio numerico / esercizio}
[AGGIUNTA] Esempio: dato uno schema R(A,B,C) con dipendenze funzionali A->B e B->C, mostrare la chiusura A^+ e discutere se A è superchiave. Soluzione sintetica: A^+={A,B,C}, quindi A determina tutti gli attributi di R ed è superchiave.
\subsection*{Note su assunzioni o integrazioni}
[ASSUNZIONE] L'estrazione OCR può perdere formattazione o simboli della slide originale.
[AGGIUNTA] Le spiegazioni sono estensioni didattiche per rendere espliciti passaggi impliciti nel materiale proiettato.
\paragraph{Riferimento fonte} File: Lezione 01 - Dipendenze funzionali e Normalizzazione di DB Relazionali.pdf. Numero slide: 63.
\section{Equivalenza di insiemi di dipendenze --- Slide 64}
\subsection*{Estratto originale dalla slide 64}
\begin{lstlisting}
Equivalenza di insiemi di dipendenze
funzionali
-    Definizione:
      Un insieme F di dipendenze funzionali copre un altro
      insieme E di dipendenze funzionali, se ogni DF in E è
      presente anche in F +, cioè se ogni dipendenza in E può
      essere inferita a partire da F.
-    Definizione:
      Due insiemi E ed F di dipendenze funzionali sono
      equivalenti se E+ = F+; ossia E è equivalente ad F se
      sussistono entrambe le condizioni: E copre F e F
      copre E.
-    Si può determinare se F copre E calcolando X+
    rispetto a F per ogni DF X -> Y in E, e quindi
    verificando se questo X+ comprende gli attributi
    presenti in Y.
\end{lstlisting}
\subsection*{Spiegazione estesa}
[AGGIUNTA] Questa spiegazione collega il contenuto testuale della slide al contesto del corso. La slide enfatizza i concetti mostrati nell'estratto e il loro ruolo nella progettazione di basi di dati relazionali, con particolare attenzione a correttezza semantica, qualità dello schema e compromessi progettuali.
[AGGIUNTA] Fonte integrazione: riformulazione didattica del testo presente nella slide 64 del file Lezione 01 - Dipendenze funzionali e Normalizzazione di DB Relazionali.pdf.
\subsection*{Derivazioni}
[AGGIUNTA] Nella slide compaiono espressioni simboliche/relazioni. Derivazione didattica generale:
\begin{align*}
X^+ &= X \\
X^+ &\leftarrow X^+ \cup Y \quad \text{se } (U \to V) \in F \text{ e } U \subseteq X^+ \text{ con } Y=V \\
\text{Stop} &\text{ quando non si aggiungono nuovi attributi in } X^+
\end{align*}
[AGGIUNTA] Questa procedura esplicita il metodo di chiusura usato nelle slide su dipendenze e chiavi.
\begin{itemize}[leftmargin=*]
\item Forma osservata nella slide: \texttt{equivalenti se E+ = F+; ossia E è equivalente ad F se}
\item Forma osservata nella slide: \texttt{rispetto a F per ogni DF X -> Y in E, e quindi}
\end{itemize}
\subsection*{Esempio numerico / esercizio}
[AGGIUNTA] Esempio: dato uno schema R(A,B,C) con dipendenze funzionali A->B e B->C, mostrare la chiusura A^+ e discutere se A è superchiave. Soluzione sintetica: A^+={A,B,C}, quindi A determina tutti gli attributi di R ed è superchiave.
\subsection*{Note su assunzioni o integrazioni}
[ASSUNZIONE] L'estrazione OCR può perdere formattazione o simboli della slide originale.
[AGGIUNTA] Le spiegazioni sono estensioni didattiche per rendere espliciti passaggi impliciti nel materiale proiettato.
\paragraph{Riferimento fonte} File: Lezione 01 - Dipendenze funzionali e Normalizzazione di DB Relazionali.pdf. Numero slide: 64.
\section{Forme Normali --- Slide 65}
\subsection*{Estratto originale dalla slide 65}
\begin{lstlisting}
Forme Normali
\end{lstlisting}
\subsection*{Spiegazione estesa}
[AGGIUNTA] Questa spiegazione collega il contenuto testuale della slide al contesto del corso. La slide enfatizza i concetti mostrati nell'estratto e il loro ruolo nella progettazione di basi di dati relazionali, con particolare attenzione a correttezza semantica, qualità dello schema e compromessi progettuali.
[AGGIUNTA] Fonte integrazione: riformulazione didattica del testo presente nella slide 65 del file Lezione 01 - Dipendenze funzionali e Normalizzazione di DB Relazionali.pdf.
\subsection*{Derivazioni}
[AGGIUNTA] In questa slide non sono state rilevate equazioni esplicite nel testo estratto; non è quindi richiesta una derivazione algebrica puntuale.
\subsection*{Esempio numerico / esercizio}
[AGGIUNTA] Esempio: dato uno schema R(A,B,C) con dipendenze funzionali A->B e B->C, mostrare la chiusura A^+ e discutere se A è superchiave. Soluzione sintetica: A^+={A,B,C}, quindi A determina tutti gli attributi di R ed è superchiave.
\subsection*{Note su assunzioni o integrazioni}
[ASSUNZIONE] L'estrazione OCR può perdere formattazione o simboli della slide originale.
[AGGIUNTA] Le spiegazioni sono estensioni didattiche per rendere espliciti passaggi impliciti nel materiale proiettato.
\paragraph{Riferimento fonte} File: Lezione 01 - Dipendenze funzionali e Normalizzazione di DB Relazionali.pdf. Numero slide: 65.
\section{Normalizzazione dei dati --- Slide 66}
\subsection*{Estratto originale dalla slide 66}
\begin{lstlisting}
Normalizzazione dei dati
-    La normalizzazione dei dati può essere vista
    come un processo che consente di decomporre
    schemi di relazione non ottimali in schemi più
    piccoli, tali da garantire la mancanza di "update
    anomalies".
\end{lstlisting}
\subsection*{Spiegazione estesa}
[AGGIUNTA] Questa spiegazione collega il contenuto testuale della slide al contesto del corso. La slide enfatizza i concetti mostrati nell'estratto e il loro ruolo nella progettazione di basi di dati relazionali, con particolare attenzione a correttezza semantica, qualità dello schema e compromessi progettuali.
[AGGIUNTA] Fonte integrazione: riformulazione didattica del testo presente nella slide 66 del file Lezione 01 - Dipendenze funzionali e Normalizzazione di DB Relazionali.pdf.
\subsection*{Derivazioni}
[AGGIUNTA] In questa slide non sono state rilevate equazioni esplicite nel testo estratto; non è quindi richiesta una derivazione algebrica puntuale.
\subsection*{Esempio numerico / esercizio}
[AGGIUNTA] Esempio: dato uno schema R(A,B,C) con dipendenze funzionali A->B e B->C, mostrare la chiusura A^+ e discutere se A è superchiave. Soluzione sintetica: A^+={A,B,C}, quindi A determina tutti gli attributi di R ed è superchiave.
\subsection*{Note su assunzioni o integrazioni}
[ASSUNZIONE] L'estrazione OCR può perdere formattazione o simboli della slide originale.
[AGGIUNTA] Le spiegazioni sono estensioni didattiche per rendere espliciti passaggi impliciti nel materiale proiettato.
\paragraph{Riferimento fonte} File: Lezione 01 - Dipendenze funzionali e Normalizzazione di DB Relazionali.pdf. Numero slide: 66.
\section{Forme normali --- Slide 67}
\subsection*{Estratto originale dalla slide 67}
\begin{lstlisting}
Forme normali
-    Le forme normali forniscono
      Un ambito formale per analizzare schemi di
      relazione, basato sulle chiavi e sulle dipendenze
      funzionali tra attributi.
      Una serie di test da condurre sugli schemi di
      relazione individuali:
        se un test fallisce, la relazione che viola il test deve
       essere decomposta in relazioni che individualmente
       superano il test.
\end{lstlisting}
\subsection*{Spiegazione estesa}
[AGGIUNTA] Questa spiegazione collega il contenuto testuale della slide al contesto del corso. La slide enfatizza i concetti mostrati nell'estratto e il loro ruolo nella progettazione di basi di dati relazionali, con particolare attenzione a correttezza semantica, qualità dello schema e compromessi progettuali.
[AGGIUNTA] Fonte integrazione: riformulazione didattica del testo presente nella slide 67 del file Lezione 01 - Dipendenze funzionali e Normalizzazione di DB Relazionali.pdf.
\subsection*{Derivazioni}
[AGGIUNTA] In questa slide non sono state rilevate equazioni esplicite nel testo estratto; non è quindi richiesta una derivazione algebrica puntuale.
\subsection*{Esempio numerico / esercizio}
[AGGIUNTA] Esempio: dato uno schema R(A,B,C) con dipendenze funzionali A->B e B->C, mostrare la chiusura A^+ e discutere se A è superchiave. Soluzione sintetica: A^+={A,B,C}, quindi A determina tutti gli attributi di R ed è superchiave.
\subsection*{Note su assunzioni o integrazioni}
[ASSUNZIONE] L'estrazione OCR può perdere formattazione o simboli della slide originale.
[AGGIUNTA] Le spiegazioni sono estensioni didattiche per rendere espliciti passaggi impliciti nel materiale proiettato.
\paragraph{Riferimento fonte} File: Lezione 01 - Dipendenze funzionali e Normalizzazione di DB Relazionali.pdf. Numero slide: 67.
\section{Forme normali (2) --- Slide 68}
\subsection*{Estratto originale dalla slide 68}
\begin{lstlisting}
Forme normali (2)
-  Le prime tre forme normali (1NF, 2NF, 3NF) e la
  Boyce-Codd (BCNF) sono definite considerando
  solo vincoli di dipendenza funzionale e di chiave.
-  La 4NF considera la dipendenza multivalued.

-  La 5NF considera la join dependency.
\end{lstlisting}
\subsection*{Spiegazione estesa}
[AGGIUNTA] Questa spiegazione collega il contenuto testuale della slide al contesto del corso. La slide enfatizza i concetti mostrati nell'estratto e il loro ruolo nella progettazione di basi di dati relazionali, con particolare attenzione a correttezza semantica, qualità dello schema e compromessi progettuali.
[AGGIUNTA] Fonte integrazione: riformulazione didattica del testo presente nella slide 68 del file Lezione 01 - Dipendenze funzionali e Normalizzazione di DB Relazionali.pdf.
\subsection*{Derivazioni}
[AGGIUNTA] In questa slide non sono state rilevate equazioni esplicite nel testo estratto; non è quindi richiesta una derivazione algebrica puntuale.
\subsection*{Esempio numerico / esercizio}
[AGGIUNTA] Esempio: dato uno schema R(A,B,C) con dipendenze funzionali A->B e B->C, mostrare la chiusura A^+ e discutere se A è superchiave. Soluzione sintetica: A^+={A,B,C}, quindi A determina tutti gli attributi di R ed è superchiave.
\subsection*{Note su assunzioni o integrazioni}
[ASSUNZIONE] L'estrazione OCR può perdere formattazione o simboli della slide originale.
[AGGIUNTA] Le spiegazioni sono estensioni didattiche per rendere espliciti passaggi impliciti nel materiale proiettato.
\paragraph{Riferimento fonte} File: Lezione 01 - Dipendenze funzionali e Normalizzazione di DB Relazionali.pdf. Numero slide: 68.
\section{Definizioni --- Slide 69}
\subsection*{Estratto originale dalla slide 69}
\begin{lstlisting}
Definizioni
-  Una superchiave di uno schema di relazione
  R = {A1, A2, ...,An} è un insieme di attributi
  S ÍR con la proprietà che non esistono due tuple
  t1e t2 in qualche stato di relazione legale di r di R
  tali che t1[S] = t2[S].
-  Una chiave k è una superchiave con la proprietà

  aggiuntiva che la rimozione da k di un qualche
  attributo fa perdere a k la proprietà di
  superchiave.
\end{lstlisting}
\subsection*{Spiegazione estesa}
[AGGIUNTA] Questa spiegazione collega il contenuto testuale della slide al contesto del corso. La slide enfatizza i concetti mostrati nell'estratto e il loro ruolo nella progettazione di basi di dati relazionali, con particolare attenzione a correttezza semantica, qualità dello schema e compromessi progettuali.
[AGGIUNTA] Fonte integrazione: riformulazione didattica del testo presente nella slide 69 del file Lezione 01 - Dipendenze funzionali e Normalizzazione di DB Relazionali.pdf.
\subsection*{Derivazioni}
[AGGIUNTA] Nella slide compaiono espressioni simboliche/relazioni. Derivazione didattica generale:
\begin{align*}
X^+ &= X \\
X^+ &\leftarrow X^+ \cup Y \quad \text{se } (U \to V) \in F \text{ e } U \subseteq X^+ \text{ con } Y=V \\
\text{Stop} &\text{ quando non si aggiungono nuovi attributi in } X^+
\end{align*}
[AGGIUNTA] Questa procedura esplicita il metodo di chiusura usato nelle slide su dipendenze e chiavi.
\begin{itemize}[leftmargin=*]
\item Forma osservata nella slide: \texttt{R = (A1, A2, ...,An) è un insieme di attributi}
\item Forma osservata nella slide: \texttt{tali che t1[S] = t2[S].}
\end{itemize}
\subsection*{Esempio numerico / esercizio}
[AGGIUNTA] Esempio: dato uno schema R(A,B,C) con dipendenze funzionali A->B e B->C, mostrare la chiusura A^+ e discutere se A è superchiave. Soluzione sintetica: A^+={A,B,C}, quindi A determina tutti gli attributi di R ed è superchiave.
\subsection*{Note su assunzioni o integrazioni}
[ASSUNZIONE] L'estrazione OCR può perdere formattazione o simboli della slide originale.
[AGGIUNTA] Le spiegazioni sono estensioni didattiche per rendere espliciti passaggi impliciti nel materiale proiettato.
\paragraph{Riferimento fonte} File: Lezione 01 - Dipendenze funzionali e Normalizzazione di DB Relazionali.pdf. Numero slide: 69.
\section{Definizioni (2) --- Slide 70}
\subsection*{Estratto originale dalla slide 70}
\begin{lstlisting}
Definizioni (2)
-    Tutte le chiavi "minimali" sono dette chiavi
    candidate. Una tra esse viene scelta
    arbitrariamente ed è detta chiave primaria.

-    Un attributo A di uno schema di relazione R è
    primo se e solo se fa parte di almeno una chiave
    candidata di R.
      In caso contrario A è detto non-primo.
\end{lstlisting}
\subsection*{Spiegazione estesa}
[AGGIUNTA] Questa spiegazione collega il contenuto testuale della slide al contesto del corso. La slide enfatizza i concetti mostrati nell'estratto e il loro ruolo nella progettazione di basi di dati relazionali, con particolare attenzione a correttezza semantica, qualità dello schema e compromessi progettuali.
[AGGIUNTA] Fonte integrazione: riformulazione didattica del testo presente nella slide 70 del file Lezione 01 - Dipendenze funzionali e Normalizzazione di DB Relazionali.pdf.
\subsection*{Derivazioni}
[AGGIUNTA] In questa slide non sono state rilevate equazioni esplicite nel testo estratto; non è quindi richiesta una derivazione algebrica puntuale.
\subsection*{Esempio numerico / esercizio}
[AGGIUNTA] Esempio: dato uno schema R(A,B,C) con dipendenze funzionali A->B e B->C, mostrare la chiusura A^+ e discutere se A è superchiave. Soluzione sintetica: A^+={A,B,C}, quindi A determina tutti gli attributi di R ed è superchiave.
\subsection*{Note su assunzioni o integrazioni}
[ASSUNZIONE] L'estrazione OCR può perdere formattazione o simboli della slide originale.
[AGGIUNTA] Le spiegazioni sono estensioni didattiche per rendere espliciti passaggi impliciti nel materiale proiettato.
\paragraph{Riferimento fonte} File: Lezione 01 - Dipendenze funzionali e Normalizzazione di DB Relazionali.pdf. Numero slide: 70.
\section{Definizioni: Esempio --- Slide 71}
\subsection*{Estratto originale dalla slide 71}
\begin{lstlisting}
Definizioni: Esempio


-  {SSN} è una chiave
-  {SSN}, {SSN,ENAME}, {SSN, ENAME, BDATE}

  etc. sono superchiavi
-  {SSN} è l'unica chiave candidata

-  {SSN} è la chiave primaria.
\end{lstlisting}
\subsection*{Spiegazione estesa}
[AGGIUNTA] Questa spiegazione collega il contenuto testuale della slide al contesto del corso. La slide enfatizza i concetti mostrati nell'estratto e il loro ruolo nella progettazione di basi di dati relazionali, con particolare attenzione a correttezza semantica, qualità dello schema e compromessi progettuali.
[AGGIUNTA] Fonte integrazione: riformulazione didattica del testo presente nella slide 71 del file Lezione 01 - Dipendenze funzionali e Normalizzazione di DB Relazionali.pdf.
\subsection*{Derivazioni}
[AGGIUNTA] In questa slide non sono state rilevate equazioni esplicite nel testo estratto; non è quindi richiesta una derivazione algebrica puntuale.
\subsection*{Esempio numerico / esercizio}
[AGGIUNTA] Esempio: dato uno schema R(A,B,C) con dipendenze funzionali A->B e B->C, mostrare la chiusura A^+ e discutere se A è superchiave. Soluzione sintetica: A^+={A,B,C}, quindi A determina tutti gli attributi di R ed è superchiave.
\subsection*{Note su assunzioni o integrazioni}
[ASSUNZIONE] L'estrazione OCR può perdere formattazione o simboli della slide originale.
[AGGIUNTA] Le spiegazioni sono estensioni didattiche per rendere espliciti passaggi impliciti nel materiale proiettato.
\paragraph{Riferimento fonte} File: Lezione 01 - Dipendenze funzionali e Normalizzazione di DB Relazionali.pdf. Numero slide: 71.
\section{Esempio di attributi primi --- Slide 72}
\subsection*{Estratto originale dalla slide 72}
\begin{lstlisting}
Esempio di attributi primi
-  Nello schema di relazione
           Stradario = {Via, Comune, CAP}
-  si ha:

           {Via, CAP} è una chiave
           dato che CAP -> Comune

-  Quindi {Via, CAP} sono attributi primi.
-  Tuttavia, anche {Via, Comune} è chiave, quindi

  Comune è primo.
\end{lstlisting}
\subsection*{Spiegazione estesa}
[AGGIUNTA] Questa spiegazione collega il contenuto testuale della slide al contesto del corso. La slide enfatizza i concetti mostrati nell'estratto e il loro ruolo nella progettazione di basi di dati relazionali, con particolare attenzione a correttezza semantica, qualità dello schema e compromessi progettuali.
[AGGIUNTA] Fonte integrazione: riformulazione didattica del testo presente nella slide 72 del file Lezione 01 - Dipendenze funzionali e Normalizzazione di DB Relazionali.pdf.
\subsection*{Derivazioni}
[AGGIUNTA] Nella slide compaiono espressioni simboliche/relazioni. Derivazione didattica generale:
\begin{align*}
X^+ &= X \\
X^+ &\leftarrow X^+ \cup Y \quad \text{se } (U \to V) \in F \text{ e } U \subseteq X^+ \text{ con } Y=V \\
\text{Stop} &\text{ quando non si aggiungono nuovi attributi in } X^+
\end{align*}
[AGGIUNTA] Questa procedura esplicita il metodo di chiusura usato nelle slide su dipendenze e chiavi.
\begin{itemize}[leftmargin=*]
\item Forma osservata nella slide: \texttt{Stradario = (Via, Comune, CAP)}
\item Forma osservata nella slide: \texttt{dato che CAP -> Comune}
\end{itemize}
\subsection*{Esempio numerico / esercizio}
[AGGIUNTA] Esempio: dato uno schema R(A,B,C) con dipendenze funzionali A->B e B->C, mostrare la chiusura A^+ e discutere se A è superchiave. Soluzione sintetica: A^+={A,B,C}, quindi A determina tutti gli attributi di R ed è superchiave.
\subsection*{Note su assunzioni o integrazioni}
[ASSUNZIONE] L'estrazione OCR può perdere formattazione o simboli della slide originale.
[AGGIUNTA] Le spiegazioni sono estensioni didattiche per rendere espliciti passaggi impliciti nel materiale proiettato.
\paragraph{Riferimento fonte} File: Lezione 01 - Dipendenze funzionali e Normalizzazione di DB Relazionali.pdf. Numero slide: 72.
\section{Prima forma normale (1NF) --- Slide 73}
\subsection*{Estratto originale dalla slide 73}
\begin{lstlisting}
Prima forma normale (1NF)
-    La 1NF è stata definita per non consentire
    attributi multivalued, composti e loro
    combinazioni.
      Ora fa parte della normale definizione di relazione
      del modello relazionale.


-    Gli unici valori consentiti da 1NF sono valori
    atomici (o indivisibili).
\end{lstlisting}
\subsection*{Spiegazione estesa}
[AGGIUNTA] Questa spiegazione collega il contenuto testuale della slide al contesto del corso. La slide enfatizza i concetti mostrati nell'estratto e il loro ruolo nella progettazione di basi di dati relazionali, con particolare attenzione a correttezza semantica, qualità dello schema e compromessi progettuali.
[AGGIUNTA] Fonte integrazione: riformulazione didattica del testo presente nella slide 73 del file Lezione 01 - Dipendenze funzionali e Normalizzazione di DB Relazionali.pdf.
\subsection*{Derivazioni}
[AGGIUNTA] In questa slide non sono state rilevate equazioni esplicite nel testo estratto; non è quindi richiesta una derivazione algebrica puntuale.
\subsection*{Esempio numerico / esercizio}
[AGGIUNTA] Esempio: dato uno schema R(A,B,C) con dipendenze funzionali A->B e B->C, mostrare la chiusura A^+ e discutere se A è superchiave. Soluzione sintetica: A^+={A,B,C}, quindi A determina tutti gli attributi di R ed è superchiave.
\subsection*{Note su assunzioni o integrazioni}
[ASSUNZIONE] L'estrazione OCR può perdere formattazione o simboli della slide originale.
[AGGIUNTA] Le spiegazioni sono estensioni didattiche per rendere espliciti passaggi impliciti nel materiale proiettato.
\paragraph{Riferimento fonte} File: Lezione 01 - Dipendenze funzionali e Normalizzazione di DB Relazionali.pdf. Numero slide: 73.
\section{Prima forma normale: Esempio --- Slide 74}
\subsection*{Estratto originale dalla slide 74}
\begin{lstlisting}
Prima forma normale: Esempio



-    L'idea è di rimuovere DLOCATIONS che viola la 1NF e
    porlo in una relazione separata insieme con la chiave
    primaria.
    La chiave primaria di questa relazione è la combinazione
    di {DNUMBER, DLOCATION}.
\end{lstlisting}
\subsection*{Spiegazione estesa}
[AGGIUNTA] Questa spiegazione collega il contenuto testuale della slide al contesto del corso. La slide enfatizza i concetti mostrati nell'estratto e il loro ruolo nella progettazione di basi di dati relazionali, con particolare attenzione a correttezza semantica, qualità dello schema e compromessi progettuali.
[AGGIUNTA] Fonte integrazione: riformulazione didattica del testo presente nella slide 74 del file Lezione 01 - Dipendenze funzionali e Normalizzazione di DB Relazionali.pdf.
\subsection*{Derivazioni}
[AGGIUNTA] In questa slide non sono state rilevate equazioni esplicite nel testo estratto; non è quindi richiesta una derivazione algebrica puntuale.
\subsection*{Esempio numerico / esercizio}
[AGGIUNTA] Esempio: dato uno schema R(A,B,C) con dipendenze funzionali A->B e B->C, mostrare la chiusura A^+ e discutere se A è superchiave. Soluzione sintetica: A^+={A,B,C}, quindi A determina tutti gli attributi di R ed è superchiave.
\subsection*{Note su assunzioni o integrazioni}
[ASSUNZIONE] L'estrazione OCR può perdere formattazione o simboli della slide originale.
[AGGIUNTA] Le spiegazioni sono estensioni didattiche per rendere espliciti passaggi impliciti nel materiale proiettato.
\paragraph{Riferimento fonte} File: Lezione 01 - Dipendenze funzionali e Normalizzazione di DB Relazionali.pdf. Numero slide: 74.
\section{Prima forma normale: Esempio (2) --- Slide 75}
\subsection*{Estratto originale dalla slide 75}
\begin{lstlisting}
Prima forma normale: Esempio (2)
-    Un secondo modo di normalizzare è di avere una
    tupla per ogni locazione; in tal caso la chiave
    primaria è {DNUMBER, DLOCATION}, ma si ha
    ridondanza tra le tuple.
\end{lstlisting}
\subsection*{Spiegazione estesa}
[AGGIUNTA] Questa spiegazione collega il contenuto testuale della slide al contesto del corso. La slide enfatizza i concetti mostrati nell'estratto e il loro ruolo nella progettazione di basi di dati relazionali, con particolare attenzione a correttezza semantica, qualità dello schema e compromessi progettuali.
[AGGIUNTA] Fonte integrazione: riformulazione didattica del testo presente nella slide 75 del file Lezione 01 - Dipendenze funzionali e Normalizzazione di DB Relazionali.pdf.
\subsection*{Derivazioni}
[AGGIUNTA] In questa slide non sono state rilevate equazioni esplicite nel testo estratto; non è quindi richiesta una derivazione algebrica puntuale.
\subsection*{Esempio numerico / esercizio}
[AGGIUNTA] Esempio: dato uno schema R(A,B,C) con dipendenze funzionali A->B e B->C, mostrare la chiusura A^+ e discutere se A è superchiave. Soluzione sintetica: A^+={A,B,C}, quindi A determina tutti gli attributi di R ed è superchiave.
\subsection*{Note su assunzioni o integrazioni}
[ASSUNZIONE] L'estrazione OCR può perdere formattazione o simboli della slide originale.
[AGGIUNTA] Le spiegazioni sono estensioni didattiche per rendere espliciti passaggi impliciti nel materiale proiettato.
\paragraph{Riferimento fonte} File: Lezione 01 - Dipendenze funzionali e Normalizzazione di DB Relazionali.pdf. Numero slide: 75.
\section{Prima forma normale: Esempio (3) --- Slide 76}
\subsection*{Estratto originale dalla slide 76}
\begin{lstlisting}
Prima forma normale: Esempio (3)
-    Un terzo modo consiste nello stabilire un
    massimo numero di valori per l'attributo
    DLOCATION - ad es. 3 - e rimpiazzare l'attributo
    con i 3 attributi atomici DLOCATION1,
    DLOCATION2 e DLOCATION3.

SVANTAGGIO: introduzione di valori nulli.
\end{lstlisting}
\subsection*{Spiegazione estesa}
[AGGIUNTA] Questa spiegazione collega il contenuto testuale della slide al contesto del corso. La slide enfatizza i concetti mostrati nell'estratto e il loro ruolo nella progettazione di basi di dati relazionali, con particolare attenzione a correttezza semantica, qualità dello schema e compromessi progettuali.
[AGGIUNTA] Fonte integrazione: riformulazione didattica del testo presente nella slide 76 del file Lezione 01 - Dipendenze funzionali e Normalizzazione di DB Relazionali.pdf.
\subsection*{Derivazioni}
[AGGIUNTA] In questa slide non sono state rilevate equazioni esplicite nel testo estratto; non è quindi richiesta una derivazione algebrica puntuale.
\subsection*{Esempio numerico / esercizio}
[AGGIUNTA] Esempio: dato uno schema R(A,B,C) con dipendenze funzionali A->B e B->C, mostrare la chiusura A^+ e discutere se A è superchiave. Soluzione sintetica: A^+={A,B,C}, quindi A determina tutti gli attributi di R ed è superchiave.
\subsection*{Note su assunzioni o integrazioni}
[ASSUNZIONE] L'estrazione OCR può perdere formattazione o simboli della slide originale.
[AGGIUNTA] Le spiegazioni sono estensioni didattiche per rendere espliciti passaggi impliciti nel materiale proiettato.
\paragraph{Riferimento fonte} File: Lezione 01 - Dipendenze funzionali e Normalizzazione di DB Relazionali.pdf. Numero slide: 76.
\section{Prima forma normale: Esempio 2 --- Slide 77}
\subsection*{Estratto originale dalla slide 77}
\begin{lstlisting}
Prima forma normale: Esempio 2
-    La 1NF proibisce anche attributi composti
    (relazioni annidate):




        SSN è chiave
        primaria.
        PNUMBER è
        chiave primaria
        parziale della
        relazione
        annidata.
\end{lstlisting}
\subsection*{Spiegazione estesa}
[AGGIUNTA] Questa spiegazione collega il contenuto testuale della slide al contesto del corso. La slide enfatizza i concetti mostrati nell'estratto e il loro ruolo nella progettazione di basi di dati relazionali, con particolare attenzione a correttezza semantica, qualità dello schema e compromessi progettuali.
[AGGIUNTA] Fonte integrazione: riformulazione didattica del testo presente nella slide 77 del file Lezione 01 - Dipendenze funzionali e Normalizzazione di DB Relazionali.pdf.
\subsection*{Derivazioni}
[AGGIUNTA] In questa slide non sono state rilevate equazioni esplicite nel testo estratto; non è quindi richiesta una derivazione algebrica puntuale.
\subsection*{Esempio numerico / esercizio}
[AGGIUNTA] Esempio: dato uno schema R(A,B,C) con dipendenze funzionali A->B e B->C, mostrare la chiusura A^+ e discutere se A è superchiave. Soluzione sintetica: A^+={A,B,C}, quindi A determina tutti gli attributi di R ed è superchiave.
\subsection*{Note su assunzioni o integrazioni}
[ASSUNZIONE] L'estrazione OCR può perdere formattazione o simboli della slide originale.
[AGGIUNTA] Le spiegazioni sono estensioni didattiche per rendere espliciti passaggi impliciti nel materiale proiettato.
\paragraph{Riferimento fonte} File: Lezione 01 - Dipendenze funzionali e Normalizzazione di DB Relazionali.pdf. Numero slide: 77.
\section{Prima forma normale: Esempio 2 (2) --- Slide 78}
\subsection*{Estratto originale dalla slide 78}
\begin{lstlisting}
Prima forma normale: Esempio 2 (2)
-    Per normalizzare:
      rimuovere gli attributi della relazione annidata in
     una nuova relazione e propagare la chiave primaria
     in essa.
\end{lstlisting}
\subsection*{Spiegazione estesa}
[AGGIUNTA] Questa spiegazione collega il contenuto testuale della slide al contesto del corso. La slide enfatizza i concetti mostrati nell'estratto e il loro ruolo nella progettazione di basi di dati relazionali, con particolare attenzione a correttezza semantica, qualità dello schema e compromessi progettuali.
[AGGIUNTA] Fonte integrazione: riformulazione didattica del testo presente nella slide 78 del file Lezione 01 - Dipendenze funzionali e Normalizzazione di DB Relazionali.pdf.
\subsection*{Derivazioni}
[AGGIUNTA] In questa slide non sono state rilevate equazioni esplicite nel testo estratto; non è quindi richiesta una derivazione algebrica puntuale.
\subsection*{Esempio numerico / esercizio}
[AGGIUNTA] Esempio: dato uno schema R(A,B,C) con dipendenze funzionali A->B e B->C, mostrare la chiusura A^+ e discutere se A è superchiave. Soluzione sintetica: A^+={A,B,C}, quindi A determina tutti gli attributi di R ed è superchiave.
\subsection*{Note su assunzioni o integrazioni}
[ASSUNZIONE] L'estrazione OCR può perdere formattazione o simboli della slide originale.
[AGGIUNTA] Le spiegazioni sono estensioni didattiche per rendere espliciti passaggi impliciti nel materiale proiettato.
\paragraph{Riferimento fonte} File: Lezione 01 - Dipendenze funzionali e Normalizzazione di DB Relazionali.pdf. Numero slide: 78.
\section{Seconda Forma Normale (2NF) --- Slide 79}
\subsection*{Estratto originale dalla slide 79}
\begin{lstlisting}
Seconda Forma Normale (2NF)




                  x
\end{lstlisting}
\subsection*{Spiegazione estesa}
[AGGIUNTA] Questa spiegazione collega il contenuto testuale della slide al contesto del corso. La slide enfatizza i concetti mostrati nell'estratto e il loro ruolo nella progettazione di basi di dati relazionali, con particolare attenzione a correttezza semantica, qualità dello schema e compromessi progettuali.
[AGGIUNTA] Fonte integrazione: riformulazione didattica del testo presente nella slide 79 del file Lezione 01 - Dipendenze funzionali e Normalizzazione di DB Relazionali.pdf.
\subsection*{Derivazioni}
[AGGIUNTA] In questa slide non sono state rilevate equazioni esplicite nel testo estratto; non è quindi richiesta una derivazione algebrica puntuale.
\subsection*{Esempio numerico / esercizio}
[AGGIUNTA] Esempio: dato uno schema R(A,B,C) con dipendenze funzionali A->B e B->C, mostrare la chiusura A^+ e discutere se A è superchiave. Soluzione sintetica: A^+={A,B,C}, quindi A determina tutti gli attributi di R ed è superchiave.
\subsection*{Note su assunzioni o integrazioni}
[ASSUNZIONE] L'estrazione OCR può perdere formattazione o simboli della slide originale.
[AGGIUNTA] Le spiegazioni sono estensioni didattiche per rendere espliciti passaggi impliciti nel materiale proiettato.
\paragraph{Riferimento fonte} File: Lezione 01 - Dipendenze funzionali e Normalizzazione di DB Relazionali.pdf. Numero slide: 79.
\section{Esempio dipendenze piene --- Slide 80}
\subsection*{Estratto originale dalla slide 80}
\begin{lstlisting}
Esempio dipendenze piene




-    {SSN, PNUMBER} -> HOURS: piena infatti
      SSN ->
          x HOURS
      PNUMBER ->
              x HOURS
-    {SSN, PNUMBER} -> ENAME: parziale infatti
      SSN -> ENAME
\end{lstlisting}
\subsection*{Spiegazione estesa}
[AGGIUNTA] Questa spiegazione collega il contenuto testuale della slide al contesto del corso. La slide enfatizza i concetti mostrati nell'estratto e il loro ruolo nella progettazione di basi di dati relazionali, con particolare attenzione a correttezza semantica, qualità dello schema e compromessi progettuali.
[AGGIUNTA] Fonte integrazione: riformulazione didattica del testo presente nella slide 80 del file Lezione 01 - Dipendenze funzionali e Normalizzazione di DB Relazionali.pdf.
\subsection*{Derivazioni}
[AGGIUNTA] Nella slide compaiono espressioni simboliche/relazioni. Derivazione didattica generale:
\begin{align*}
X^+ &= X \\
X^+ &\leftarrow X^+ \cup Y \quad \text{se } (U \to V) \in F \text{ e } U \subseteq X^+ \text{ con } Y=V \\
\text{Stop} &\text{ quando non si aggiungono nuovi attributi in } X^+
\end{align*}
[AGGIUNTA] Questa procedura esplicita il metodo di chiusura usato nelle slide su dipendenze e chiavi.
\begin{itemize}[leftmargin=*]
\item Forma osservata nella slide: \texttt{-    (SSN, PNUMBER) -> HOURS: piena infatti}
\item Forma osservata nella slide: \texttt{SSN ->}
\item Forma osservata nella slide: \texttt{PNUMBER ->}
\end{itemize}
\subsection*{Esempio numerico / esercizio}
[AGGIUNTA] Esempio: dato uno schema R(A,B,C) con dipendenze funzionali A->B e B->C, mostrare la chiusura A^+ e discutere se A è superchiave. Soluzione sintetica: A^+={A,B,C}, quindi A determina tutti gli attributi di R ed è superchiave.
\subsection*{Note su assunzioni o integrazioni}
[ASSUNZIONE] L'estrazione OCR può perdere formattazione o simboli della slide originale.
[AGGIUNTA] Le spiegazioni sono estensioni didattiche per rendere espliciti passaggi impliciti nel materiale proiettato.
\paragraph{Riferimento fonte} File: Lezione 01 - Dipendenze funzionali e Normalizzazione di DB Relazionali.pdf. Numero slide: 80.
\section{Seconda Forma Normale (2NF) --- Slide 81}
\subsection*{Estratto originale dalla slide 81}
\begin{lstlisting}
Seconda Forma Normale (2NF)
-    Uno schema di relazione R è in 2NF se:
      R è in prima forma normale.

      ogni attributo non-primo A in R è pienamente
     funzionalmente dipendente dalla chiave primaria di
     R.
\end{lstlisting}
\subsection*{Spiegazione estesa}
[AGGIUNTA] Questa spiegazione collega il contenuto testuale della slide al contesto del corso. La slide enfatizza i concetti mostrati nell'estratto e il loro ruolo nella progettazione di basi di dati relazionali, con particolare attenzione a correttezza semantica, qualità dello schema e compromessi progettuali.
[AGGIUNTA] Fonte integrazione: riformulazione didattica del testo presente nella slide 81 del file Lezione 01 - Dipendenze funzionali e Normalizzazione di DB Relazionali.pdf.
\subsection*{Derivazioni}
[AGGIUNTA] In questa slide non sono state rilevate equazioni esplicite nel testo estratto; non è quindi richiesta una derivazione algebrica puntuale.
\subsection*{Esempio numerico / esercizio}
[AGGIUNTA] Esempio: dato uno schema R(A,B,C) con dipendenze funzionali A->B e B->C, mostrare la chiusura A^+ e discutere se A è superchiave. Soluzione sintetica: A^+={A,B,C}, quindi A determina tutti gli attributi di R ed è superchiave.
\subsection*{Note su assunzioni o integrazioni}
[ASSUNZIONE] L'estrazione OCR può perdere formattazione o simboli della slide originale.
[AGGIUNTA] Le spiegazioni sono estensioni didattiche per rendere espliciti passaggi impliciti nel materiale proiettato.
\paragraph{Riferimento fonte} File: Lezione 01 - Dipendenze funzionali e Normalizzazione di DB Relazionali.pdf. Numero slide: 81.
\section{Seconda Forma Normale: Esempio --- Slide 82}
\subsection*{Estratto originale dalla slide 82}
\begin{lstlisting}
Seconda Forma Normale: Esempio




-  EMP_PROJ è in 1NF ma non in 2NF.
-  Infatti gli attributi non primi:
      PNAME, PLOCATION violano 2NF in virtù di fd3.
      ENAME viola 2NF in virtù di fd2.
-    Le dipendenze funzionali fd2 e fd3 rendono
    ENAME, PNAME, PLOCATION parzialmente
    dipendenti dalla chiave primaria
    {SSN, PNUMBER}
\end{lstlisting}
\subsection*{Spiegazione estesa}
[AGGIUNTA] Questa spiegazione collega il contenuto testuale della slide al contesto del corso. La slide enfatizza i concetti mostrati nell'estratto e il loro ruolo nella progettazione di basi di dati relazionali, con particolare attenzione a correttezza semantica, qualità dello schema e compromessi progettuali.
[AGGIUNTA] Fonte integrazione: riformulazione didattica del testo presente nella slide 82 del file Lezione 01 - Dipendenze funzionali e Normalizzazione di DB Relazionali.pdf.
\subsection*{Derivazioni}
[AGGIUNTA] In questa slide non sono state rilevate equazioni esplicite nel testo estratto; non è quindi richiesta una derivazione algebrica puntuale.
\subsection*{Esempio numerico / esercizio}
[AGGIUNTA] Esempio: dato uno schema R(A,B,C) con dipendenze funzionali A->B e B->C, mostrare la chiusura A^+ e discutere se A è superchiave. Soluzione sintetica: A^+={A,B,C}, quindi A determina tutti gli attributi di R ed è superchiave.
\subsection*{Note su assunzioni o integrazioni}
[ASSUNZIONE] L'estrazione OCR può perdere formattazione o simboli della slide originale.
[AGGIUNTA] Le spiegazioni sono estensioni didattiche per rendere espliciti passaggi impliciti nel materiale proiettato.
\paragraph{Riferimento fonte} File: Lezione 01 - Dipendenze funzionali e Normalizzazione di DB Relazionali.pdf. Numero slide: 82.
\section{Seconda Forma Normale --- Slide 83}
\subsection*{Estratto originale dalla slide 83}
\begin{lstlisting}
Seconda Forma Normale
-    Se uno schema di relazione non è in 2NF, può
    essere ulteriormente normalizzato in un certo
    numero di relazioni in 2NF, in cui gli attributi non
    primi sono associati solo con la parte di chiave
    primaria da cui sono funzionalmente dipendenti
    pienamente.
\end{lstlisting}
\subsection*{Spiegazione estesa}
[AGGIUNTA] Questa spiegazione collega il contenuto testuale della slide al contesto del corso. La slide enfatizza i concetti mostrati nell'estratto e il loro ruolo nella progettazione di basi di dati relazionali, con particolare attenzione a correttezza semantica, qualità dello schema e compromessi progettuali.
[AGGIUNTA] Fonte integrazione: riformulazione didattica del testo presente nella slide 83 del file Lezione 01 - Dipendenze funzionali e Normalizzazione di DB Relazionali.pdf.
\subsection*{Derivazioni}
[AGGIUNTA] In questa slide non sono state rilevate equazioni esplicite nel testo estratto; non è quindi richiesta una derivazione algebrica puntuale.
\subsection*{Esempio numerico / esercizio}
[AGGIUNTA] Esempio: dato uno schema R(A,B,C) con dipendenze funzionali A->B e B->C, mostrare la chiusura A^+ e discutere se A è superchiave. Soluzione sintetica: A^+={A,B,C}, quindi A determina tutti gli attributi di R ed è superchiave.
\subsection*{Note su assunzioni o integrazioni}
[ASSUNZIONE] L'estrazione OCR può perdere formattazione o simboli della slide originale.
[AGGIUNTA] Le spiegazioni sono estensioni didattiche per rendere espliciti passaggi impliciti nel materiale proiettato.
\paragraph{Riferimento fonte} File: Lezione 01 - Dipendenze funzionali e Normalizzazione di DB Relazionali.pdf. Numero slide: 83.
\section{Esempio di 2NF --- Slide 84}
\subsection*{Estratto originale dalla slide 84}
\begin{lstlisting}
Esempio di 2NF




     Normalizzazione in 2NF della relazione EMP_PROJ
\end{lstlisting}
\subsection*{Spiegazione estesa}
[AGGIUNTA] Questa spiegazione collega il contenuto testuale della slide al contesto del corso. La slide enfatizza i concetti mostrati nell'estratto e il loro ruolo nella progettazione di basi di dati relazionali, con particolare attenzione a correttezza semantica, qualità dello schema e compromessi progettuali.
[AGGIUNTA] Fonte integrazione: riformulazione didattica del testo presente nella slide 84 del file Lezione 01 - Dipendenze funzionali e Normalizzazione di DB Relazionali.pdf.
\subsection*{Derivazioni}
[AGGIUNTA] In questa slide non sono state rilevate equazioni esplicite nel testo estratto; non è quindi richiesta una derivazione algebrica puntuale.
\subsection*{Esempio numerico / esercizio}
[AGGIUNTA] Esempio: dato uno schema R(A,B,C) con dipendenze funzionali A->B e B->C, mostrare la chiusura A^+ e discutere se A è superchiave. Soluzione sintetica: A^+={A,B,C}, quindi A determina tutti gli attributi di R ed è superchiave.
\subsection*{Note su assunzioni o integrazioni}
[ASSUNZIONE] L'estrazione OCR può perdere formattazione o simboli della slide originale.
[AGGIUNTA] Le spiegazioni sono estensioni didattiche per rendere espliciti passaggi impliciti nel materiale proiettato.
\paragraph{Riferimento fonte} File: Lezione 01 - Dipendenze funzionali e Normalizzazione di DB Relazionali.pdf. Numero slide: 84.
\section{Terza Forma Normale (3NF) --- Slide 85}
\subsection*{Estratto originale dalla slide 85}
\begin{lstlisting}
Terza Forma Normale (3NF)
-    È basata sul concetto di dipendenza transitiva.

-    Una dipendenza funzionale X -> Y in uno schema
    R è una dipendenza transitiva se esiste un
    insieme di attributi Z che non è sottoinsieme di
    alcuna chiave di R e valgono X -> Z e Z -> Y.
\end{lstlisting}
\subsection*{Spiegazione estesa}
[AGGIUNTA] Questa spiegazione collega il contenuto testuale della slide al contesto del corso. La slide enfatizza i concetti mostrati nell'estratto e il loro ruolo nella progettazione di basi di dati relazionali, con particolare attenzione a correttezza semantica, qualità dello schema e compromessi progettuali.
[AGGIUNTA] Fonte integrazione: riformulazione didattica del testo presente nella slide 85 del file Lezione 01 - Dipendenze funzionali e Normalizzazione di DB Relazionali.pdf.
\subsection*{Derivazioni}
[AGGIUNTA] Nella slide compaiono espressioni simboliche/relazioni. Derivazione didattica generale:
\begin{align*}
X^+ &= X \\
X^+ &\leftarrow X^+ \cup Y \quad \text{se } (U \to V) \in F \text{ e } U \subseteq X^+ \text{ con } Y=V \\
\text{Stop} &\text{ quando non si aggiungono nuovi attributi in } X^+
\end{align*}
[AGGIUNTA] Questa procedura esplicita il metodo di chiusura usato nelle slide su dipendenze e chiavi.
\begin{itemize}[leftmargin=*]
\item Forma osservata nella slide: \texttt{-    Una dipendenza funzionale X -> Y in uno schema}
\item Forma osservata nella slide: \texttt{alcuna chiave di R e valgono X -> Z e Z -> Y.}
\end{itemize}
\subsection*{Esempio numerico / esercizio}
[AGGIUNTA] Esempio: dato uno schema R(A,B,C) con dipendenze funzionali A->B e B->C, mostrare la chiusura A^+ e discutere se A è superchiave. Soluzione sintetica: A^+={A,B,C}, quindi A determina tutti gli attributi di R ed è superchiave.
\subsection*{Note su assunzioni o integrazioni}
[ASSUNZIONE] L'estrazione OCR può perdere formattazione o simboli della slide originale.
[AGGIUNTA] Le spiegazioni sono estensioni didattiche per rendere espliciti passaggi impliciti nel materiale proiettato.
\paragraph{Riferimento fonte} File: Lezione 01 - Dipendenze funzionali e Normalizzazione di DB Relazionali.pdf. Numero slide: 85.
\section{Terza Forma Normale: Esempio --- Slide 86}
\subsection*{Estratto originale dalla slide 86}
\begin{lstlisting}
Terza Forma Normale: Esempio




-    SSN -> DMGRSSN è transitiva attraverso
    DNUMBER infatti
      SSN -> DNUMBER
      DNUMBER -> DMGRSSN
      DNUMBER non è sottoinsieme della chiave di
     EMP_DEPT
\end{lstlisting}
\subsection*{Spiegazione estesa}
[AGGIUNTA] Questa spiegazione collega il contenuto testuale della slide al contesto del corso. La slide enfatizza i concetti mostrati nell'estratto e il loro ruolo nella progettazione di basi di dati relazionali, con particolare attenzione a correttezza semantica, qualità dello schema e compromessi progettuali.
[AGGIUNTA] Fonte integrazione: riformulazione didattica del testo presente nella slide 86 del file Lezione 01 - Dipendenze funzionali e Normalizzazione di DB Relazionali.pdf.
\subsection*{Derivazioni}
[AGGIUNTA] Nella slide compaiono espressioni simboliche/relazioni. Derivazione didattica generale:
\begin{align*}
X^+ &= X \\
X^+ &\leftarrow X^+ \cup Y \quad \text{se } (U \to V) \in F \text{ e } U \subseteq X^+ \text{ con } Y=V \\
\text{Stop} &\text{ quando non si aggiungono nuovi attributi in } X^+
\end{align*}
[AGGIUNTA] Questa procedura esplicita il metodo di chiusura usato nelle slide su dipendenze e chiavi.
\begin{itemize}[leftmargin=*]
\item Forma osservata nella slide: \texttt{-    SSN -> DMGRSSN è transitiva attraverso}
\item Forma osservata nella slide: \texttt{SSN -> DNUMBER}
\item Forma osservata nella slide: \texttt{DNUMBER -> DMGRSSN}
\end{itemize}
\subsection*{Esempio numerico / esercizio}
[AGGIUNTA] Esempio: dato uno schema R(A,B,C) con dipendenze funzionali A->B e B->C, mostrare la chiusura A^+ e discutere se A è superchiave. Soluzione sintetica: A^+={A,B,C}, quindi A determina tutti gli attributi di R ed è superchiave.
\subsection*{Note su assunzioni o integrazioni}
[ASSUNZIONE] L'estrazione OCR può perdere formattazione o simboli della slide originale.
[AGGIUNTA] Le spiegazioni sono estensioni didattiche per rendere espliciti passaggi impliciti nel materiale proiettato.
\paragraph{Riferimento fonte} File: Lezione 01 - Dipendenze funzionali e Normalizzazione di DB Relazionali.pdf. Numero slide: 86.
\section{Terza Forma Normale: Esempio (2) --- Slide 87}
\subsection*{Estratto originale dalla slide 87}
\begin{lstlisting}
Terza Forma Normale: Esempio (2)
\end{lstlisting}
\subsection*{Spiegazione estesa}
[AGGIUNTA] Questa spiegazione collega il contenuto testuale della slide al contesto del corso. La slide enfatizza i concetti mostrati nell'estratto e il loro ruolo nella progettazione di basi di dati relazionali, con particolare attenzione a correttezza semantica, qualità dello schema e compromessi progettuali.
[AGGIUNTA] Fonte integrazione: riformulazione didattica del testo presente nella slide 87 del file Lezione 01 - Dipendenze funzionali e Normalizzazione di DB Relazionali.pdf.
\subsection*{Derivazioni}
[AGGIUNTA] In questa slide non sono state rilevate equazioni esplicite nel testo estratto; non è quindi richiesta una derivazione algebrica puntuale.
\subsection*{Esempio numerico / esercizio}
[AGGIUNTA] Esempio: dato uno schema R(A,B,C) con dipendenze funzionali A->B e B->C, mostrare la chiusura A^+ e discutere se A è superchiave. Soluzione sintetica: A^+={A,B,C}, quindi A determina tutti gli attributi di R ed è superchiave.
\subsection*{Note su assunzioni o integrazioni}
[ASSUNZIONE] L'estrazione OCR può perdere formattazione o simboli della slide originale.
[AGGIUNTA] Le spiegazioni sono estensioni didattiche per rendere espliciti passaggi impliciti nel materiale proiettato.
\paragraph{Riferimento fonte} File: Lezione 01 - Dipendenze funzionali e Normalizzazione di DB Relazionali.pdf. Numero slide: 87.
\section{Terza Forma Normale --- Slide 88}
\subsection*{Estratto originale dalla slide 88}
\begin{lstlisting}
Terza Forma Normale
-    Secondo le definizione di Codd:
    uno schema di relazione R è in 3NF se:
       è in 2NF.
      nessun attributo non-primo di R è transitivamente
      dipendente dalla chiave primaria.
\end{lstlisting}
\subsection*{Spiegazione estesa}
[AGGIUNTA] Questa spiegazione collega il contenuto testuale della slide al contesto del corso. La slide enfatizza i concetti mostrati nell'estratto e il loro ruolo nella progettazione di basi di dati relazionali, con particolare attenzione a correttezza semantica, qualità dello schema e compromessi progettuali.
[AGGIUNTA] Fonte integrazione: riformulazione didattica del testo presente nella slide 88 del file Lezione 01 - Dipendenze funzionali e Normalizzazione di DB Relazionali.pdf.
\subsection*{Derivazioni}
[AGGIUNTA] In questa slide non sono state rilevate equazioni esplicite nel testo estratto; non è quindi richiesta una derivazione algebrica puntuale.
\subsection*{Esempio numerico / esercizio}
[AGGIUNTA] Esempio: dato uno schema R(A,B,C) con dipendenze funzionali A->B e B->C, mostrare la chiusura A^+ e discutere se A è superchiave. Soluzione sintetica: A^+={A,B,C}, quindi A determina tutti gli attributi di R ed è superchiave.
\subsection*{Note su assunzioni o integrazioni}
[ASSUNZIONE] L'estrazione OCR può perdere formattazione o simboli della slide originale.
[AGGIUNTA] Le spiegazioni sono estensioni didattiche per rendere espliciti passaggi impliciti nel materiale proiettato.
\paragraph{Riferimento fonte} File: Lezione 01 - Dipendenze funzionali e Normalizzazione di DB Relazionali.pdf. Numero slide: 88.
\section{Sommario --- Slide 89}
\subsection*{Estratto originale dalla slide 89}
\begin{lstlisting}
Sommario
  Forma normale   Test                            Normalizzazione

  1NF             La relazione non deve           Forma una nuova relazione
                  presentare attributi non        per ogni attributo non
                  atomici.                        atomico.

  2NF             La chiave primaria              Crea una relazione per ogni
                  contiene attributi multipli e   chiave parziale con i suoi
                  attributi non chiave sono       attributi dipendenti.
                  funzionalmente                  (mantieni una relazione per
                  dipendenti da parte della       la chiave primaria e gli
                  chiave primaria.                attributi che dipendono
                                                  funzionalmente da questa)
  3NF             Non deve esistere una           Decomponi la relazione in
                  dipendenza transitiva tra       un insieme di relazioni che
                  un attributo non chiave e       includono gli attributi non
                  la chiave primaria.             chiave che funzionalmente
                                                  determinano altri attributi
                                                  non chiave.
\end{lstlisting}
\subsection*{Spiegazione estesa}
[AGGIUNTA] Questa spiegazione collega il contenuto testuale della slide al contesto del corso. La slide enfatizza i concetti mostrati nell'estratto e il loro ruolo nella progettazione di basi di dati relazionali, con particolare attenzione a correttezza semantica, qualità dello schema e compromessi progettuali.
[AGGIUNTA] Fonte integrazione: riformulazione didattica del testo presente nella slide 89 del file Lezione 01 - Dipendenze funzionali e Normalizzazione di DB Relazionali.pdf.
\subsection*{Derivazioni}
[AGGIUNTA] In questa slide non sono state rilevate equazioni esplicite nel testo estratto; non è quindi richiesta una derivazione algebrica puntuale.
\subsection*{Esempio numerico / esercizio}
[AGGIUNTA] Esempio: dato uno schema R(A,B,C) con dipendenze funzionali A->B e B->C, mostrare la chiusura A^+ e discutere se A è superchiave. Soluzione sintetica: A^+={A,B,C}, quindi A determina tutti gli attributi di R ed è superchiave.
\subsection*{Note su assunzioni o integrazioni}
[ASSUNZIONE] L'estrazione OCR può perdere formattazione o simboli della slide originale.
[AGGIUNTA] Le spiegazioni sono estensioni didattiche per rendere espliciti passaggi impliciti nel materiale proiettato.
\paragraph{Riferimento fonte} File: Lezione 01 - Dipendenze funzionali e Normalizzazione di DB Relazionali.pdf. Numero slide: 89.
\section{Slide 90 --- Slide 90}
\subsection*{Estratto originale dalla slide 90}
[INCERTO] Testo non estratto automaticamente da questa slide (possibile contenuto grafico o testo non OCR).
\subsection*{Spiegazione estesa}
[AGGIUNTA] Questa spiegazione collega il contenuto testuale della slide al contesto del corso. La slide enfatizza i concetti mostrati nell'estratto e il loro ruolo nella progettazione di basi di dati relazionali, con particolare attenzione a correttezza semantica, qualità dello schema e compromessi progettuali.
[AGGIUNTA] Fonte integrazione: riformulazione didattica del testo presente nella slide 90 del file Lezione 01 - Dipendenze funzionali e Normalizzazione di DB Relazionali.pdf.
[INCERTO] Poiché l'estratto è vuoto/ambiguo, l'interpretazione è limitata al contesto delle slide adiacenti.
\subsection*{Derivazioni}
[AGGIUNTA] In questa slide non sono state rilevate equazioni esplicite nel testo estratto; non è quindi richiesta una derivazione algebrica puntuale.
\subsection*{Esempio numerico / esercizio}
[AGGIUNTA] Esempio: dato uno schema R(A,B,C) con dipendenze funzionali A->B e B->C, mostrare la chiusura A^+ e discutere se A è superchiave. Soluzione sintetica: A^+={A,B,C}, quindi A determina tutti gli attributi di R ed è superchiave.
\subsection*{Note su assunzioni o integrazioni}
[ASSUNZIONE] L'estrazione OCR può perdere formattazione o simboli della slide originale.
[AGGIUNTA] Le spiegazioni sono estensioni didattiche per rendere espliciti passaggi impliciti nel materiale proiettato.
\paragraph{Riferimento fonte} File: Lezione 01 - Dipendenze funzionali e Normalizzazione di DB Relazionali.pdf. Numero slide: 90.
\section{Definizione generale di Seconda e Terza --- Slide 91}
\subsection*{Estratto originale dalla slide 91}
\begin{lstlisting}
Definizione generale di Seconda e Terza
Forma Normale
-  La definizione generale di 2NF e 3NF non riguarda
  solo la chiave primaria ma tutte le chiavi candidate.
-  Un attributo è primo se è parte di una qualsiasi
  chiave candidata.

  Uno schema di relazione R è in 2NF se ogni attributo
  non-primo A in R non è parzialmente dipendente da
  alcuna chiave di R.
  Uno schema di relazione R è in 3NF se, ogni volta
  che vale in R una dipendenza funzionale
  X -> A, o
      X è una superchiave di R, oppure
      A è un attributo primo di R.
\end{lstlisting}
\subsection*{Spiegazione estesa}
[AGGIUNTA] Questa spiegazione collega il contenuto testuale della slide al contesto del corso. La slide enfatizza i concetti mostrati nell'estratto e il loro ruolo nella progettazione di basi di dati relazionali, con particolare attenzione a correttezza semantica, qualità dello schema e compromessi progettuali.
[AGGIUNTA] Fonte integrazione: riformulazione didattica del testo presente nella slide 91 del file Lezione 01 - Dipendenze funzionali e Normalizzazione di DB Relazionali.pdf.
\subsection*{Derivazioni}
[AGGIUNTA] Nella slide compaiono espressioni simboliche/relazioni. Derivazione didattica generale:
\begin{align*}
X^+ &= X \\
X^+ &\leftarrow X^+ \cup Y \quad \text{se } (U \to V) \in F \text{ e } U \subseteq X^+ \text{ con } Y=V \\
\text{Stop} &\text{ quando non si aggiungono nuovi attributi in } X^+
\end{align*}
[AGGIUNTA] Questa procedura esplicita il metodo di chiusura usato nelle slide su dipendenze e chiavi.
\begin{itemize}[leftmargin=*]
\item Forma osservata nella slide: \texttt{X -> A, o}
\end{itemize}
\subsection*{Esempio numerico / esercizio}
[AGGIUNTA] Esempio: dato uno schema R(A,B,C) con dipendenze funzionali A->B e B->C, mostrare la chiusura A^+ e discutere se A è superchiave. Soluzione sintetica: A^+={A,B,C}, quindi A determina tutti gli attributi di R ed è superchiave.
\subsection*{Note su assunzioni o integrazioni}
[ASSUNZIONE] L'estrazione OCR può perdere formattazione o simboli della slide originale.
[AGGIUNTA] Le spiegazioni sono estensioni didattiche per rendere espliciti passaggi impliciti nel materiale proiettato.
\paragraph{Riferimento fonte} File: Lezione 01 - Dipendenze funzionali e Normalizzazione di DB Relazionali.pdf. Numero slide: 91.
\section{Forma Normale di Boyce-Codd (BCNF) --- Slide 92}
\subsection*{Estratto originale dalla slide 92}
\begin{lstlisting}
Forma Normale di Boyce-Codd (BCNF)

    Uno schema R è in BCNF se ogniqualvolta vale
    in R una dipendenza funzionale X -> A, allora X è
    una superchiave di R.




              Una relazione in 3NF, ma non in BCNF
\end{lstlisting}
\subsection*{Spiegazione estesa}
[AGGIUNTA] Questa spiegazione collega il contenuto testuale della slide al contesto del corso. La slide enfatizza i concetti mostrati nell'estratto e il loro ruolo nella progettazione di basi di dati relazionali, con particolare attenzione a correttezza semantica, qualità dello schema e compromessi progettuali.
[AGGIUNTA] Fonte integrazione: riformulazione didattica del testo presente nella slide 92 del file Lezione 01 - Dipendenze funzionali e Normalizzazione di DB Relazionali.pdf.
\subsection*{Derivazioni}
[AGGIUNTA] Nella slide compaiono espressioni simboliche/relazioni. Derivazione didattica generale:
\begin{align*}
X^+ &= X \\
X^+ &\leftarrow X^+ \cup Y \quad \text{se } (U \to V) \in F \text{ e } U \subseteq X^+ \text{ con } Y=V \\
\text{Stop} &\text{ quando non si aggiungono nuovi attributi in } X^+
\end{align*}
[AGGIUNTA] Questa procedura esplicita il metodo di chiusura usato nelle slide su dipendenze e chiavi.
\begin{itemize}[leftmargin=*]
\item Forma osservata nella slide: \texttt{in R una dipendenza funzionale X -> A, allora X è}
\end{itemize}
\subsection*{Esempio numerico / esercizio}
[AGGIUNTA] Esempio: dato uno schema R(A,B,C) con dipendenze funzionali A->B e B->C, mostrare la chiusura A^+ e discutere se A è superchiave. Soluzione sintetica: A^+={A,B,C}, quindi A determina tutti gli attributi di R ed è superchiave.
\subsection*{Note su assunzioni o integrazioni}
[ASSUNZIONE] L'estrazione OCR può perdere formattazione o simboli della slide originale.
[AGGIUNTA] Le spiegazioni sono estensioni didattiche per rendere espliciti passaggi impliciti nel materiale proiettato.
\paragraph{Riferimento fonte} File: Lezione 01 - Dipendenze funzionali e Normalizzazione di DB Relazionali.pdf. Numero slide: 92.
\section{Forme Normali: Esempio --- Slide 93}
\subsection*{Estratto originale dalla slide 93}
\begin{lstlisting}
Forme Normali: Esempio
-    Lo schema di relazione LOTS descrive lotti di terreno in
    vendita in varie contee di uno stato.




-  Chiavi candidate: PROPERTY_ID# e
                    {COUNTY_NAME, LOT#}
-  Chiave primaria: PROPERTY_ID#

-  Inoltre valgono:

      FD3:   COUNTY_NAME -> TAX-RATE
      FD4:   AREA -> PRICE
\end{lstlisting}
\subsection*{Spiegazione estesa}
[AGGIUNTA] Questa spiegazione collega il contenuto testuale della slide al contesto del corso. La slide enfatizza i concetti mostrati nell'estratto e il loro ruolo nella progettazione di basi di dati relazionali, con particolare attenzione a correttezza semantica, qualità dello schema e compromessi progettuali.
[AGGIUNTA] Fonte integrazione: riformulazione didattica del testo presente nella slide 93 del file Lezione 01 - Dipendenze funzionali e Normalizzazione di DB Relazionali.pdf.
\subsection*{Derivazioni}
[AGGIUNTA] Nella slide compaiono espressioni simboliche/relazioni. Derivazione didattica generale:
\begin{align*}
X^+ &= X \\
X^+ &\leftarrow X^+ \cup Y \quad \text{se } (U \to V) \in F \text{ e } U \subseteq X^+ \text{ con } Y=V \\
\text{Stop} &\text{ quando non si aggiungono nuovi attributi in } X^+
\end{align*}
[AGGIUNTA] Questa procedura esplicita il metodo di chiusura usato nelle slide su dipendenze e chiavi.
\begin{itemize}[leftmargin=*]
\item Forma osservata nella slide: \texttt{FD3:   COUNTY\_NAME -> TAX-RATE}
\item Forma osservata nella slide: \texttt{FD4:   AREA -> PRICE}
\end{itemize}
\subsection*{Esempio numerico / esercizio}
[AGGIUNTA] Esempio: dato uno schema R(A,B,C) con dipendenze funzionali A->B e B->C, mostrare la chiusura A^+ e discutere se A è superchiave. Soluzione sintetica: A^+={A,B,C}, quindi A determina tutti gli attributi di R ed è superchiave.
\subsection*{Note su assunzioni o integrazioni}
[ASSUNZIONE] L'estrazione OCR può perdere formattazione o simboli della slide originale.
[AGGIUNTA] Le spiegazioni sono estensioni didattiche per rendere espliciti passaggi impliciti nel materiale proiettato.
\paragraph{Riferimento fonte} File: Lezione 01 - Dipendenze funzionali e Normalizzazione di DB Relazionali.pdf. Numero slide: 93.
\section{Forme Normali: Esempio (2) --- Slide 94}
\subsection*{Estratto originale dalla slide 94}
\begin{lstlisting}
Forme Normali: Esempio (2)
-    LOTS viola la definizione generale di 2NF perché
    TAX_RATE è parzialmente dipendente dalla chiave
    candidata {COUNTY_NAME, LOT#} in virtù di fd3.
-    Per normalizzare decomponiamo LOTS in LOTS1 e
    LOTS2.
\end{lstlisting}
\subsection*{Spiegazione estesa}
[AGGIUNTA] Questa spiegazione collega il contenuto testuale della slide al contesto del corso. La slide enfatizza i concetti mostrati nell'estratto e il loro ruolo nella progettazione di basi di dati relazionali, con particolare attenzione a correttezza semantica, qualità dello schema e compromessi progettuali.
[AGGIUNTA] Fonte integrazione: riformulazione didattica del testo presente nella slide 94 del file Lezione 01 - Dipendenze funzionali e Normalizzazione di DB Relazionali.pdf.
\subsection*{Derivazioni}
[AGGIUNTA] In questa slide non sono state rilevate equazioni esplicite nel testo estratto; non è quindi richiesta una derivazione algebrica puntuale.
\subsection*{Esempio numerico / esercizio}
[AGGIUNTA] Esempio: dato uno schema R(A,B,C) con dipendenze funzionali A->B e B->C, mostrare la chiusura A^+ e discutere se A è superchiave. Soluzione sintetica: A^+={A,B,C}, quindi A determina tutti gli attributi di R ed è superchiave.
\subsection*{Note su assunzioni o integrazioni}
[ASSUNZIONE] L'estrazione OCR può perdere formattazione o simboli della slide originale.
[AGGIUNTA] Le spiegazioni sono estensioni didattiche per rendere espliciti passaggi impliciti nel materiale proiettato.
\paragraph{Riferimento fonte} File: Lezione 01 - Dipendenze funzionali e Normalizzazione di DB Relazionali.pdf. Numero slide: 94.
\section{Forme Normali: Esempio (3) --- Slide 95}
\subsection*{Estratto originale dalla slide 95}
\begin{lstlisting}
Forme Normali: Esempio (3)
-    In accordo alla definizione generale di 3NF,
    LOTS2 è in 3NF, mentre fd4 in LOTS1 viola la
    3NF.

-    Infatti AREA -> PRICE, ma
      AREA non è superchiave di LOTS1, e
      PRICE non è attributo primo.



-    Per normalizzare: decomponiamo LOTS1
    rimuovendo PRICE e ponendolo insieme ad
    AREA in un altro schema di relazione.
\end{lstlisting}
\subsection*{Spiegazione estesa}
[AGGIUNTA] Questa spiegazione collega il contenuto testuale della slide al contesto del corso. La slide enfatizza i concetti mostrati nell'estratto e il loro ruolo nella progettazione di basi di dati relazionali, con particolare attenzione a correttezza semantica, qualità dello schema e compromessi progettuali.
[AGGIUNTA] Fonte integrazione: riformulazione didattica del testo presente nella slide 95 del file Lezione 01 - Dipendenze funzionali e Normalizzazione di DB Relazionali.pdf.
\subsection*{Derivazioni}
[AGGIUNTA] Nella slide compaiono espressioni simboliche/relazioni. Derivazione didattica generale:
\begin{align*}
X^+ &= X \\
X^+ &\leftarrow X^+ \cup Y \quad \text{se } (U \to V) \in F \text{ e } U \subseteq X^+ \text{ con } Y=V \\
\text{Stop} &\text{ quando non si aggiungono nuovi attributi in } X^+
\end{align*}
[AGGIUNTA] Questa procedura esplicita il metodo di chiusura usato nelle slide su dipendenze e chiavi.
\begin{itemize}[leftmargin=*]
\item Forma osservata nella slide: \texttt{-    Infatti AREA -> PRICE, ma}
\end{itemize}
\subsection*{Esempio numerico / esercizio}
[AGGIUNTA] Esempio: dato uno schema R(A,B,C) con dipendenze funzionali A->B e B->C, mostrare la chiusura A^+ e discutere se A è superchiave. Soluzione sintetica: A^+={A,B,C}, quindi A determina tutti gli attributi di R ed è superchiave.
\subsection*{Note su assunzioni o integrazioni}
[ASSUNZIONE] L'estrazione OCR può perdere formattazione o simboli della slide originale.
[AGGIUNTA] Le spiegazioni sono estensioni didattiche per rendere espliciti passaggi impliciti nel materiale proiettato.
\paragraph{Riferimento fonte} File: Lezione 01 - Dipendenze funzionali e Normalizzazione di DB Relazionali.pdf. Numero slide: 95.
\section{Forme Normali: Esempio (4) --- Slide 96}
\subsection*{Estratto originale dalla slide 96}
\begin{lstlisting}
Forme Normali: Esempio (4)



           Decomposizione di LOTS1 in 3NF




           Sommario delle decomposizioni
\end{lstlisting}
\subsection*{Spiegazione estesa}
[AGGIUNTA] Questa spiegazione collega il contenuto testuale della slide al contesto del corso. La slide enfatizza i concetti mostrati nell'estratto e il loro ruolo nella progettazione di basi di dati relazionali, con particolare attenzione a correttezza semantica, qualità dello schema e compromessi progettuali.
[AGGIUNTA] Fonte integrazione: riformulazione didattica del testo presente nella slide 96 del file Lezione 01 - Dipendenze funzionali e Normalizzazione di DB Relazionali.pdf.
\subsection*{Derivazioni}
[AGGIUNTA] In questa slide non sono state rilevate equazioni esplicite nel testo estratto; non è quindi richiesta una derivazione algebrica puntuale.
\subsection*{Esempio numerico / esercizio}
[AGGIUNTA] Esempio: dato uno schema R(A,B,C) con dipendenze funzionali A->B e B->C, mostrare la chiusura A^+ e discutere se A è superchiave. Soluzione sintetica: A^+={A,B,C}, quindi A determina tutti gli attributi di R ed è superchiave.
\subsection*{Note su assunzioni o integrazioni}
[ASSUNZIONE] L'estrazione OCR può perdere formattazione o simboli della slide originale.
[AGGIUNTA] Le spiegazioni sono estensioni didattiche per rendere espliciti passaggi impliciti nel materiale proiettato.
\paragraph{Riferimento fonte} File: Lezione 01 - Dipendenze funzionali e Normalizzazione di DB Relazionali.pdf. Numero slide: 96.
\section{Forme Normali: Esempio (5) --- Slide 97}
\subsection*{Estratto originale dalla slide 97}
\begin{lstlisting}
Forme Normali: Esempio (5)
-     Supponiamo che ci siano migliaia di lotti: i lotti
     appartengono a solo due contee
       Marion County  con aree: 0,5 0,6 0,7 0,8 0,9 1,0
       Liberty County con aree: 1,1 1,2  ... 1,9 2,0




                                                   aggiungiamo fd5
    fd2 viene persa
\end{lstlisting}
\subsection*{Spiegazione estesa}
[AGGIUNTA] Questa spiegazione collega il contenuto testuale della slide al contesto del corso. La slide enfatizza i concetti mostrati nell'estratto e il loro ruolo nella progettazione di basi di dati relazionali, con particolare attenzione a correttezza semantica, qualità dello schema e compromessi progettuali.
[AGGIUNTA] Fonte integrazione: riformulazione didattica del testo presente nella slide 97 del file Lezione 01 - Dipendenze funzionali e Normalizzazione di DB Relazionali.pdf.
\subsection*{Derivazioni}
[AGGIUNTA] In questa slide non sono state rilevate equazioni esplicite nel testo estratto; non è quindi richiesta una derivazione algebrica puntuale.
\subsection*{Esempio numerico / esercizio}
[AGGIUNTA] Esempio: dato uno schema R(A,B,C) con dipendenze funzionali A->B e B->C, mostrare la chiusura A^+ e discutere se A è superchiave. Soluzione sintetica: A^+={A,B,C}, quindi A determina tutti gli attributi di R ed è superchiave.
\subsection*{Note su assunzioni o integrazioni}
[ASSUNZIONE] L'estrazione OCR può perdere formattazione o simboli della slide originale.
[AGGIUNTA] Le spiegazioni sono estensioni didattiche per rendere espliciti passaggi impliciti nel materiale proiettato.
\paragraph{Riferimento fonte} File: Lezione 01 - Dipendenze funzionali e Normalizzazione di DB Relazionali.pdf. Numero slide: 97.
\section{Forme Normali: Esempio (6) --- Slide 98}
\subsection*{Estratto originale dalla slide 98}
\begin{lstlisting}
Forme Normali: Esempio (6)

                                                     S     C    I




                                             {STUDENT, COURSE} è
                                             una chiave candidata

              FD1: {STUDENT, COURSE} -> INSTRUCTOR
                      FD2: INSTRUCTOR -> COURSE
                          È in 3NF ma non in BCNF
                        Esistono varie decomposizioni
           {STUDENT, INSTRUCTOR} e {STUDENT, COURSE}
           {COURSE, INSTRUCTOR} e {COURSE , STUDENT}
        {INSTRUCTOR, COURSE} e {INSTRUCTOR, STUDENT}
             Tutte e tre perdono la dipendenza funzionale FD1
\end{lstlisting}
\subsection*{Spiegazione estesa}
[AGGIUNTA] Questa spiegazione collega il contenuto testuale della slide al contesto del corso. La slide enfatizza i concetti mostrati nell'estratto e il loro ruolo nella progettazione di basi di dati relazionali, con particolare attenzione a correttezza semantica, qualità dello schema e compromessi progettuali.
[AGGIUNTA] Fonte integrazione: riformulazione didattica del testo presente nella slide 98 del file Lezione 01 - Dipendenze funzionali e Normalizzazione di DB Relazionali.pdf.
\subsection*{Derivazioni}
[AGGIUNTA] Nella slide compaiono espressioni simboliche/relazioni. Derivazione didattica generale:
\begin{align*}
X^+ &= X \\
X^+ &\leftarrow X^+ \cup Y \quad \text{se } (U \to V) \in F \text{ e } U \subseteq X^+ \text{ con } Y=V \\
\text{Stop} &\text{ quando non si aggiungono nuovi attributi in } X^+
\end{align*}
[AGGIUNTA] Questa procedura esplicita il metodo di chiusura usato nelle slide su dipendenze e chiavi.
\begin{itemize}[leftmargin=*]
\item Forma osservata nella slide: \texttt{FD1: (STUDENT, COURSE) -> INSTRUCTOR}
\item Forma osservata nella slide: \texttt{FD2: INSTRUCTOR -> COURSE}
\end{itemize}
\subsection*{Esempio numerico / esercizio}
[AGGIUNTA] Esempio: dato uno schema R(A,B,C) con dipendenze funzionali A->B e B->C, mostrare la chiusura A^+ e discutere se A è superchiave. Soluzione sintetica: A^+={A,B,C}, quindi A determina tutti gli attributi di R ed è superchiave.
\subsection*{Note su assunzioni o integrazioni}
[ASSUNZIONE] L'estrazione OCR può perdere formattazione o simboli della slide originale.
[AGGIUNTA] Le spiegazioni sono estensioni didattiche per rendere espliciti passaggi impliciti nel materiale proiettato.
\paragraph{Riferimento fonte} File: Lezione 01 - Dipendenze funzionali e Normalizzazione di DB Relazionali.pdf. Numero slide: 98.
\section{BASI DI DATI 2 --- Slide 1}
\subsection*{Estratto originale dalla slide 1}
\begin{lstlisting}
BASI DI DATI 2


ALGORITMI PER LA PROGETTAZIONE DI
DB RELAZIONALI E ULTERIORI
DIPENDENZE

                Prof.ssa G. Tortora
\end{lstlisting}
\subsection*{Spiegazione estesa}
[AGGIUNTA] Questa spiegazione collega il contenuto testuale della slide al contesto del corso. La slide enfatizza i concetti mostrati nell'estratto e il loro ruolo nella progettazione di basi di dati relazionali, con particolare attenzione a correttezza semantica, qualità dello schema e compromessi progettuali.
[AGGIUNTA] Fonte integrazione: riformulazione didattica del testo presente nella slide 1 del file Lezione 02 - Algoritmi per la progettazione di DB Relazionali e ulteriori dipendenze.pdf.
\subsection*{Derivazioni}
[AGGIUNTA] In questa slide non sono state rilevate equazioni esplicite nel testo estratto; non è quindi richiesta una derivazione algebrica puntuale.
\subsection*{Esempio numerico / esercizio}
[AGGIUNTA] Esempio: dato uno schema R(A,B,C) con dipendenze funzionali A->B e B->C, mostrare la chiusura A^+ e discutere se A è superchiave. Soluzione sintetica: A^+={A,B,C}, quindi A determina tutti gli attributi di R ed è superchiave.
\subsection*{Note su assunzioni o integrazioni}
[ASSUNZIONE] L'estrazione OCR può perdere formattazione o simboli della slide originale.
[AGGIUNTA] Le spiegazioni sono estensioni didattiche per rendere espliciti passaggi impliciti nel materiale proiettato.
\paragraph{Riferimento fonte} File: Lezione 02 - Algoritmi per la progettazione di DB Relazionali e ulteriori dipendenze.pdf. Numero slide: 1.
\section{Top-Down o Bottom-Up? --- Slide 2}
\subsection*{Estratto originale dalla slide 2}
\begin{lstlisting}
Top-Down o Bottom-Up?
-    Per la progettazione di un db relazionale,
    esistono due differenti filosofie:
      L'approccio Top-Down
        Si parte dal diagramma ER (o EER), si mappa in uno schema
        relazionale e si applicano le procedure di normalizzazione.
      L'approccio Bottom-Up
        È più formale, poiché considera lo schema del db solo in termini
        di dipendenze funzionali.
        Dopo aver definito tutte le FD, si applica un algoritmo di
        normalizzazione per sintetizzare gli schemi di relazione in 3NF o
        in BCNF.
-    Studiamo delle tecniche per il secondo tipo di
    approccio.
\end{lstlisting}
\subsection*{Spiegazione estesa}
[AGGIUNTA] Questa spiegazione collega il contenuto testuale della slide al contesto del corso. La slide enfatizza i concetti mostrati nell'estratto e il loro ruolo nella progettazione di basi di dati relazionali, con particolare attenzione a correttezza semantica, qualità dello schema e compromessi progettuali.
[AGGIUNTA] Fonte integrazione: riformulazione didattica del testo presente nella slide 2 del file Lezione 02 - Algoritmi per la progettazione di DB Relazionali e ulteriori dipendenze.pdf.
\subsection*{Derivazioni}
[AGGIUNTA] In questa slide non sono state rilevate equazioni esplicite nel testo estratto; non è quindi richiesta una derivazione algebrica puntuale.
\subsection*{Esempio numerico / esercizio}
[AGGIUNTA] Esempio: dato uno schema R(A,B,C) con dipendenze funzionali A->B e B->C, mostrare la chiusura A^+ e discutere se A è superchiave. Soluzione sintetica: A^+={A,B,C}, quindi A determina tutti gli attributi di R ed è superchiave.
\subsection*{Note su assunzioni o integrazioni}
[ASSUNZIONE] L'estrazione OCR può perdere formattazione o simboli della slide originale.
[AGGIUNTA] Le spiegazioni sono estensioni didattiche per rendere espliciti passaggi impliciti nel materiale proiettato.
\paragraph{Riferimento fonte} File: Lezione 02 - Algoritmi per la progettazione di DB Relazionali e ulteriori dipendenze.pdf. Numero slide: 2.
\section{Algoritmi per la progettazione di --- Slide 3}
\subsection*{Estratto originale dalla slide 3}
\begin{lstlisting}
Algoritmi per la progettazione di
    schemi di db relazionali
\end{lstlisting}
\subsection*{Spiegazione estesa}
[AGGIUNTA] Questa spiegazione collega il contenuto testuale della slide al contesto del corso. La slide enfatizza i concetti mostrati nell'estratto e il loro ruolo nella progettazione di basi di dati relazionali, con particolare attenzione a correttezza semantica, qualità dello schema e compromessi progettuali.
[AGGIUNTA] Fonte integrazione: riformulazione didattica del testo presente nella slide 3 del file Lezione 02 - Algoritmi per la progettazione di DB Relazionali e ulteriori dipendenze.pdf.
\subsection*{Derivazioni}
[AGGIUNTA] In questa slide non sono state rilevate equazioni esplicite nel testo estratto; non è quindi richiesta una derivazione algebrica puntuale.
\subsection*{Esempio numerico / esercizio}
[AGGIUNTA] Esempio: dato uno schema R(A,B,C) con dipendenze funzionali A->B e B->C, mostrare la chiusura A^+ e discutere se A è superchiave. Soluzione sintetica: A^+={A,B,C}, quindi A determina tutti gli attributi di R ed è superchiave.
\subsection*{Note su assunzioni o integrazioni}
[ASSUNZIONE] L'estrazione OCR può perdere formattazione o simboli della slide originale.
[AGGIUNTA] Le spiegazioni sono estensioni didattiche per rendere espliciti passaggi impliciti nel materiale proiettato.
\paragraph{Riferimento fonte} File: Lezione 02 - Algoritmi per la progettazione di DB Relazionali e ulteriori dipendenze.pdf. Numero slide: 3.
\section{Decomposizione delle relazioni --- Slide 4}
\subsection*{Estratto originale dalla slide 4}
\begin{lstlisting}
Decomposizione delle relazioni
-    Gli algoritmi che presentiamo partono da un
    singolo schema di relazione universale
    R = {A1, A2, ...,An} che include tutti gli attributi del
    database.

-    I progettisti specificano l'insieme F delle
    dipendenze funzionali che devono valere sugli
    attributi di R.

-    Usando le dipendenze funzionali, gli algoritmi
    decompongono lo schema di relazione R in un
    insieme di schemi D = {R1, R2, ..., Rn}
      D è detto decomposizione di R.
\end{lstlisting}
\subsection*{Spiegazione estesa}
[AGGIUNTA] Questa spiegazione collega il contenuto testuale della slide al contesto del corso. La slide enfatizza i concetti mostrati nell'estratto e il loro ruolo nella progettazione di basi di dati relazionali, con particolare attenzione a correttezza semantica, qualità dello schema e compromessi progettuali.
[AGGIUNTA] Fonte integrazione: riformulazione didattica del testo presente nella slide 4 del file Lezione 02 - Algoritmi per la progettazione di DB Relazionali e ulteriori dipendenze.pdf.
\subsection*{Derivazioni}
[AGGIUNTA] Nella slide compaiono espressioni simboliche/relazioni. Derivazione didattica generale:
\begin{align*}
X^+ &= X \\
X^+ &\leftarrow X^+ \cup Y \quad \text{se } (U \to V) \in F \text{ e } U \subseteq X^+ \text{ con } Y=V \\
\text{Stop} &\text{ quando non si aggiungono nuovi attributi in } X^+
\end{align*}
[AGGIUNTA] Questa procedura esplicita il metodo di chiusura usato nelle slide su dipendenze e chiavi.
\begin{itemize}[leftmargin=*]
\item Forma osservata nella slide: \texttt{R = (A1, A2, ...,An) che include tutti gli attributi del}
\item Forma osservata nella slide: \texttt{insieme di schemi D = (R1, R2, ..., Rn)}
\end{itemize}
\subsection*{Esempio numerico / esercizio}
[AGGIUNTA] Esempio: dato uno schema R(A,B,C) con dipendenze funzionali A->B e B->C, mostrare la chiusura A^+ e discutere se A è superchiave. Soluzione sintetica: A^+={A,B,C}, quindi A determina tutti gli attributi di R ed è superchiave.
\subsection*{Note su assunzioni o integrazioni}
[ASSUNZIONE] L'estrazione OCR può perdere formattazione o simboli della slide originale.
[AGGIUNTA] Le spiegazioni sono estensioni didattiche per rendere espliciti passaggi impliciti nel materiale proiettato.
\paragraph{Riferimento fonte} File: Lezione 02 - Algoritmi per la progettazione di DB Relazionali e ulteriori dipendenze.pdf. Numero slide: 4.
\section{Conservazione degli attributi --- Slide 5}
\subsection*{Estratto originale dalla slide 5}
\begin{lstlisting}
Conservazione degli attributi
-    Ciascun attributo di R deve apparire in almeno
    una relazione in D:
                    n

                   UR =R
                   i=1
                         i



-    Questa condizione è detta di conservazione
    degli attributi.
\end{lstlisting}
\subsection*{Spiegazione estesa}
[AGGIUNTA] Questa spiegazione collega il contenuto testuale della slide al contesto del corso. La slide enfatizza i concetti mostrati nell'estratto e il loro ruolo nella progettazione di basi di dati relazionali, con particolare attenzione a correttezza semantica, qualità dello schema e compromessi progettuali.
[AGGIUNTA] Fonte integrazione: riformulazione didattica del testo presente nella slide 5 del file Lezione 02 - Algoritmi per la progettazione di DB Relazionali e ulteriori dipendenze.pdf.
\subsection*{Derivazioni}
[AGGIUNTA] Nella slide compaiono espressioni simboliche/relazioni. Derivazione didattica generale:
\begin{align*}
X^+ &= X \\
X^+ &\leftarrow X^+ \cup Y \quad \text{se } (U \to V) \in F \text{ e } U \subseteq X^+ \text{ con } Y=V \\
\text{Stop} &\text{ quando non si aggiungono nuovi attributi in } X^+
\end{align*}
[AGGIUNTA] Questa procedura esplicita il metodo di chiusura usato nelle slide su dipendenze e chiavi.
\begin{itemize}[leftmargin=*]
\item Forma osservata nella slide: \texttt{UR =R}
\item Forma osservata nella slide: \texttt{i=1}
\end{itemize}
\subsection*{Esempio numerico / esercizio}
[AGGIUNTA] Esempio: dato uno schema R(A,B,C) con dipendenze funzionali A->B e B->C, mostrare la chiusura A^+ e discutere se A è superchiave. Soluzione sintetica: A^+={A,B,C}, quindi A determina tutti gli attributi di R ed è superchiave.
\subsection*{Note su assunzioni o integrazioni}
[ASSUNZIONE] L'estrazione OCR può perdere formattazione o simboli della slide originale.
[AGGIUNTA] Le spiegazioni sono estensioni didattiche per rendere espliciti passaggi impliciti nel materiale proiettato.
\paragraph{Riferimento fonte} File: Lezione 02 - Algoritmi per la progettazione di DB Relazionali e ulteriori dipendenze.pdf. Numero slide: 5.
\section{Decomposizione delle relazioni (2) --- Slide 6}
\subsection*{Estratto originale dalla slide 6}
\begin{lstlisting}
Decomposizione delle relazioni (2)
-    Altro obiettivo è che ogni relazione individuale R i
    in D deve essere in BCNF (o in 3NF).
      Da notare che qualsiasi schema di relazione con soli
      due attributi è automaticamente in BCNF.


-    Questa condizione non è sufficiente per garantire
    di avere un buon db, in quanto il fenomeno delle
    tuple spurie può comunque presentarsi anche
    con relazioni in BCNF.
\end{lstlisting}
\subsection*{Spiegazione estesa}
[AGGIUNTA] Questa spiegazione collega il contenuto testuale della slide al contesto del corso. La slide enfatizza i concetti mostrati nell'estratto e il loro ruolo nella progettazione di basi di dati relazionali, con particolare attenzione a correttezza semantica, qualità dello schema e compromessi progettuali.
[AGGIUNTA] Fonte integrazione: riformulazione didattica del testo presente nella slide 6 del file Lezione 02 - Algoritmi per la progettazione di DB Relazionali e ulteriori dipendenze.pdf.
\subsection*{Derivazioni}
[AGGIUNTA] In questa slide non sono state rilevate equazioni esplicite nel testo estratto; non è quindi richiesta una derivazione algebrica puntuale.
\subsection*{Esempio numerico / esercizio}
[AGGIUNTA] Esempio: dato uno schema R(A,B,C) con dipendenze funzionali A->B e B->C, mostrare la chiusura A^+ e discutere se A è superchiave. Soluzione sintetica: A^+={A,B,C}, quindi A determina tutti gli attributi di R ed è superchiave.
\subsection*{Note su assunzioni o integrazioni}
[ASSUNZIONE] L'estrazione OCR può perdere formattazione o simboli della slide originale.
[AGGIUNTA] Le spiegazioni sono estensioni didattiche per rendere espliciti passaggi impliciti nel materiale proiettato.
\paragraph{Riferimento fonte} File: Lezione 02 - Algoritmi per la progettazione di DB Relazionali e ulteriori dipendenze.pdf. Numero slide: 6.
\section{Conservazione delle dipendenze --- Slide 7}
\subsection*{Estratto originale dalla slide 7}
\begin{lstlisting}
Conservazione delle dipendenze
-    Sarebbe utile che le dipendenze funzionali
    X -> Y specificate in F, apparissero direttamente
    in una delle Ri della decomposizione D o possano
    essere inferite da dipendenze in altre R i.
      Questa è (informalmente) la condizione di
     conservazione delle dipendenze.

-    È fondamentale non perdere delle dipendenze,
    poiché queste rappresentano dei vincoli sul
    database.
\end{lstlisting}
\subsection*{Spiegazione estesa}
[AGGIUNTA] Questa spiegazione collega il contenuto testuale della slide al contesto del corso. La slide enfatizza i concetti mostrati nell'estratto e il loro ruolo nella progettazione di basi di dati relazionali, con particolare attenzione a correttezza semantica, qualità dello schema e compromessi progettuali.
[AGGIUNTA] Fonte integrazione: riformulazione didattica del testo presente nella slide 7 del file Lezione 02 - Algoritmi per la progettazione di DB Relazionali e ulteriori dipendenze.pdf.
\subsection*{Derivazioni}
[AGGIUNTA] Nella slide compaiono espressioni simboliche/relazioni. Derivazione didattica generale:
\begin{align*}
X^+ &= X \\
X^+ &\leftarrow X^+ \cup Y \quad \text{se } (U \to V) \in F \text{ e } U \subseteq X^+ \text{ con } Y=V \\
\text{Stop} &\text{ quando non si aggiungono nuovi attributi in } X^+
\end{align*}
[AGGIUNTA] Questa procedura esplicita il metodo di chiusura usato nelle slide su dipendenze e chiavi.
\begin{itemize}[leftmargin=*]
\item Forma osservata nella slide: \texttt{X -> Y specificate in F, apparissero direttamente}
\end{itemize}
\subsection*{Esempio numerico / esercizio}
[AGGIUNTA] Esempio: dato uno schema R(A,B,C) con dipendenze funzionali A->B e B->C, mostrare la chiusura A^+ e discutere se A è superchiave. Soluzione sintetica: A^+={A,B,C}, quindi A determina tutti gli attributi di R ed è superchiave.
\subsection*{Note su assunzioni o integrazioni}
[ASSUNZIONE] L'estrazione OCR può perdere formattazione o simboli della slide originale.
[AGGIUNTA] Le spiegazioni sono estensioni didattiche per rendere espliciti passaggi impliciti nel materiale proiettato.
\paragraph{Riferimento fonte} File: Lezione 02 - Algoritmi per la progettazione di DB Relazionali e ulteriori dipendenze.pdf. Numero slide: 7.
\section{Conservazione delle dipendenze (2) --- Slide 8}
\subsection*{Estratto originale dalla slide 8}
\begin{lstlisting}
Conservazione delle dipendenze (2)
-    Non è necessario che esattamente le stesse
    dipendenze specificate in F appaiono nelle
    relazioni della decomposizione D:
      è sufficiente che l'unione delle dipendenze che
      valgono nelle singole relazioni di D sia equivalente
      ad F.

-    Per formalizzare tali concetti abbiamo bisogno di
    alcune definizioni ...
\end{lstlisting}
\subsection*{Spiegazione estesa}
[AGGIUNTA] Questa spiegazione collega il contenuto testuale della slide al contesto del corso. La slide enfatizza i concetti mostrati nell'estratto e il loro ruolo nella progettazione di basi di dati relazionali, con particolare attenzione a correttezza semantica, qualità dello schema e compromessi progettuali.
[AGGIUNTA] Fonte integrazione: riformulazione didattica del testo presente nella slide 8 del file Lezione 02 - Algoritmi per la progettazione di DB Relazionali e ulteriori dipendenze.pdf.
\subsection*{Derivazioni}
[AGGIUNTA] In questa slide non sono state rilevate equazioni esplicite nel testo estratto; non è quindi richiesta una derivazione algebrica puntuale.
\subsection*{Esempio numerico / esercizio}
[AGGIUNTA] Esempio: dato uno schema R(A,B,C) con dipendenze funzionali A->B e B->C, mostrare la chiusura A^+ e discutere se A è superchiave. Soluzione sintetica: A^+={A,B,C}, quindi A determina tutti gli attributi di R ed è superchiave.
\subsection*{Note su assunzioni o integrazioni}
[ASSUNZIONE] L'estrazione OCR può perdere formattazione o simboli della slide originale.
[AGGIUNTA] Le spiegazioni sono estensioni didattiche per rendere espliciti passaggi impliciti nel materiale proiettato.
\paragraph{Riferimento fonte} File: Lezione 02 - Algoritmi per la progettazione di DB Relazionali e ulteriori dipendenze.pdf. Numero slide: 8.
\section{Proiezione --- Slide 9}
\subsection*{Estratto originale dalla slide 9}
\begin{lstlisting}
Proiezione
-    Definizione:
    Dato un insieme di dipendenze F in R, la
    proiezione di F su Ri , denotata da     (F)
                                           Ri

    (dove Ri è un sottoinsieme di R), è l'insieme di
    dipendenze X -> Y in F+ tale che gli attributi in
    X   Y sono tutti contenuti in Ri .
\end{lstlisting}
\subsection*{Spiegazione estesa}
[AGGIUNTA] Questa spiegazione collega il contenuto testuale della slide al contesto del corso. La slide enfatizza i concetti mostrati nell'estratto e il loro ruolo nella progettazione di basi di dati relazionali, con particolare attenzione a correttezza semantica, qualità dello schema e compromessi progettuali.
[AGGIUNTA] Fonte integrazione: riformulazione didattica del testo presente nella slide 9 del file Lezione 02 - Algoritmi per la progettazione di DB Relazionali e ulteriori dipendenze.pdf.
\subsection*{Derivazioni}
[AGGIUNTA] Nella slide compaiono espressioni simboliche/relazioni. Derivazione didattica generale:
\begin{align*}
X^+ &= X \\
X^+ &\leftarrow X^+ \cup Y \quad \text{se } (U \to V) \in F \text{ e } U \subseteq X^+ \text{ con } Y=V \\
\text{Stop} &\text{ quando non si aggiungono nuovi attributi in } X^+
\end{align*}
[AGGIUNTA] Questa procedura esplicita il metodo di chiusura usato nelle slide su dipendenze e chiavi.
\begin{itemize}[leftmargin=*]
\item Forma osservata nella slide: \texttt{dipendenze X -> Y in F+ tale che gli attributi in}
\end{itemize}
\subsection*{Esempio numerico / esercizio}
[AGGIUNTA] Esempio: dato uno schema R(A,B,C) con dipendenze funzionali A->B e B->C, mostrare la chiusura A^+ e discutere se A è superchiave. Soluzione sintetica: A^+={A,B,C}, quindi A determina tutti gli attributi di R ed è superchiave.
\subsection*{Note su assunzioni o integrazioni}
[ASSUNZIONE] L'estrazione OCR può perdere formattazione o simboli della slide originale.
[AGGIUNTA] Le spiegazioni sono estensioni didattiche per rendere espliciti passaggi impliciti nel materiale proiettato.
\paragraph{Riferimento fonte} File: Lezione 02 - Algoritmi per la progettazione di DB Relazionali e ulteriori dipendenze.pdf. Numero slide: 9.
\section{Decomposizione dependency-preserving --- Slide 10}
\subsection*{Estratto originale dalla slide 10}
\begin{lstlisting}
Decomposizione dependency-preserving

-    Diciamo che una decomposizione
    D = {R1, R2, ..., Rm} di R è
    dependency-preserving rispetto a F se l'unione
    delle proiezioni di F su ciascun R i in D è
    equivalente a F, cioè:

       ((    R1 (F))   ...   (       Rm (F)) )+ = F+




-    Se una decomposizione non è dependency-
    preserving, si ha la perdita di qualche dipendenza
    nella decomposizione.
\end{lstlisting}
\subsection*{Spiegazione estesa}
[AGGIUNTA] Questa spiegazione collega il contenuto testuale della slide al contesto del corso. La slide enfatizza i concetti mostrati nell'estratto e il loro ruolo nella progettazione di basi di dati relazionali, con particolare attenzione a correttezza semantica, qualità dello schema e compromessi progettuali.
[AGGIUNTA] Fonte integrazione: riformulazione didattica del testo presente nella slide 10 del file Lezione 02 - Algoritmi per la progettazione di DB Relazionali e ulteriori dipendenze.pdf.
\subsection*{Derivazioni}
[AGGIUNTA] Nella slide compaiono espressioni simboliche/relazioni. Derivazione didattica generale:
\begin{align*}
X^+ &= X \\
X^+ &\leftarrow X^+ \cup Y \quad \text{se } (U \to V) \in F \text{ e } U \subseteq X^+ \text{ con } Y=V \\
\text{Stop} &\text{ quando non si aggiungono nuovi attributi in } X^+
\end{align*}
[AGGIUNTA] Questa procedura esplicita il metodo di chiusura usato nelle slide su dipendenze e chiavi.
\begin{itemize}[leftmargin=*]
\item Forma osservata nella slide: \texttt{D = (R1, R2, ..., Rm) di R è}
\item Forma osservata nella slide: \texttt{((    R1 (F))   ...   (       Rm (F)) )+ = F+}
\end{itemize}
\subsection*{Esempio numerico / esercizio}
[AGGIUNTA] Esempio: dato uno schema R(A,B,C) con dipendenze funzionali A->B e B->C, mostrare la chiusura A^+ e discutere se A è superchiave. Soluzione sintetica: A^+={A,B,C}, quindi A determina tutti gli attributi di R ed è superchiave.
\subsection*{Note su assunzioni o integrazioni}
[ASSUNZIONE] L'estrazione OCR può perdere formattazione o simboli della slide originale.
[AGGIUNTA] Le spiegazioni sono estensioni didattiche per rendere espliciti passaggi impliciti nel materiale proiettato.
\paragraph{Riferimento fonte} File: Lezione 02 - Algoritmi per la progettazione di DB Relazionali e ulteriori dipendenze.pdf. Numero slide: 10.
\section{Decomposizione dependency-preserving: --- Slide 11}
\subsection*{Estratto originale dalla slide 11}
\begin{lstlisting}
Decomposizione dependency-preserving:
Esempio
-    Decomposizione non dependency-preserving:




                     La FD2 viene persa
\end{lstlisting}
\subsection*{Spiegazione estesa}
[AGGIUNTA] Questa spiegazione collega il contenuto testuale della slide al contesto del corso. La slide enfatizza i concetti mostrati nell'estratto e il loro ruolo nella progettazione di basi di dati relazionali, con particolare attenzione a correttezza semantica, qualità dello schema e compromessi progettuali.
[AGGIUNTA] Fonte integrazione: riformulazione didattica del testo presente nella slide 11 del file Lezione 02 - Algoritmi per la progettazione di DB Relazionali e ulteriori dipendenze.pdf.
\subsection*{Derivazioni}
[AGGIUNTA] In questa slide non sono state rilevate equazioni esplicite nel testo estratto; non è quindi richiesta una derivazione algebrica puntuale.
\subsection*{Esempio numerico / esercizio}
[AGGIUNTA] Esempio: dato uno schema R(A,B,C) con dipendenze funzionali A->B e B->C, mostrare la chiusura A^+ e discutere se A è superchiave. Soluzione sintetica: A^+={A,B,C}, quindi A determina tutti gli attributi di R ed è superchiave.
\subsection*{Note su assunzioni o integrazioni}
[ASSUNZIONE] L'estrazione OCR può perdere formattazione o simboli della slide originale.
[AGGIUNTA] Le spiegazioni sono estensioni didattiche per rendere espliciti passaggi impliciti nel materiale proiettato.
\paragraph{Riferimento fonte} File: Lezione 02 - Algoritmi per la progettazione di DB Relazionali e ulteriori dipendenze.pdf. Numero slide: 11.
\section{Decomposizione dependency-preserving: --- Slide 12}
\subsection*{Estratto originale dalla slide 12}
\begin{lstlisting}
Decomposizione dependency-preserving:
Esempio 2
-    Decomposizione dependency-preserving:
\end{lstlisting}
\subsection*{Spiegazione estesa}
[AGGIUNTA] Questa spiegazione collega il contenuto testuale della slide al contesto del corso. La slide enfatizza i concetti mostrati nell'estratto e il loro ruolo nella progettazione di basi di dati relazionali, con particolare attenzione a correttezza semantica, qualità dello schema e compromessi progettuali.
[AGGIUNTA] Fonte integrazione: riformulazione didattica del testo presente nella slide 12 del file Lezione 02 - Algoritmi per la progettazione di DB Relazionali e ulteriori dipendenze.pdf.
\subsection*{Derivazioni}
[AGGIUNTA] In questa slide non sono state rilevate equazioni esplicite nel testo estratto; non è quindi richiesta una derivazione algebrica puntuale.
\subsection*{Esempio numerico / esercizio}
[AGGIUNTA] Esempio: dato uno schema R(A,B,C) con dipendenze funzionali A->B e B->C, mostrare la chiusura A^+ e discutere se A è superchiave. Soluzione sintetica: A^+={A,B,C}, quindi A determina tutti gli attributi di R ed è superchiave.
\subsection*{Note su assunzioni o integrazioni}
[ASSUNZIONE] L'estrazione OCR può perdere formattazione o simboli della slide originale.
[AGGIUNTA] Le spiegazioni sono estensioni didattiche per rendere espliciti passaggi impliciti nel materiale proiettato.
\paragraph{Riferimento fonte} File: Lezione 02 - Algoritmi per la progettazione di DB Relazionali e ulteriori dipendenze.pdf. Numero slide: 12.
\section{Proposizione 1 --- Slide 13}
\subsection*{Estratto originale dalla slide 13}
\begin{lstlisting}
Proposizione 1
-    È sempre possibile trovare una decomposizione
    dependency-preserving D rispetto a F tale che
    ogni Ri in D è in 3NF.
\end{lstlisting}
\subsection*{Spiegazione estesa}
[AGGIUNTA] Questa spiegazione collega il contenuto testuale della slide al contesto del corso. La slide enfatizza i concetti mostrati nell'estratto e il loro ruolo nella progettazione di basi di dati relazionali, con particolare attenzione a correttezza semantica, qualità dello schema e compromessi progettuali.
[AGGIUNTA] Fonte integrazione: riformulazione didattica del testo presente nella slide 13 del file Lezione 02 - Algoritmi per la progettazione di DB Relazionali e ulteriori dipendenze.pdf.
\subsection*{Derivazioni}
[AGGIUNTA] In questa slide non sono state rilevate equazioni esplicite nel testo estratto; non è quindi richiesta una derivazione algebrica puntuale.
\subsection*{Esempio numerico / esercizio}
[AGGIUNTA] Esempio: dato uno schema R(A,B,C) con dipendenze funzionali A->B e B->C, mostrare la chiusura A^+ e discutere se A è superchiave. Soluzione sintetica: A^+={A,B,C}, quindi A determina tutti gli attributi di R ed è superchiave.
\subsection*{Note su assunzioni o integrazioni}
[ASSUNZIONE] L'estrazione OCR può perdere formattazione o simboli della slide originale.
[AGGIUNTA] Le spiegazioni sono estensioni didattiche per rendere espliciti passaggi impliciti nel materiale proiettato.
\paragraph{Riferimento fonte} File: Lezione 02 - Algoritmi per la progettazione di DB Relazionali e ulteriori dipendenze.pdf. Numero slide: 13.
\section{Copertura minimale --- Slide 14}
\subsection*{Estratto originale dalla slide 14}
\begin{lstlisting}
Copertura minimale
-    Un insieme di dipendenze funzionali F è minimale
    se:
      Ogni dipendenza in F ha un singolo attributo sul lato
      destro.
      Non è possibile rimpiazzare una dipendenza
      X -> A con una dipendenza Y -> A dove Y è un
      sottoinsieme proprio di X ed avere ancora un insieme di
      dipendenze che è equivalente ad F.
      Non possiamo rimuovere una dipendenza da F ed
      ottenere un insieme di dipendenze equivalente ad F.
-    Definizione: Una copertura minimale di un
    insieme E di dipendenze funzionali è un insieme
    minimale F di dipendenze che è equivalente ad
    E.
\end{lstlisting}
\subsection*{Spiegazione estesa}
[AGGIUNTA] Questa spiegazione collega il contenuto testuale della slide al contesto del corso. La slide enfatizza i concetti mostrati nell'estratto e il loro ruolo nella progettazione di basi di dati relazionali, con particolare attenzione a correttezza semantica, qualità dello schema e compromessi progettuali.
[AGGIUNTA] Fonte integrazione: riformulazione didattica del testo presente nella slide 14 del file Lezione 02 - Algoritmi per la progettazione di DB Relazionali e ulteriori dipendenze.pdf.
\subsection*{Derivazioni}
[AGGIUNTA] Nella slide compaiono espressioni simboliche/relazioni. Derivazione didattica generale:
\begin{align*}
X^+ &= X \\
X^+ &\leftarrow X^+ \cup Y \quad \text{se } (U \to V) \in F \text{ e } U \subseteq X^+ \text{ con } Y=V \\
\text{Stop} &\text{ quando non si aggiungono nuovi attributi in } X^+
\end{align*}
[AGGIUNTA] Questa procedura esplicita il metodo di chiusura usato nelle slide su dipendenze e chiavi.
\begin{itemize}[leftmargin=*]
\item Forma osservata nella slide: \texttt{X -> A con una dipendenza Y -> A dove Y è un}
\end{itemize}
\subsection*{Esempio numerico / esercizio}
[AGGIUNTA] Esempio: dato uno schema R(A,B,C) con dipendenze funzionali A->B e B->C, mostrare la chiusura A^+ e discutere se A è superchiave. Soluzione sintetica: A^+={A,B,C}, quindi A determina tutti gli attributi di R ed è superchiave.
\subsection*{Note su assunzioni o integrazioni}
[ASSUNZIONE] L'estrazione OCR può perdere formattazione o simboli della slide originale.
[AGGIUNTA] Le spiegazioni sono estensioni didattiche per rendere espliciti passaggi impliciti nel materiale proiettato.
\paragraph{Riferimento fonte} File: Lezione 02 - Algoritmi per la progettazione di DB Relazionali e ulteriori dipendenze.pdf. Numero slide: 14.
\section{Algoritmo di sintesi relazionale --- Slide 15}
\subsection*{Estratto originale dalla slide 15}
\begin{lstlisting}
Algoritmo di sintesi relazionale
-    Algoritmo 1
    Decomposizione dependency-preserving in
    schemi di relazione 3NF:
    1.   trovare una copertura minimale G di F.
    2.   per ogni parte sinistra X di una dipendenza
         funzionale che appare in G, creare uno schema di
         relazione {X   A1   A2   ...   Am} in D dove
         X -> A1, X -> A2, ... ,X -> Am sono le sole dipendenze
         in G aventi X come parte sinistra.
    3.   mettere in uno schema di relazione singolo tutti gli
         attributi rimanenti, per garantire la proprietà di
         attribute-preservation.
\end{lstlisting}
\subsection*{Spiegazione estesa}
[AGGIUNTA] Questa spiegazione collega il contenuto testuale della slide al contesto del corso. La slide enfatizza i concetti mostrati nell'estratto e il loro ruolo nella progettazione di basi di dati relazionali, con particolare attenzione a correttezza semantica, qualità dello schema e compromessi progettuali.
[AGGIUNTA] Fonte integrazione: riformulazione didattica del testo presente nella slide 15 del file Lezione 02 - Algoritmi per la progettazione di DB Relazionali e ulteriori dipendenze.pdf.
\subsection*{Derivazioni}
[AGGIUNTA] Nella slide compaiono espressioni simboliche/relazioni. Derivazione didattica generale:
\begin{align*}
X^+ &= X \\
X^+ &\leftarrow X^+ \cup Y \quad \text{se } (U \to V) \in F \text{ e } U \subseteq X^+ \text{ con } Y=V \\
\text{Stop} &\text{ quando non si aggiungono nuovi attributi in } X^+
\end{align*}
[AGGIUNTA] Questa procedura esplicita il metodo di chiusura usato nelle slide su dipendenze e chiavi.
\begin{itemize}[leftmargin=*]
\item Forma osservata nella slide: \texttt{X -> A1, X -> A2, ... ,X -> Am sono le sole dipendenze}
\end{itemize}
\subsection*{Esempio numerico / esercizio}
[AGGIUNTA] Esempio: dato uno schema R(A,B,C) con dipendenze funzionali A->B e B->C, mostrare la chiusura A^+ e discutere se A è superchiave. Soluzione sintetica: A^+={A,B,C}, quindi A determina tutti gli attributi di R ed è superchiave.
\subsection*{Note su assunzioni o integrazioni}
[ASSUNZIONE] L'estrazione OCR può perdere formattazione o simboli della slide originale.
[AGGIUNTA] Le spiegazioni sono estensioni didattiche per rendere espliciti passaggi impliciti nel materiale proiettato.
\paragraph{Riferimento fonte} File: Lezione 02 - Algoritmi per la progettazione di DB Relazionali e ulteriori dipendenze.pdf. Numero slide: 15.
\section{Algoritmo per trovare una copertura --- Slide 16}
\subsection*{Estratto originale dalla slide 16}
\begin{lstlisting}
Algoritmo per trovare una copertura
minimale G per F
-    Il passo 1 dell'algoritmo richiede di trovare una
    copertura minimale.
    Ecco l'algoritmo per fare ciò:
-    Algoritmo 1.a
    Consta di quattro passi.
    G è l'insieme minimale di dipendenze funzionali
    equivalente ad F.

    1.   porre G:= F;
    2.   rimpiazzare ogni dipendenza funzionale
         X-> {A1, A2, ..., An} in G, con n dipendenze funzionali
         X -> A1, X -> A2, ... , X -> An ;
    ...
\end{lstlisting}
\subsection*{Spiegazione estesa}
[AGGIUNTA] Questa spiegazione collega il contenuto testuale della slide al contesto del corso. La slide enfatizza i concetti mostrati nell'estratto e il loro ruolo nella progettazione di basi di dati relazionali, con particolare attenzione a correttezza semantica, qualità dello schema e compromessi progettuali.
[AGGIUNTA] Fonte integrazione: riformulazione didattica del testo presente nella slide 16 del file Lezione 02 - Algoritmi per la progettazione di DB Relazionali e ulteriori dipendenze.pdf.
\subsection*{Derivazioni}
[AGGIUNTA] Nella slide compaiono espressioni simboliche/relazioni. Derivazione didattica generale:
\begin{align*}
X^+ &= X \\
X^+ &\leftarrow X^+ \cup Y \quad \text{se } (U \to V) \in F \text{ e } U \subseteq X^+ \text{ con } Y=V \\
\text{Stop} &\text{ quando non si aggiungono nuovi attributi in } X^+
\end{align*}
[AGGIUNTA] Questa procedura esplicita il metodo di chiusura usato nelle slide su dipendenze e chiavi.
\begin{itemize}[leftmargin=*]
\item Forma osservata nella slide: \texttt{1.   porre G:= F;}
\item Forma osservata nella slide: \texttt{X-> (A1, A2, ..., An) in G, con n dipendenze funzionali}
\item Forma osservata nella slide: \texttt{X -> A1, X -> A2, ... , X -> An ;}
\end{itemize}
\subsection*{Esempio numerico / esercizio}
[AGGIUNTA] Esempio: dato uno schema R(A,B,C) con dipendenze funzionali A->B e B->C, mostrare la chiusura A^+ e discutere se A è superchiave. Soluzione sintetica: A^+={A,B,C}, quindi A determina tutti gli attributi di R ed è superchiave.
\subsection*{Note su assunzioni o integrazioni}
[ASSUNZIONE] L'estrazione OCR può perdere formattazione o simboli della slide originale.
[AGGIUNTA] Le spiegazioni sono estensioni didattiche per rendere espliciti passaggi impliciti nel materiale proiettato.
\paragraph{Riferimento fonte} File: Lezione 02 - Algoritmi per la progettazione di DB Relazionali e ulteriori dipendenze.pdf. Numero slide: 16.
\section{Algoritmo per trovare una copertura --- Slide 17}
\subsection*{Estratto originale dalla slide 17}
\begin{lstlisting}
Algoritmo per trovare una copertura
minimale G per F (2)
  3.   per ogni dipendenza funzionale X -> A in G
        per ogni attributo B che è un elemento di X
        {
          se { {G - {X -> A} }   { (X - {B}) -> A} } è
              equivalente a G
          allora sostituire X -> A con
              (X - {B}) -> A in G;
         }
  4.   per ogni dipendenza funzionale rimanente X -> A in G
         {
            se { G - {X -> A} } è equivalente a G
            allora rimuovere X -> A da G;
         }
\end{lstlisting}
\subsection*{Spiegazione estesa}
[AGGIUNTA] Questa spiegazione collega il contenuto testuale della slide al contesto del corso. La slide enfatizza i concetti mostrati nell'estratto e il loro ruolo nella progettazione di basi di dati relazionali, con particolare attenzione a correttezza semantica, qualità dello schema e compromessi progettuali.
[AGGIUNTA] Fonte integrazione: riformulazione didattica del testo presente nella slide 17 del file Lezione 02 - Algoritmi per la progettazione di DB Relazionali e ulteriori dipendenze.pdf.
\subsection*{Derivazioni}
[AGGIUNTA] Nella slide compaiono espressioni simboliche/relazioni. Derivazione didattica generale:
\begin{align*}
X^+ &= X \\
X^+ &\leftarrow X^+ \cup Y \quad \text{se } (U \to V) \in F \text{ e } U \subseteq X^+ \text{ con } Y=V \\
\text{Stop} &\text{ quando non si aggiungono nuovi attributi in } X^+
\end{align*}
[AGGIUNTA] Questa procedura esplicita il metodo di chiusura usato nelle slide su dipendenze e chiavi.
\begin{itemize}[leftmargin=*]
\item Forma osservata nella slide: \texttt{3.   per ogni dipendenza funzionale X -> A in G}
\item Forma osservata nella slide: \texttt{se ( (G - (X -> A) )   ( (X - (B)) -> A) ) è}
\item Forma osservata nella slide: \texttt{allora sostituire X -> A con}
\end{itemize}
\subsection*{Esempio numerico / esercizio}
[AGGIUNTA] Esempio: dato uno schema R(A,B,C) con dipendenze funzionali A->B e B->C, mostrare la chiusura A^+ e discutere se A è superchiave. Soluzione sintetica: A^+={A,B,C}, quindi A determina tutti gli attributi di R ed è superchiave.
\subsection*{Note su assunzioni o integrazioni}
[ASSUNZIONE] L'estrazione OCR può perdere formattazione o simboli della slide originale.
[AGGIUNTA] Le spiegazioni sono estensioni didattiche per rendere espliciti passaggi impliciti nel materiale proiettato.
\paragraph{Riferimento fonte} File: Lezione 02 - Algoritmi per la progettazione di DB Relazionali e ulteriori dipendenze.pdf. Numero slide: 17.
\section{Decomposizione e Join Lossless --- Slide 18}
\subsection*{Estratto originale dalla slide 18}
\begin{lstlisting}
Decomposizione e Join Lossless
(o Non-additive)
-    La lossless (o non-additive) join (LJ) è un'altra
    proprietà di cui dovrebbe godere una
    decomposizione.

-    Una decomposizione con LJ garantisce che non
    vengono generate tuple spurie applicando
    un'operazione di NATURAL-JOIN alle relazioni
    nella decomposizione.
\end{lstlisting}
\subsection*{Spiegazione estesa}
[AGGIUNTA] Questa spiegazione collega il contenuto testuale della slide al contesto del corso. La slide enfatizza i concetti mostrati nell'estratto e il loro ruolo nella progettazione di basi di dati relazionali, con particolare attenzione a correttezza semantica, qualità dello schema e compromessi progettuali.
[AGGIUNTA] Fonte integrazione: riformulazione didattica del testo presente nella slide 18 del file Lezione 02 - Algoritmi per la progettazione di DB Relazionali e ulteriori dipendenze.pdf.
\subsection*{Derivazioni}
[AGGIUNTA] In questa slide non sono state rilevate equazioni esplicite nel testo estratto; non è quindi richiesta una derivazione algebrica puntuale.
\subsection*{Esempio numerico / esercizio}
[AGGIUNTA] Esempio: dato uno schema R(A,B,C) con dipendenze funzionali A->B e B->C, mostrare la chiusura A^+ e discutere se A è superchiave. Soluzione sintetica: A^+={A,B,C}, quindi A determina tutti gli attributi di R ed è superchiave.
\subsection*{Note su assunzioni o integrazioni}
[ASSUNZIONE] L'estrazione OCR può perdere formattazione o simboli della slide originale.
[AGGIUNTA] Le spiegazioni sono estensioni didattiche per rendere espliciti passaggi impliciti nel materiale proiettato.
\paragraph{Riferimento fonte} File: Lezione 02 - Algoritmi per la progettazione di DB Relazionali e ulteriori dipendenze.pdf. Numero slide: 18.
\section{Lossless Join --- Slide 19}
\subsection*{Estratto originale dalla slide 19}
\begin{lstlisting}
Lossless Join
-    Definizione:
    Una decomposizione D = {R1, R2, ..., Rm} di R ha
    la proprietà di lossless join rispetto all'insieme di
    dipendenze di F su R, se per ogni stato di
    relazione r di R che soddisfa F, vale:
        *(  R1(r), ...,   Rm(r) ) = r

    dove con * indichiamo la natural-join.
\end{lstlisting}
\subsection*{Spiegazione estesa}
[AGGIUNTA] Questa spiegazione collega il contenuto testuale della slide al contesto del corso. La slide enfatizza i concetti mostrati nell'estratto e il loro ruolo nella progettazione di basi di dati relazionali, con particolare attenzione a correttezza semantica, qualità dello schema e compromessi progettuali.
[AGGIUNTA] Fonte integrazione: riformulazione didattica del testo presente nella slide 19 del file Lezione 02 - Algoritmi per la progettazione di DB Relazionali e ulteriori dipendenze.pdf.
\subsection*{Derivazioni}
[AGGIUNTA] Nella slide compaiono espressioni simboliche/relazioni. Derivazione didattica generale:
\begin{align*}
X^+ &= X \\
X^+ &\leftarrow X^+ \cup Y \quad \text{se } (U \to V) \in F \text{ e } U \subseteq X^+ \text{ con } Y=V \\
\text{Stop} &\text{ quando non si aggiungono nuovi attributi in } X^+
\end{align*}
[AGGIUNTA] Questa procedura esplicita il metodo di chiusura usato nelle slide su dipendenze e chiavi.
\begin{itemize}[leftmargin=*]
\item Forma osservata nella slide: \texttt{Una decomposizione D = (R1, R2, ..., Rm) di R ha}
\item Forma osservata nella slide: \texttt{*(  R1(r), ...,   Rm(r) ) = r}
\end{itemize}
\subsection*{Esempio numerico / esercizio}
[AGGIUNTA] Esempio: dato uno schema R(A,B,C) con dipendenze funzionali A->B e B->C, mostrare la chiusura A^+ e discutere se A è superchiave. Soluzione sintetica: A^+={A,B,C}, quindi A determina tutti gli attributi di R ed è superchiave.
\subsection*{Note su assunzioni o integrazioni}
[ASSUNZIONE] L'estrazione OCR può perdere formattazione o simboli della slide originale.
[AGGIUNTA] Le spiegazioni sono estensioni didattiche per rendere espliciti passaggi impliciti nel materiale proiettato.
\paragraph{Riferimento fonte} File: Lezione 02 - Algoritmi per la progettazione di DB Relazionali e ulteriori dipendenze.pdf. Numero slide: 19.
\section{Test della proprietà di lossless join --- Slide 20}
\subsection*{Estratto originale dalla slide 20}
\begin{lstlisting}
Test della proprietà di lossless join
-    Algoritmo 2
    Verifica in cinque passi se una decomposizione
    D ha la proprietà di lossless join rispetto ad un
    insieme F di dipendenze funzionali:

    1.   Creare una matrice S con una riga i per ogni
         relazione Ri nella decomposizione D, e una colonna j
         per ogni attributo Aj in R;
    2.   porre S(i,j) = bij per tutte le entrate della matrice;
         /* ogni bij è un simbolo distinto associato agli indici
         (i,j) */
    ...
\end{lstlisting}
\subsection*{Spiegazione estesa}
[AGGIUNTA] Questa spiegazione collega il contenuto testuale della slide al contesto del corso. La slide enfatizza i concetti mostrati nell'estratto e il loro ruolo nella progettazione di basi di dati relazionali, con particolare attenzione a correttezza semantica, qualità dello schema e compromessi progettuali.
[AGGIUNTA] Fonte integrazione: riformulazione didattica del testo presente nella slide 20 del file Lezione 02 - Algoritmi per la progettazione di DB Relazionali e ulteriori dipendenze.pdf.
\subsection*{Derivazioni}
[AGGIUNTA] Nella slide compaiono espressioni simboliche/relazioni. Derivazione didattica generale:
\begin{align*}
X^+ &= X \\
X^+ &\leftarrow X^+ \cup Y \quad \text{se } (U \to V) \in F \text{ e } U \subseteq X^+ \text{ con } Y=V \\
\text{Stop} &\text{ quando non si aggiungono nuovi attributi in } X^+
\end{align*}
[AGGIUNTA] Questa procedura esplicita il metodo di chiusura usato nelle slide su dipendenze e chiavi.
\begin{itemize}[leftmargin=*]
\item Forma osservata nella slide: \texttt{2.   porre S(i,j) = bij per tutte le entrate della matrice;}
\end{itemize}
\subsection*{Esempio numerico / esercizio}
[AGGIUNTA] Esempio: dato uno schema R(A,B,C) con dipendenze funzionali A->B e B->C, mostrare la chiusura A^+ e discutere se A è superchiave. Soluzione sintetica: A^+={A,B,C}, quindi A determina tutti gli attributi di R ed è superchiave.
\subsection*{Note su assunzioni o integrazioni}
[ASSUNZIONE] L'estrazione OCR può perdere formattazione o simboli della slide originale.
[AGGIUNTA] Le spiegazioni sono estensioni didattiche per rendere espliciti passaggi impliciti nel materiale proiettato.
\paragraph{Riferimento fonte} File: Lezione 02 - Algoritmi per la progettazione di DB Relazionali e ulteriori dipendenze.pdf. Numero slide: 20.
\section{Test della proprietà di lossless join (2) --- Slide 21}
\subsection*{Estratto originale dalla slide 21}
\begin{lstlisting}
Test della proprietà di lossless join (2)
  3.   per ogni riga i che rappresenta lo schema di relazione
       Ri
         per ogni colonna j che rappresenta l'attributo Aj
              se Ri include l'attributo Aj
              allora porre S(i,j) = aj ;
  4.   ripetere quanto segue finché l'esecuzione del ciclo
       non modifica più S
         per ogni dipendenza funzionale X -> Y in F
              per tutte le righe in S, che hanno gli stessi
              simboli nelle colonne corrispondenti
              agli attributi in X:
       ...
\end{lstlisting}
\subsection*{Spiegazione estesa}
[AGGIUNTA] Questa spiegazione collega il contenuto testuale della slide al contesto del corso. La slide enfatizza i concetti mostrati nell'estratto e il loro ruolo nella progettazione di basi di dati relazionali, con particolare attenzione a correttezza semantica, qualità dello schema e compromessi progettuali.
[AGGIUNTA] Fonte integrazione: riformulazione didattica del testo presente nella slide 21 del file Lezione 02 - Algoritmi per la progettazione di DB Relazionali e ulteriori dipendenze.pdf.
\subsection*{Derivazioni}
[AGGIUNTA] Nella slide compaiono espressioni simboliche/relazioni. Derivazione didattica generale:
\begin{align*}
X^+ &= X \\
X^+ &\leftarrow X^+ \cup Y \quad \text{se } (U \to V) \in F \text{ e } U \subseteq X^+ \text{ con } Y=V \\
\text{Stop} &\text{ quando non si aggiungono nuovi attributi in } X^+
\end{align*}
[AGGIUNTA] Questa procedura esplicita il metodo di chiusura usato nelle slide su dipendenze e chiavi.
\begin{itemize}[leftmargin=*]
\item Forma osservata nella slide: \texttt{allora porre S(i,j) = aj ;}
\item Forma osservata nella slide: \texttt{per ogni dipendenza funzionale X -> Y in F}
\end{itemize}
\subsection*{Esempio numerico / esercizio}
[AGGIUNTA] Esempio: dato uno schema R(A,B,C) con dipendenze funzionali A->B e B->C, mostrare la chiusura A^+ e discutere se A è superchiave. Soluzione sintetica: A^+={A,B,C}, quindi A determina tutti gli attributi di R ed è superchiave.
\subsection*{Note su assunzioni o integrazioni}
[ASSUNZIONE] L'estrazione OCR può perdere formattazione o simboli della slide originale.
[AGGIUNTA] Le spiegazioni sono estensioni didattiche per rendere espliciti passaggi impliciti nel materiale proiettato.
\paragraph{Riferimento fonte} File: Lezione 02 - Algoritmi per la progettazione di DB Relazionali e ulteriori dipendenze.pdf. Numero slide: 21.
\section{Test della proprietà di lossless join (3) --- Slide 22}
\subsection*{Estratto originale dalla slide 22}
\begin{lstlisting}
Test della proprietà di lossless join (3)
                  rendere uguali i simboli in ogni colonna
                  che corrispondono ad un attributo in Y,
                  come segue:
                          se una riga ha un simbolo 'a' nella
                          colonna, porre le altre righe allo
                          stesso simbolo 'a' nella colonna.
                          Se non esiste in nessuna riga un
                          simbolo 'a' per l'attributo, scegliere
                          uno dei simboli 'b' che appaiono in
                          una delle righe e porre le altre righe
                          a quel simbolo 'b' nella colonna;
      5.   se una riga contiene solo simboli 'a',
           allora la decomposizione ha la proprietà lossless
           join,
           altrimenti no.
\end{lstlisting}
\subsection*{Spiegazione estesa}
[AGGIUNTA] Questa spiegazione collega il contenuto testuale della slide al contesto del corso. La slide enfatizza i concetti mostrati nell'estratto e il loro ruolo nella progettazione di basi di dati relazionali, con particolare attenzione a correttezza semantica, qualità dello schema e compromessi progettuali.
[AGGIUNTA] Fonte integrazione: riformulazione didattica del testo presente nella slide 22 del file Lezione 02 - Algoritmi per la progettazione di DB Relazionali e ulteriori dipendenze.pdf.
\subsection*{Derivazioni}
[AGGIUNTA] In questa slide non sono state rilevate equazioni esplicite nel testo estratto; non è quindi richiesta una derivazione algebrica puntuale.
\subsection*{Esempio numerico / esercizio}
[AGGIUNTA] Esempio: dato uno schema R(A,B,C) con dipendenze funzionali A->B e B->C, mostrare la chiusura A^+ e discutere se A è superchiave. Soluzione sintetica: A^+={A,B,C}, quindi A determina tutti gli attributi di R ed è superchiave.
\subsection*{Note su assunzioni o integrazioni}
[ASSUNZIONE] L'estrazione OCR può perdere formattazione o simboli della slide originale.
[AGGIUNTA] Le spiegazioni sono estensioni didattiche per rendere espliciti passaggi impliciti nel materiale proiettato.
\paragraph{Riferimento fonte} File: Lezione 02 - Algoritmi per la progettazione di DB Relazionali e ulteriori dipendenze.pdf. Numero slide: 22.
\section{Descrizione dell'algoritmo 2 --- Slide 23}
\subsection*{Estratto originale dalla slide 23}
\begin{lstlisting}
Descrizione dell'algoritmo 2
-    Data una relazione R, decomposta in
    R1, R2, ..., Rm, l'algoritmo crea uno stato di
    relazione r nella matrice S. La riga i-esima
    rappresenta la relazione Ri, con un simbolo 'a' in
    corrispondenza degli attributi di R i, ed un simbolo
    'b' nelle colonne rimanenti.
\end{lstlisting}
\subsection*{Spiegazione estesa}
[AGGIUNTA] Questa spiegazione collega il contenuto testuale della slide al contesto del corso. La slide enfatizza i concetti mostrati nell'estratto e il loro ruolo nella progettazione di basi di dati relazionali, con particolare attenzione a correttezza semantica, qualità dello schema e compromessi progettuali.
[AGGIUNTA] Fonte integrazione: riformulazione didattica del testo presente nella slide 23 del file Lezione 02 - Algoritmi per la progettazione di DB Relazionali e ulteriori dipendenze.pdf.
\subsection*{Derivazioni}
[AGGIUNTA] In questa slide non sono state rilevate equazioni esplicite nel testo estratto; non è quindi richiesta una derivazione algebrica puntuale.
\subsection*{Esempio numerico / esercizio}
[AGGIUNTA] Esempio: dato uno schema R(A,B,C) con dipendenze funzionali A->B e B->C, mostrare la chiusura A^+ e discutere se A è superchiave. Soluzione sintetica: A^+={A,B,C}, quindi A determina tutti gli attributi di R ed è superchiave.
\subsection*{Note su assunzioni o integrazioni}
[ASSUNZIONE] L'estrazione OCR può perdere formattazione o simboli della slide originale.
[AGGIUNTA] Le spiegazioni sono estensioni didattiche per rendere espliciti passaggi impliciti nel materiale proiettato.
\paragraph{Riferimento fonte} File: Lezione 02 - Algoritmi per la progettazione di DB Relazionali e ulteriori dipendenze.pdf. Numero slide: 23.
\section{Descrizione dell'algoritmo 2 (2) --- Slide 24}
\subsection*{Estratto originale dalla slide 24}
\begin{lstlisting}
Descrizione dell'algoritmo 2 (2)
-    L'algoritmo trasforma le righe della matrice (nel
    ciclo del passo 4) in modo che rappresentino
    tuple soddisfacenti tutte le FD in F.

-    Alla fine del loop, per ogni coppia di righe in S per
    i cui attributi vale la parte sinistra X di una FD
    X -> Y, deve valere anche la parte destra Y di tale
    dipendenza.
\end{lstlisting}
\subsection*{Spiegazione estesa}
[AGGIUNTA] Questa spiegazione collega il contenuto testuale della slide al contesto del corso. La slide enfatizza i concetti mostrati nell'estratto e il loro ruolo nella progettazione di basi di dati relazionali, con particolare attenzione a correttezza semantica, qualità dello schema e compromessi progettuali.
[AGGIUNTA] Fonte integrazione: riformulazione didattica del testo presente nella slide 24 del file Lezione 02 - Algoritmi per la progettazione di DB Relazionali e ulteriori dipendenze.pdf.
\subsection*{Derivazioni}
[AGGIUNTA] Nella slide compaiono espressioni simboliche/relazioni. Derivazione didattica generale:
\begin{align*}
X^+ &= X \\
X^+ &\leftarrow X^+ \cup Y \quad \text{se } (U \to V) \in F \text{ e } U \subseteq X^+ \text{ con } Y=V \\
\text{Stop} &\text{ quando non si aggiungono nuovi attributi in } X^+
\end{align*}
[AGGIUNTA] Questa procedura esplicita il metodo di chiusura usato nelle slide su dipendenze e chiavi.
\begin{itemize}[leftmargin=*]
\item Forma osservata nella slide: \texttt{X -> Y, deve valere anche la parte destra Y di tale}
\end{itemize}
\subsection*{Esempio numerico / esercizio}
[AGGIUNTA] Esempio: dato uno schema R(A,B,C) con dipendenze funzionali A->B e B->C, mostrare la chiusura A^+ e discutere se A è superchiave. Soluzione sintetica: A^+={A,B,C}, quindi A determina tutti gli attributi di R ed è superchiave.
\subsection*{Note su assunzioni o integrazioni}
[ASSUNZIONE] L'estrazione OCR può perdere formattazione o simboli della slide originale.
[AGGIUNTA] Le spiegazioni sono estensioni didattiche per rendere espliciti passaggi impliciti nel materiale proiettato.
\paragraph{Riferimento fonte} File: Lezione 02 - Algoritmi per la progettazione di DB Relazionali e ulteriori dipendenze.pdf. Numero slide: 24.
\section{Descrizione dell'algoritmo 2 (3) --- Slide 25}
\subsection*{Estratto originale dalla slide 25}
\begin{lstlisting}
Descrizione dell'algoritmo 2 (3)
-    Può essere dimostrato che dopo il ciclo del passo
    4, la decomposizione D ha la proprietà di lossless
    join se e solo se una riga ha tutti simboli 'a'.

-    Se non si hanno tutti simboli 'a', lo stato di
    relazione r (e quindi la decomposizione D)
    soddisfa le FD di F, ma non gode della proprietà
    di lossless join.
\end{lstlisting}
\subsection*{Spiegazione estesa}
[AGGIUNTA] Questa spiegazione collega il contenuto testuale della slide al contesto del corso. La slide enfatizza i concetti mostrati nell'estratto e il loro ruolo nella progettazione di basi di dati relazionali, con particolare attenzione a correttezza semantica, qualità dello schema e compromessi progettuali.
[AGGIUNTA] Fonte integrazione: riformulazione didattica del testo presente nella slide 25 del file Lezione 02 - Algoritmi per la progettazione di DB Relazionali e ulteriori dipendenze.pdf.
\subsection*{Derivazioni}
[AGGIUNTA] In questa slide non sono state rilevate equazioni esplicite nel testo estratto; non è quindi richiesta una derivazione algebrica puntuale.
\subsection*{Esempio numerico / esercizio}
[AGGIUNTA] Esempio: dato uno schema R(A,B,C) con dipendenze funzionali A->B e B->C, mostrare la chiusura A^+ e discutere se A è superchiave. Soluzione sintetica: A^+={A,B,C}, quindi A determina tutti gli attributi di R ed è superchiave.
\subsection*{Note su assunzioni o integrazioni}
[ASSUNZIONE] L'estrazione OCR può perdere formattazione o simboli della slide originale.
[AGGIUNTA] Le spiegazioni sono estensioni didattiche per rendere espliciti passaggi impliciti nel materiale proiettato.
\paragraph{Riferimento fonte} File: Lezione 02 - Algoritmi per la progettazione di DB Relazionali e ulteriori dipendenze.pdf. Numero slide: 25.
\section{Algoritmo 2: Esempio --- Slide 26}
\subsection*{Estratto originale dalla slide 26}
\begin{lstlisting}
Algoritmo 2: Esempio
\end{lstlisting}
\subsection*{Spiegazione estesa}
[AGGIUNTA] Questa spiegazione collega il contenuto testuale della slide al contesto del corso. La slide enfatizza i concetti mostrati nell'estratto e il loro ruolo nella progettazione di basi di dati relazionali, con particolare attenzione a correttezza semantica, qualità dello schema e compromessi progettuali.
[AGGIUNTA] Fonte integrazione: riformulazione didattica del testo presente nella slide 26 del file Lezione 02 - Algoritmi per la progettazione di DB Relazionali e ulteriori dipendenze.pdf.
\subsection*{Derivazioni}
[AGGIUNTA] In questa slide non sono state rilevate equazioni esplicite nel testo estratto; non è quindi richiesta una derivazione algebrica puntuale.
\subsection*{Esempio numerico / esercizio}
[AGGIUNTA] Esempio: dato uno schema R(A,B,C) con dipendenze funzionali A->B e B->C, mostrare la chiusura A^+ e discutere se A è superchiave. Soluzione sintetica: A^+={A,B,C}, quindi A determina tutti gli attributi di R ed è superchiave.
\subsection*{Note su assunzioni o integrazioni}
[ASSUNZIONE] L'estrazione OCR può perdere formattazione o simboli della slide originale.
[AGGIUNTA] Le spiegazioni sono estensioni didattiche per rendere espliciti passaggi impliciti nel materiale proiettato.
\paragraph{Riferimento fonte} File: Lezione 02 - Algoritmi per la progettazione di DB Relazionali e ulteriori dipendenze.pdf. Numero slide: 26.
\section{Trovare una decomposizione --- Slide 27}
\subsection*{Estratto originale dalla slide 27}
\begin{lstlisting}
Trovare una decomposizione
-    L'algoritmo proposto consente di verificare se una
    decomposizione ha la proprietà di lossless join.

-    Vogliamo trovare una decomposizione di uno
    schema di relazione R che sia in BCNF e che
    presenti la proprietà di lossless join rispetto ad un
    determinato insieme di dipendenze funzionali.
\end{lstlisting}
\subsection*{Spiegazione estesa}
[AGGIUNTA] Questa spiegazione collega il contenuto testuale della slide al contesto del corso. La slide enfatizza i concetti mostrati nell'estratto e il loro ruolo nella progettazione di basi di dati relazionali, con particolare attenzione a correttezza semantica, qualità dello schema e compromessi progettuali.
[AGGIUNTA] Fonte integrazione: riformulazione didattica del testo presente nella slide 27 del file Lezione 02 - Algoritmi per la progettazione di DB Relazionali e ulteriori dipendenze.pdf.
\subsection*{Derivazioni}
[AGGIUNTA] In questa slide non sono state rilevate equazioni esplicite nel testo estratto; non è quindi richiesta una derivazione algebrica puntuale.
\subsection*{Esempio numerico / esercizio}
[AGGIUNTA] Esempio: dato uno schema R(A,B,C) con dipendenze funzionali A->B e B->C, mostrare la chiusura A^+ e discutere se A è superchiave. Soluzione sintetica: A^+={A,B,C}, quindi A determina tutti gli attributi di R ed è superchiave.
\subsection*{Note su assunzioni o integrazioni}
[ASSUNZIONE] L'estrazione OCR può perdere formattazione o simboli della slide originale.
[AGGIUNTA] Le spiegazioni sono estensioni didattiche per rendere espliciti passaggi impliciti nel materiale proiettato.
\paragraph{Riferimento fonte} File: Lezione 02 - Algoritmi per la progettazione di DB Relazionali e ulteriori dipendenze.pdf. Numero slide: 27.
\section{Proprietà delle decomposizioni lossless join --- Slide 28}
\subsection*{Estratto originale dalla slide 28}
\begin{lstlisting}
Proprietà delle decomposizioni lossless join

-    Proprietà LJ1:
    Una decomposizione D = {R1, R2} di R ha la
    proprietà lossless join rispetto a un insieme di
    dipendenze funzionali F su R se e solo se:
      la FD ( (R1   R2) -> (R1 - R2) ) è in F+,
    oppure
      la FD ( (R1   R2) -> (R2 - R1) ) è in F+


-    Questo è un test molto più immediato
    dell'algoritmo 2 ma funziona solo per
    decomposizioni in due schemi.
\end{lstlisting}
\subsection*{Spiegazione estesa}
[AGGIUNTA] Questa spiegazione collega il contenuto testuale della slide al contesto del corso. La slide enfatizza i concetti mostrati nell'estratto e il loro ruolo nella progettazione di basi di dati relazionali, con particolare attenzione a correttezza semantica, qualità dello schema e compromessi progettuali.
[AGGIUNTA] Fonte integrazione: riformulazione didattica del testo presente nella slide 28 del file Lezione 02 - Algoritmi per la progettazione di DB Relazionali e ulteriori dipendenze.pdf.
\subsection*{Derivazioni}
[AGGIUNTA] Nella slide compaiono espressioni simboliche/relazioni. Derivazione didattica generale:
\begin{align*}
X^+ &= X \\
X^+ &\leftarrow X^+ \cup Y \quad \text{se } (U \to V) \in F \text{ e } U \subseteq X^+ \text{ con } Y=V \\
\text{Stop} &\text{ quando non si aggiungono nuovi attributi in } X^+
\end{align*}
[AGGIUNTA] Questa procedura esplicita il metodo di chiusura usato nelle slide su dipendenze e chiavi.
\begin{itemize}[leftmargin=*]
\item Forma osservata nella slide: \texttt{Una decomposizione D = (R1, R2) di R ha la}
\item Forma osservata nella slide: \texttt{la FD ( (R1   R2) -> (R1 - R2) ) è in F+,}
\item Forma osservata nella slide: \texttt{la FD ( (R1   R2) -> (R2 - R1) ) è in F+}
\end{itemize}
\subsection*{Esempio numerico / esercizio}
[AGGIUNTA] Esempio: dato uno schema R(A,B,C) con dipendenze funzionali A->B e B->C, mostrare la chiusura A^+ e discutere se A è superchiave. Soluzione sintetica: A^+={A,B,C}, quindi A determina tutti gli attributi di R ed è superchiave.
\subsection*{Note su assunzioni o integrazioni}
[ASSUNZIONE] L'estrazione OCR può perdere formattazione o simboli della slide originale.
[AGGIUNTA] Le spiegazioni sono estensioni didattiche per rendere espliciti passaggi impliciti nel materiale proiettato.
\paragraph{Riferimento fonte} File: Lezione 02 - Algoritmi per la progettazione di DB Relazionali e ulteriori dipendenze.pdf. Numero slide: 28.
\section{Proprietà delle decomposizioni lossless join (2) --- Slide 29}
\subsection*{Estratto originale dalla slide 29}
\begin{lstlisting}
Proprietà delle decomposizioni lossless join (2)

-    Proprietà LJ2:
    Se una decomposizione D = {R1, R2, ..., Rm} di R
    ha la proprietà LJ rispetto a un insieme di FD F
    su R, e se una decomposizione
    D1= {Q1, Q2, ..., Qk} di Ri ha la proprietà LJ
    rispetto alla proiezione di F su R i allora la
    decomposizione
    D2 = {R1, R2, ..., Ri-1, Q1, Q2, ..., Qk, Ri+1, ... Rm} di R
    ha la proprietà LJ rispetto a F.

-    Usando queste proprietà possiamo sviluppare un
    algoritmo per creare decomposizioni LJ con schemi
    di relazione in BCNF.
\end{lstlisting}
\subsection*{Spiegazione estesa}
[AGGIUNTA] Questa spiegazione collega il contenuto testuale della slide al contesto del corso. La slide enfatizza i concetti mostrati nell'estratto e il loro ruolo nella progettazione di basi di dati relazionali, con particolare attenzione a correttezza semantica, qualità dello schema e compromessi progettuali.
[AGGIUNTA] Fonte integrazione: riformulazione didattica del testo presente nella slide 29 del file Lezione 02 - Algoritmi per la progettazione di DB Relazionali e ulteriori dipendenze.pdf.
\subsection*{Derivazioni}
[AGGIUNTA] Nella slide compaiono espressioni simboliche/relazioni. Derivazione didattica generale:
\begin{align*}
X^+ &= X \\
X^+ &\leftarrow X^+ \cup Y \quad \text{se } (U \to V) \in F \text{ e } U \subseteq X^+ \text{ con } Y=V \\
\text{Stop} &\text{ quando non si aggiungono nuovi attributi in } X^+
\end{align*}
[AGGIUNTA] Questa procedura esplicita il metodo di chiusura usato nelle slide su dipendenze e chiavi.
\begin{itemize}[leftmargin=*]
\item Forma osservata nella slide: \texttt{Se una decomposizione D = (R1, R2, ..., Rm) di R}
\item Forma osservata nella slide: \texttt{D1= (Q1, Q2, ..., Qk) di Ri ha la proprietà LJ}
\item Forma osservata nella slide: \texttt{D2 = (R1, R2, ..., Ri-1, Q1, Q2, ..., Qk, Ri+1, ... Rm) di R}
\end{itemize}
\subsection*{Esempio numerico / esercizio}
[AGGIUNTA] Esempio: dato uno schema R(A,B,C) con dipendenze funzionali A->B e B->C, mostrare la chiusura A^+ e discutere se A è superchiave. Soluzione sintetica: A^+={A,B,C}, quindi A determina tutti gli attributi di R ed è superchiave.
\subsection*{Note su assunzioni o integrazioni}
[ASSUNZIONE] L'estrazione OCR può perdere formattazione o simboli della slide originale.
[AGGIUNTA] Le spiegazioni sono estensioni didattiche per rendere espliciti passaggi impliciti nel materiale proiettato.
\paragraph{Riferimento fonte} File: Lezione 02 - Algoritmi per la progettazione di DB Relazionali e ulteriori dipendenze.pdf. Numero slide: 29.
\section{Decomposizione LJ in schemi in BCNF --- Slide 30}
\subsection*{Estratto originale dalla slide 30}
\begin{lstlisting}
Decomposizione LJ in schemi in BCNF

-     Algoritmo 3
     Consente di ottenere in due passi una
     decomposizione LJ con schemi di relazione in
     BCNF:
1.   porre D := R;
2.   finché ci sta uno schema di relazione Q in D che non è in BCNF
        {
        si scelga uno schema di relazione Q in D che non sia in BCNF;
        si trovi una dipendenza funzionale X -> Y che violi BCNF;
        si sostituisca Q in D con due schemi di relazione (Q-Y) e (X   Y).
        }
\end{lstlisting}
\subsection*{Spiegazione estesa}
[AGGIUNTA] Questa spiegazione collega il contenuto testuale della slide al contesto del corso. La slide enfatizza i concetti mostrati nell'estratto e il loro ruolo nella progettazione di basi di dati relazionali, con particolare attenzione a correttezza semantica, qualità dello schema e compromessi progettuali.
[AGGIUNTA] Fonte integrazione: riformulazione didattica del testo presente nella slide 30 del file Lezione 02 - Algoritmi per la progettazione di DB Relazionali e ulteriori dipendenze.pdf.
\subsection*{Derivazioni}
[AGGIUNTA] Nella slide compaiono espressioni simboliche/relazioni. Derivazione didattica generale:
\begin{align*}
X^+ &= X \\
X^+ &\leftarrow X^+ \cup Y \quad \text{se } (U \to V) \in F \text{ e } U \subseteq X^+ \text{ con } Y=V \\
\text{Stop} &\text{ quando non si aggiungono nuovi attributi in } X^+
\end{align*}
[AGGIUNTA] Questa procedura esplicita il metodo di chiusura usato nelle slide su dipendenze e chiavi.
\begin{itemize}[leftmargin=*]
\item Forma osservata nella slide: \texttt{1.   porre D := R;}
\item Forma osservata nella slide: \texttt{si trovi una dipendenza funzionale X -> Y che violi BCNF;}
\end{itemize}
\subsection*{Esempio numerico / esercizio}
[AGGIUNTA] Esempio: dato uno schema R(A,B,C) con dipendenze funzionali A->B e B->C, mostrare la chiusura A^+ e discutere se A è superchiave. Soluzione sintetica: A^+={A,B,C}, quindi A determina tutti gli attributi di R ed è superchiave.
\subsection*{Note su assunzioni o integrazioni}
[ASSUNZIONE] L'estrazione OCR può perdere formattazione o simboli della slide originale.
[AGGIUNTA] Le spiegazioni sono estensioni didattiche per rendere espliciti passaggi impliciti nel materiale proiettato.
\paragraph{Riferimento fonte} File: Lezione 02 - Algoritmi per la progettazione di DB Relazionali e ulteriori dipendenze.pdf. Numero slide: 30.
\section{Decomposizione LJ e dependency-preserving --- Slide 31}
\subsection*{Estratto originale dalla slide 31}
\begin{lstlisting}
Decomposizione LJ e dependency-preserving
in schemi di relazione in 3NF
-    Algoritmo 4
    Consente di ottenere in tre passi una
    decomposizione LJ e dependency-preserving con
    schemi di relazione in 3NF:
    1.  Trovare una copertura minimale G per F;
    2.  per ogni parte sinistra X di una FD in G, creare uno
        schema di relazione in D con attributi
    {X   A1   A2   ...   Am} dove X -> A1, X -> A2, ..., X -> Am
        sono le sole dipendenze in G aventi X come parte
        sinistra.
    3.  Se nessuno degli schemi di relazione in D contiene una
        chiave di R, creare un altro schema di relazione che
        contiene attributi che formano una chiave per R.
\end{lstlisting}
\subsection*{Spiegazione estesa}
[AGGIUNTA] Questa spiegazione collega il contenuto testuale della slide al contesto del corso. La slide enfatizza i concetti mostrati nell'estratto e il loro ruolo nella progettazione di basi di dati relazionali, con particolare attenzione a correttezza semantica, qualità dello schema e compromessi progettuali.
[AGGIUNTA] Fonte integrazione: riformulazione didattica del testo presente nella slide 31 del file Lezione 02 - Algoritmi per la progettazione di DB Relazionali e ulteriori dipendenze.pdf.
\subsection*{Derivazioni}
[AGGIUNTA] Nella slide compaiono espressioni simboliche/relazioni. Derivazione didattica generale:
\begin{align*}
X^+ &= X \\
X^+ &\leftarrow X^+ \cup Y \quad \text{se } (U \to V) \in F \text{ e } U \subseteq X^+ \text{ con } Y=V \\
\text{Stop} &\text{ quando non si aggiungono nuovi attributi in } X^+
\end{align*}
[AGGIUNTA] Questa procedura esplicita il metodo di chiusura usato nelle slide su dipendenze e chiavi.
\begin{itemize}[leftmargin=*]
\item Forma osservata nella slide: \texttt{(X   A1   A2   ...   Am) dove X -> A1, X -> A2, ..., X -> Am}
\end{itemize}
\subsection*{Esempio numerico / esercizio}
[AGGIUNTA] Esempio: dato uno schema R(A,B,C) con dipendenze funzionali A->B e B->C, mostrare la chiusura A^+ e discutere se A è superchiave. Soluzione sintetica: A^+={A,B,C}, quindi A determina tutti gli attributi di R ed è superchiave.
\subsection*{Note su assunzioni o integrazioni}
[ASSUNZIONE] L'estrazione OCR può perdere formattazione o simboli della slide originale.
[AGGIUNTA] Le spiegazioni sono estensioni didattiche per rendere espliciti passaggi impliciti nel materiale proiettato.
\paragraph{Riferimento fonte} File: Lezione 02 - Algoritmi per la progettazione di DB Relazionali e ulteriori dipendenze.pdf. Numero slide: 31.
\section{Trovare una chiave K per uno schema di --- Slide 32}
\subsection*{Estratto originale dalla slide 32}
\begin{lstlisting}
Trovare una chiave K per uno schema di
relazione R
-    L'algoritmo che segue consente di identificare
    una chiave K di R, come richiesto nel passo 3.
    La chiave restituita dipende dall'ordine in cui gli
    attributi sono stati rimossi al passo 2.
-    Algoritmo 4a:
    1.   porre K := R;
    2.   per ogni attributo A in K
         {
           calcolare (K - A)+ rispetto all' insieme di FD;
           se (K - A)+ contiene tutti gli attributi in R,
           allora K = K - {A};
         }
\end{lstlisting}
\subsection*{Spiegazione estesa}
[AGGIUNTA] Questa spiegazione collega il contenuto testuale della slide al contesto del corso. La slide enfatizza i concetti mostrati nell'estratto e il loro ruolo nella progettazione di basi di dati relazionali, con particolare attenzione a correttezza semantica, qualità dello schema e compromessi progettuali.
[AGGIUNTA] Fonte integrazione: riformulazione didattica del testo presente nella slide 32 del file Lezione 02 - Algoritmi per la progettazione di DB Relazionali e ulteriori dipendenze.pdf.
\subsection*{Derivazioni}
[AGGIUNTA] Nella slide compaiono espressioni simboliche/relazioni. Derivazione didattica generale:
\begin{align*}
X^+ &= X \\
X^+ &\leftarrow X^+ \cup Y \quad \text{se } (U \to V) \in F \text{ e } U \subseteq X^+ \text{ con } Y=V \\
\text{Stop} &\text{ quando non si aggiungono nuovi attributi in } X^+
\end{align*}
[AGGIUNTA] Questa procedura esplicita il metodo di chiusura usato nelle slide su dipendenze e chiavi.
\begin{itemize}[leftmargin=*]
\item Forma osservata nella slide: \texttt{1.   porre K := R;}
\item Forma osservata nella slide: \texttt{allora K = K - (A);}
\end{itemize}
\subsection*{Esempio numerico / esercizio}
[AGGIUNTA] Esempio: dato uno schema R(A,B,C) con dipendenze funzionali A->B e B->C, mostrare la chiusura A^+ e discutere se A è superchiave. Soluzione sintetica: A^+={A,B,C}, quindi A determina tutti gli attributi di R ed è superchiave.
\subsection*{Note su assunzioni o integrazioni}
[ASSUNZIONE] L'estrazione OCR può perdere formattazione o simboli della slide originale.
[AGGIUNTA] Le spiegazioni sono estensioni didattiche per rendere espliciti passaggi impliciti nel materiale proiettato.
\paragraph{Riferimento fonte} File: Lezione 02 - Algoritmi per la progettazione di DB Relazionali e ulteriori dipendenze.pdf. Numero slide: 32.
\section{Nota --- Slide 33}
\subsection*{Estratto originale dalla slide 33}
\begin{lstlisting}
Nota
-    Non è sempre possibile trovare una
    decomposizione LJ che sia dependency-
    preserving e i cui schemi di relazione siano in
    BCNF.
\end{lstlisting}
\subsection*{Spiegazione estesa}
[AGGIUNTA] Questa spiegazione collega il contenuto testuale della slide al contesto del corso. La slide enfatizza i concetti mostrati nell'estratto e il loro ruolo nella progettazione di basi di dati relazionali, con particolare attenzione a correttezza semantica, qualità dello schema e compromessi progettuali.
[AGGIUNTA] Fonte integrazione: riformulazione didattica del testo presente nella slide 33 del file Lezione 02 - Algoritmi per la progettazione di DB Relazionali e ulteriori dipendenze.pdf.
\subsection*{Derivazioni}
[AGGIUNTA] In questa slide non sono state rilevate equazioni esplicite nel testo estratto; non è quindi richiesta una derivazione algebrica puntuale.
\subsection*{Esempio numerico / esercizio}
[AGGIUNTA] Esempio: dato uno schema R(A,B,C) con dipendenze funzionali A->B e B->C, mostrare la chiusura A^+ e discutere se A è superchiave. Soluzione sintetica: A^+={A,B,C}, quindi A determina tutti gli attributi di R ed è superchiave.
\subsection*{Note su assunzioni o integrazioni}
[ASSUNZIONE] L'estrazione OCR può perdere formattazione o simboli della slide originale.
[AGGIUNTA] Le spiegazioni sono estensioni didattiche per rendere espliciti passaggi impliciti nel materiale proiettato.
\paragraph{Riferimento fonte} File: Lezione 02 - Algoritmi per la progettazione di DB Relazionali e ulteriori dipendenze.pdf. Numero slide: 33.
\section{Decomposizioni e valori null --- Slide 34}
\subsection*{Estratto originale dalla slide 34}
\begin{lstlisting}
Decomposizioni e valori null
-    La teoria delle decomposizioni LJ si basa sulla
    seguente assunzione:
    nessun valore null è ammesso per gli attributi di
    join.

-    Non esiste una teoria di progettazione di
    database relazionali con i valori null
    soddisfacente.
\end{lstlisting}
\subsection*{Spiegazione estesa}
[AGGIUNTA] Questa spiegazione collega il contenuto testuale della slide al contesto del corso. La slide enfatizza i concetti mostrati nell'estratto e il loro ruolo nella progettazione di basi di dati relazionali, con particolare attenzione a correttezza semantica, qualità dello schema e compromessi progettuali.
[AGGIUNTA] Fonte integrazione: riformulazione didattica del testo presente nella slide 34 del file Lezione 02 - Algoritmi per la progettazione di DB Relazionali e ulteriori dipendenze.pdf.
\subsection*{Derivazioni}
[AGGIUNTA] In questa slide non sono state rilevate equazioni esplicite nel testo estratto; non è quindi richiesta una derivazione algebrica puntuale.
\subsection*{Esempio numerico / esercizio}
[AGGIUNTA] Esempio: dato uno schema R(A,B,C) con dipendenze funzionali A->B e B->C, mostrare la chiusura A^+ e discutere se A è superchiave. Soluzione sintetica: A^+={A,B,C}, quindi A determina tutti gli attributi di R ed è superchiave.
\subsection*{Note su assunzioni o integrazioni}
[ASSUNZIONE] L'estrazione OCR può perdere formattazione o simboli della slide originale.
[AGGIUNTA] Le spiegazioni sono estensioni didattiche per rendere espliciti passaggi impliciti nel materiale proiettato.
\paragraph{Riferimento fonte} File: Lezione 02 - Algoritmi per la progettazione di DB Relazionali e ulteriori dipendenze.pdf. Numero slide: 34.
\section{Decomposizioni e valori null: Esempio --- Slide 35}
\subsection*{Estratto originale dalla slide 35}
\begin{lstlisting}
Decomposizioni e valori null: Esempio
-    Due impiegati non sono ancora assegnati a un
    dipartimento ...
\end{lstlisting}
\subsection*{Spiegazione estesa}
[AGGIUNTA] Questa spiegazione collega il contenuto testuale della slide al contesto del corso. La slide enfatizza i concetti mostrati nell'estratto e il loro ruolo nella progettazione di basi di dati relazionali, con particolare attenzione a correttezza semantica, qualità dello schema e compromessi progettuali.
[AGGIUNTA] Fonte integrazione: riformulazione didattica del testo presente nella slide 35 del file Lezione 02 - Algoritmi per la progettazione di DB Relazionali e ulteriori dipendenze.pdf.
\subsection*{Derivazioni}
[AGGIUNTA] In questa slide non sono state rilevate equazioni esplicite nel testo estratto; non è quindi richiesta una derivazione algebrica puntuale.
\subsection*{Esempio numerico / esercizio}
[AGGIUNTA] Esempio: dato uno schema R(A,B,C) con dipendenze funzionali A->B e B->C, mostrare la chiusura A^+ e discutere se A è superchiave. Soluzione sintetica: A^+={A,B,C}, quindi A determina tutti gli attributi di R ed è superchiave.
\subsection*{Note su assunzioni o integrazioni}
[ASSUNZIONE] L'estrazione OCR può perdere formattazione o simboli della slide originale.
[AGGIUNTA] Le spiegazioni sono estensioni didattiche per rendere espliciti passaggi impliciti nel materiale proiettato.
\paragraph{Riferimento fonte} File: Lezione 02 - Algoritmi per la progettazione di DB Relazionali e ulteriori dipendenze.pdf. Numero slide: 35.
\section{Decomposizioni e valori null: Esempio --- Slide 36}
\subsection*{Estratto originale dalla slide 36}
\begin{lstlisting}
Decomposizioni e valori null: Esempio
   -    Siano interessati ad una lista (ENAME,
       DNAME) per tutti gli impiegati.




   -    Risultato di un Natural-Join: mancano gli
       impiegati non assegnati.
\end{lstlisting}
\subsection*{Spiegazione estesa}
[AGGIUNTA] Questa spiegazione collega il contenuto testuale della slide al contesto del corso. La slide enfatizza i concetti mostrati nell'estratto e il loro ruolo nella progettazione di basi di dati relazionali, con particolare attenzione a correttezza semantica, qualità dello schema e compromessi progettuali.
[AGGIUNTA] Fonte integrazione: riformulazione didattica del testo presente nella slide 36 del file Lezione 02 - Algoritmi per la progettazione di DB Relazionali e ulteriori dipendenze.pdf.
\subsection*{Derivazioni}
[AGGIUNTA] In questa slide non sono state rilevate equazioni esplicite nel testo estratto; non è quindi richiesta una derivazione algebrica puntuale.
\subsection*{Esempio numerico / esercizio}
[AGGIUNTA] Esempio: dato uno schema R(A,B,C) con dipendenze funzionali A->B e B->C, mostrare la chiusura A^+ e discutere se A è superchiave. Soluzione sintetica: A^+={A,B,C}, quindi A determina tutti gli attributi di R ed è superchiave.
\subsection*{Note su assunzioni o integrazioni}
[ASSUNZIONE] L'estrazione OCR può perdere formattazione o simboli della slide originale.
[AGGIUNTA] Le spiegazioni sono estensioni didattiche per rendere espliciti passaggi impliciti nel materiale proiettato.
\paragraph{Riferimento fonte} File: Lezione 02 - Algoritmi per la progettazione di DB Relazionali e ulteriori dipendenze.pdf. Numero slide: 36.
\section{Decomposizioni e valori null: Esempio --- Slide 37}
\subsection*{Estratto originale dalla slide 37}
\begin{lstlisting}
Decomposizioni e valori null: Esempio




-    Risultato di un Outer-Join: abbiamo tutti gli
    impiegati, ma anche dei valori null per i
    dipartimenti.
\end{lstlisting}
\subsection*{Spiegazione estesa}
[AGGIUNTA] Questa spiegazione collega il contenuto testuale della slide al contesto del corso. La slide enfatizza i concetti mostrati nell'estratto e il loro ruolo nella progettazione di basi di dati relazionali, con particolare attenzione a correttezza semantica, qualità dello schema e compromessi progettuali.
[AGGIUNTA] Fonte integrazione: riformulazione didattica del testo presente nella slide 37 del file Lezione 02 - Algoritmi per la progettazione di DB Relazionali e ulteriori dipendenze.pdf.
\subsection*{Derivazioni}
[AGGIUNTA] In questa slide non sono state rilevate equazioni esplicite nel testo estratto; non è quindi richiesta una derivazione algebrica puntuale.
\subsection*{Esempio numerico / esercizio}
[AGGIUNTA] Esempio: dato uno schema R(A,B,C) con dipendenze funzionali A->B e B->C, mostrare la chiusura A^+ e discutere se A è superchiave. Soluzione sintetica: A^+={A,B,C}, quindi A determina tutti gli attributi di R ed è superchiave.
\subsection*{Note su assunzioni o integrazioni}
[ASSUNZIONE] L'estrazione OCR può perdere formattazione o simboli della slide originale.
[AGGIUNTA] Le spiegazioni sono estensioni didattiche per rendere espliciti passaggi impliciti nel materiale proiettato.
\paragraph{Riferimento fonte} File: Lezione 02 - Algoritmi per la progettazione di DB Relazionali e ulteriori dipendenze.pdf. Numero slide: 37.
\section{Decomposizioni e valori null (2) --- Slide 38}
\subsection*{Estratto originale dalla slide 38}
\begin{lstlisting}
Decomposizioni e valori null (2)
-    Potenziali problemi con valori null possono
    sorgere nell'uso di:
      Join, con perdita di informazioni se non si utilizza
      un Outer-Join.
      Funzioni di aggregazione, quali SUM o AVERAGE,
      con comportamenti non sempre definiti su null.
\end{lstlisting}
\subsection*{Spiegazione estesa}
[AGGIUNTA] Questa spiegazione collega il contenuto testuale della slide al contesto del corso. La slide enfatizza i concetti mostrati nell'estratto e il loro ruolo nella progettazione di basi di dati relazionali, con particolare attenzione a correttezza semantica, qualità dello schema e compromessi progettuali.
[AGGIUNTA] Fonte integrazione: riformulazione didattica del testo presente nella slide 38 del file Lezione 02 - Algoritmi per la progettazione di DB Relazionali e ulteriori dipendenze.pdf.
\subsection*{Derivazioni}
[AGGIUNTA] In questa slide non sono state rilevate equazioni esplicite nel testo estratto; non è quindi richiesta una derivazione algebrica puntuale.
\subsection*{Esempio numerico / esercizio}
[AGGIUNTA] Esempio: dato uno schema R(A,B,C) con dipendenze funzionali A->B e B->C, mostrare la chiusura A^+ e discutere se A è superchiave. Soluzione sintetica: A^+={A,B,C}, quindi A determina tutti gli attributi di R ed è superchiave.
\subsection*{Note su assunzioni o integrazioni}
[ASSUNZIONE] L'estrazione OCR può perdere formattazione o simboli della slide originale.
[AGGIUNTA] Le spiegazioni sono estensioni didattiche per rendere espliciti passaggi impliciti nel materiale proiettato.
\paragraph{Riferimento fonte} File: Lezione 02 - Algoritmi per la progettazione di DB Relazionali e ulteriori dipendenze.pdf. Numero slide: 38.
\section{Algoritmi di normalizzazione --- Slide 39}
\subsection*{Estratto originale dalla slide 39}
\begin{lstlisting}
Algoritmi di normalizzazione
Algoritmo         Input                   Output                Proprietà/Scopo       Note

                  Insieme F di            Un insieme di         Conservazione delle   Nessuna garanzia
                  dipendenze              relazioni in 3NF      dipendenze            sulla verifica della
   Algoritmo 1
                  funzionali                                                          proprietà di join
                                                                                      senza perdita.
                  Una decomposizione      Risultato booleano;   Test per la           Esiste un test più
                  D di R e un insieme F   si o no per la        decomposizione        semplice per le
   Algoritmo 2    di dipendenze           proprietà di join     non-additiva          decomposizioni
                  funzionale              non-additivo                                binarie.

                  Insieme F di            Un insieme di         Decomposizione        Nessuna garanzia
                  dipendenze              relazioni in BCNF     non-additiva          della conservazione
   Algoritmo 3
                  funzionali                                                          delle dipendenze.

                  Insieme F di            Un insieme di         Decomposizione        Può capitare di non
                  dipendenze              relazioni in 3NF      non-additiva e con    avere una BCNF; ma
                  funzionali                                    conservazione delle   si ottiene la 3NF e
   Algoritmo 4
                                                                dipendenze            tutte le proprietà
                                                                                      auspicabili.

                  Schema di relazione     Chiave K di R         Per trovare una       L'intera relazione R
                  R con un insieme F di                         chiave K (cioè un     è sempre una
   Algoritmo 4a   dipendenze                                    sottoinsieme di R)    superchiave di
                  funzionali                                                          default.
\end{lstlisting}
\subsection*{Spiegazione estesa}
[AGGIUNTA] Questa spiegazione collega il contenuto testuale della slide al contesto del corso. La slide enfatizza i concetti mostrati nell'estratto e il loro ruolo nella progettazione di basi di dati relazionali, con particolare attenzione a correttezza semantica, qualità dello schema e compromessi progettuali.
[AGGIUNTA] Fonte integrazione: riformulazione didattica del testo presente nella slide 39 del file Lezione 02 - Algoritmi per la progettazione di DB Relazionali e ulteriori dipendenze.pdf.
\subsection*{Derivazioni}
[AGGIUNTA] In questa slide non sono state rilevate equazioni esplicite nel testo estratto; non è quindi richiesta una derivazione algebrica puntuale.
\subsection*{Esempio numerico / esercizio}
[AGGIUNTA] Esempio: dato uno schema R(A,B,C) con dipendenze funzionali A->B e B->C, mostrare la chiusura A^+ e discutere se A è superchiave. Soluzione sintetica: A^+={A,B,C}, quindi A determina tutti gli attributi di R ed è superchiave.
\subsection*{Note su assunzioni o integrazioni}
[ASSUNZIONE] L'estrazione OCR può perdere formattazione o simboli della slide originale.
[AGGIUNTA] Le spiegazioni sono estensioni didattiche per rendere espliciti passaggi impliciti nel materiale proiettato.
\paragraph{Riferimento fonte} File: Lezione 02 - Algoritmi per la progettazione di DB Relazionali e ulteriori dipendenze.pdf. Numero slide: 39.
\section{Dipendenze Multivalore --- Slide 40}
\subsection*{Estratto originale dalla slide 40}
\begin{lstlisting}
Dipendenze Multivalore
\end{lstlisting}
\subsection*{Spiegazione estesa}
[AGGIUNTA] Questa spiegazione collega il contenuto testuale della slide al contesto del corso. La slide enfatizza i concetti mostrati nell'estratto e il loro ruolo nella progettazione di basi di dati relazionali, con particolare attenzione a correttezza semantica, qualità dello schema e compromessi progettuali.
[AGGIUNTA] Fonte integrazione: riformulazione didattica del testo presente nella slide 40 del file Lezione 02 - Algoritmi per la progettazione di DB Relazionali e ulteriori dipendenze.pdf.
\subsection*{Derivazioni}
[AGGIUNTA] In questa slide non sono state rilevate equazioni esplicite nel testo estratto; non è quindi richiesta una derivazione algebrica puntuale.
\subsection*{Esempio numerico / esercizio}
[AGGIUNTA] Esempio: dato uno schema R(A,B,C) con dipendenze funzionali A->B e B->C, mostrare la chiusura A^+ e discutere se A è superchiave. Soluzione sintetica: A^+={A,B,C}, quindi A determina tutti gli attributi di R ed è superchiave.
\subsection*{Note su assunzioni o integrazioni}
[ASSUNZIONE] L'estrazione OCR può perdere formattazione o simboli della slide originale.
[AGGIUNTA] Le spiegazioni sono estensioni didattiche per rendere espliciti passaggi impliciti nel materiale proiettato.
\paragraph{Riferimento fonte} File: Lezione 02 - Algoritmi per la progettazione di DB Relazionali e ulteriori dipendenze.pdf. Numero slide: 40.
\section{Dipendenze Multivalore (MVD) --- Slide 41}
\subsection*{Estratto originale dalla slide 41}
\begin{lstlisting}
Dipendenze Multivalore (MVD)
-    Sono conseguenza della 1NF, che non consente
    ad un attributo di assumere un insieme di valori.

-    In presenza di due o più attributi indipendenti
    multivalued nello stesso schema, siamo costretti
    a ripetere ogni valore di un attributo con ogni
    valore dell'altro.
\end{lstlisting}
\subsection*{Spiegazione estesa}
[AGGIUNTA] Questa spiegazione collega il contenuto testuale della slide al contesto del corso. La slide enfatizza i concetti mostrati nell'estratto e il loro ruolo nella progettazione di basi di dati relazionali, con particolare attenzione a correttezza semantica, qualità dello schema e compromessi progettuali.
[AGGIUNTA] Fonte integrazione: riformulazione didattica del testo presente nella slide 41 del file Lezione 02 - Algoritmi per la progettazione di DB Relazionali e ulteriori dipendenze.pdf.
\subsection*{Derivazioni}
[AGGIUNTA] In questa slide non sono state rilevate equazioni esplicite nel testo estratto; non è quindi richiesta una derivazione algebrica puntuale.
\subsection*{Esempio numerico / esercizio}
[AGGIUNTA] Esempio: dato uno schema R(A,B,C) con dipendenze funzionali A->B e B->C, mostrare la chiusura A^+ e discutere se A è superchiave. Soluzione sintetica: A^+={A,B,C}, quindi A determina tutti gli attributi di R ed è superchiave.
\subsection*{Note su assunzioni o integrazioni}
[ASSUNZIONE] L'estrazione OCR può perdere formattazione o simboli della slide originale.
[AGGIUNTA] Le spiegazioni sono estensioni didattiche per rendere espliciti passaggi impliciti nel materiale proiettato.
\paragraph{Riferimento fonte} File: Lezione 02 - Algoritmi per la progettazione di DB Relazionali e ulteriori dipendenze.pdf. Numero slide: 41.
\section{Dipendenze Multivalore: Esempio --- Slide 42}
\subsection*{Estratto originale dalla slide 42}
\begin{lstlisting}
Dipendenze Multivalore: Esempio




-    Questo vincolo è specificato come una
    dipendenza multivalued sulla relazione EMP:
      Informalmente, si può avere una MVD ogni volta che si
     fondono due relazioni indipendenti, A:B e A:C con
     cardinalità 1:N, nella stessa relazione.
\end{lstlisting}
\subsection*{Spiegazione estesa}
[AGGIUNTA] Questa spiegazione collega il contenuto testuale della slide al contesto del corso. La slide enfatizza i concetti mostrati nell'estratto e il loro ruolo nella progettazione di basi di dati relazionali, con particolare attenzione a correttezza semantica, qualità dello schema e compromessi progettuali.
[AGGIUNTA] Fonte integrazione: riformulazione didattica del testo presente nella slide 42 del file Lezione 02 - Algoritmi per la progettazione di DB Relazionali e ulteriori dipendenze.pdf.
\subsection*{Derivazioni}
[AGGIUNTA] In questa slide non sono state rilevate equazioni esplicite nel testo estratto; non è quindi richiesta una derivazione algebrica puntuale.
\subsection*{Esempio numerico / esercizio}
[AGGIUNTA] Esempio: dato uno schema R(A,B,C) con dipendenze funzionali A->B e B->C, mostrare la chiusura A^+ e discutere se A è superchiave. Soluzione sintetica: A^+={A,B,C}, quindi A determina tutti gli attributi di R ed è superchiave.
\subsection*{Note su assunzioni o integrazioni}
[ASSUNZIONE] L'estrazione OCR può perdere formattazione o simboli della slide originale.
[AGGIUNTA] Le spiegazioni sono estensioni didattiche per rendere espliciti passaggi impliciti nel materiale proiettato.
\paragraph{Riferimento fonte} File: Lezione 02 - Algoritmi per la progettazione di DB Relazionali e ulteriori dipendenze.pdf. Numero slide: 42.
\section{Definizione formale di MVD: --- Slide 43}
\subsection*{Estratto originale dalla slide 43}
\begin{lstlisting}
Definizione formale di MVD:
X multidetermina Y
-    Definizione:
    Una dipendenza multivalore X   Y, specificata su
    uno schema di relazione R, dove X e Y sono
    sottoinsiemi di R, specifica il seguente vincolo su
    qualche relazione r di R:
      Se esistono in r due tuple t1 e t2, con t1[X] = t2[X], allora
      dovrebbero esistere anche t 3 e t4 con le seguenti
      proprietà:
        t3[X] = t4[X] = t1[X] = t2[X]
        t3[Y] = t1[Y] e t4[Y] = t2[Y]
        Dato Z = R - (X   Y), allora
        t3[Z] = t2[Z] e t4[Z] = t1[Z]
\end{lstlisting}
\subsection*{Spiegazione estesa}
[AGGIUNTA] Questa spiegazione collega il contenuto testuale della slide al contesto del corso. La slide enfatizza i concetti mostrati nell'estratto e il loro ruolo nella progettazione di basi di dati relazionali, con particolare attenzione a correttezza semantica, qualità dello schema e compromessi progettuali.
[AGGIUNTA] Fonte integrazione: riformulazione didattica del testo presente nella slide 43 del file Lezione 02 - Algoritmi per la progettazione di DB Relazionali e ulteriori dipendenze.pdf.
\subsection*{Derivazioni}
[AGGIUNTA] Nella slide compaiono espressioni simboliche/relazioni. Derivazione didattica generale:
\begin{align*}
X^+ &= X \\
X^+ &\leftarrow X^+ \cup Y \quad \text{se } (U \to V) \in F \text{ e } U \subseteq X^+ \text{ con } Y=V \\
\text{Stop} &\text{ quando non si aggiungono nuovi attributi in } X^+
\end{align*}
[AGGIUNTA] Questa procedura esplicita il metodo di chiusura usato nelle slide su dipendenze e chiavi.
\begin{itemize}[leftmargin=*]
\item Forma osservata nella slide: \texttt{Se esistono in r due tuple t1 e t2, con t1[X] = t2[X], allora}
\item Forma osservata nella slide: \texttt{t3[X] = t4[X] = t1[X] = t2[X]}
\item Forma osservata nella slide: \texttt{t3[Y] = t1[Y] e t4[Y] = t2[Y]}
\end{itemize}
\subsection*{Esempio numerico / esercizio}
[AGGIUNTA] Esempio: dato uno schema R(A,B,C) con dipendenze funzionali A->B e B->C, mostrare la chiusura A^+ e discutere se A è superchiave. Soluzione sintetica: A^+={A,B,C}, quindi A determina tutti gli attributi di R ed è superchiave.
\subsection*{Note su assunzioni o integrazioni}
[ASSUNZIONE] L'estrazione OCR può perdere formattazione o simboli della slide originale.
[AGGIUNTA] Le spiegazioni sono estensioni didattiche per rendere espliciti passaggi impliciti nel materiale proiettato.
\paragraph{Riferimento fonte} File: Lezione 02 - Algoritmi per la progettazione di DB Relazionali e ulteriori dipendenze.pdf. Numero slide: 43.
\section{Definizione formale di MVD (2) --- Slide 44}
\subsection*{Estratto originale dalla slide 44}
\begin{lstlisting}
Definizione formale di MVD (2)
-    Data la simmetria della definizione,
    posto Z = R - (X   Y),
    si ha che se X   Y allora X   Z:
      X   Y   X   Z e perciò talvolta lo si scrive come
       X Y|Z
\end{lstlisting}
\subsection*{Spiegazione estesa}
[AGGIUNTA] Questa spiegazione collega il contenuto testuale della slide al contesto del corso. La slide enfatizza i concetti mostrati nell'estratto e il loro ruolo nella progettazione di basi di dati relazionali, con particolare attenzione a correttezza semantica, qualità dello schema e compromessi progettuali.
[AGGIUNTA] Fonte integrazione: riformulazione didattica del testo presente nella slide 44 del file Lezione 02 - Algoritmi per la progettazione di DB Relazionali e ulteriori dipendenze.pdf.
\subsection*{Derivazioni}
[AGGIUNTA] Nella slide compaiono espressioni simboliche/relazioni. Derivazione didattica generale:
\begin{align*}
X^+ &= X \\
X^+ &\leftarrow X^+ \cup Y \quad \text{se } (U \to V) \in F \text{ e } U \subseteq X^+ \text{ con } Y=V \\
\text{Stop} &\text{ quando non si aggiungono nuovi attributi in } X^+
\end{align*}
[AGGIUNTA] Questa procedura esplicita il metodo di chiusura usato nelle slide su dipendenze e chiavi.
\begin{itemize}[leftmargin=*]
\item Forma osservata nella slide: \texttt{posto Z = R - (X   Y),}
\end{itemize}
\subsection*{Esempio numerico / esercizio}
[AGGIUNTA] Esempio: dato uno schema R(A,B,C) con dipendenze funzionali A->B e B->C, mostrare la chiusura A^+ e discutere se A è superchiave. Soluzione sintetica: A^+={A,B,C}, quindi A determina tutti gli attributi di R ed è superchiave.
\subsection*{Note su assunzioni o integrazioni}
[ASSUNZIONE] L'estrazione OCR può perdere formattazione o simboli della slide originale.
[AGGIUNTA] Le spiegazioni sono estensioni didattiche per rendere espliciti passaggi impliciti nel materiale proiettato.
\paragraph{Riferimento fonte} File: Lezione 02 - Algoritmi per la progettazione di DB Relazionali e ulteriori dipendenze.pdf. Numero slide: 44.
\section{MVD: Esempio --- Slide 45}
\subsection*{Estratto originale dalla slide 45}
\begin{lstlisting}
MVD: Esempio
-    Dato un valore di X, l'insieme di valori di Y
    determinati da X è completamente determinato
    solo da X e non dipendono dai valori dei
    rimanenti attributi Z di R:




      Nella relazione EMP valgono:
       ENAME   PNAME e
       ENAME   DNAME
       (o ENAME   PNAME | DNAME)
\end{lstlisting}
\subsection*{Spiegazione estesa}
[AGGIUNTA] Questa spiegazione collega il contenuto testuale della slide al contesto del corso. La slide enfatizza i concetti mostrati nell'estratto e il loro ruolo nella progettazione di basi di dati relazionali, con particolare attenzione a correttezza semantica, qualità dello schema e compromessi progettuali.
[AGGIUNTA] Fonte integrazione: riformulazione didattica del testo presente nella slide 45 del file Lezione 02 - Algoritmi per la progettazione di DB Relazionali e ulteriori dipendenze.pdf.
\subsection*{Derivazioni}
[AGGIUNTA] In questa slide non sono state rilevate equazioni esplicite nel testo estratto; non è quindi richiesta una derivazione algebrica puntuale.
\subsection*{Esempio numerico / esercizio}
[AGGIUNTA] Esempio: dato uno schema R(A,B,C) con dipendenze funzionali A->B e B->C, mostrare la chiusura A^+ e discutere se A è superchiave. Soluzione sintetica: A^+={A,B,C}, quindi A determina tutti gli attributi di R ed è superchiave.
\subsection*{Note su assunzioni o integrazioni}
[ASSUNZIONE] L'estrazione OCR può perdere formattazione o simboli della slide originale.
[AGGIUNTA] Le spiegazioni sono estensioni didattiche per rendere espliciti passaggi impliciti nel materiale proiettato.
\paragraph{Riferimento fonte} File: Lezione 02 - Algoritmi per la progettazione di DB Relazionali e ulteriori dipendenze.pdf. Numero slide: 45.
\section{MVD banali --- Slide 46}
\subsection*{Estratto originale dalla slide 46}
\begin{lstlisting}
MVD banali
-    Una MVD X   Y in R è detta MVD banale (trivial)
    se:
      Y è un sottoinsieme di X, oppure
      X  Y=R



-    Esempio:
    In EMP-PROJECTS,
    ENAME   PNAME è
    una MVD banale.
\end{lstlisting}
\subsection*{Spiegazione estesa}
[AGGIUNTA] Questa spiegazione collega il contenuto testuale della slide al contesto del corso. La slide enfatizza i concetti mostrati nell'estratto e il loro ruolo nella progettazione di basi di dati relazionali, con particolare attenzione a correttezza semantica, qualità dello schema e compromessi progettuali.
[AGGIUNTA] Fonte integrazione: riformulazione didattica del testo presente nella slide 46 del file Lezione 02 - Algoritmi per la progettazione di DB Relazionali e ulteriori dipendenze.pdf.
\subsection*{Derivazioni}
[AGGIUNTA] Nella slide compaiono espressioni simboliche/relazioni. Derivazione didattica generale:
\begin{align*}
X^+ &= X \\
X^+ &\leftarrow X^+ \cup Y \quad \text{se } (U \to V) \in F \text{ e } U \subseteq X^+ \text{ con } Y=V \\
\text{Stop} &\text{ quando non si aggiungono nuovi attributi in } X^+
\end{align*}
[AGGIUNTA] Questa procedura esplicita il metodo di chiusura usato nelle slide su dipendenze e chiavi.
\begin{itemize}[leftmargin=*]
\item Forma osservata nella slide: \texttt{X  Y=R}
\end{itemize}
\subsection*{Esempio numerico / esercizio}
[AGGIUNTA] Esempio: dato uno schema R(A,B,C) con dipendenze funzionali A->B e B->C, mostrare la chiusura A^+ e discutere se A è superchiave. Soluzione sintetica: A^+={A,B,C}, quindi A determina tutti gli attributi di R ed è superchiave.
\subsection*{Note su assunzioni o integrazioni}
[ASSUNZIONE] L'estrazione OCR può perdere formattazione o simboli della slide originale.
[AGGIUNTA] Le spiegazioni sono estensioni didattiche per rendere espliciti passaggi impliciti nel materiale proiettato.
\paragraph{Riferimento fonte} File: Lezione 02 - Algoritmi per la progettazione di DB Relazionali e ulteriori dipendenze.pdf. Numero slide: 46.
\section{Regole di inferenza per FD e MVD --- Slide 47}
\subsection*{Estratto originale dalla slide 47}
\begin{lstlisting}
Regole di inferenza per FD e MVD
-  Come già visto per le FD, anche per le MVD
  esiste un insieme di regole di inferenza:
-  Presi R = {A1, A2, ...,An} e X, Y, Z, W Í R:


      (IR1) (Regola riflessiva per FD):
      se X   Ê Y allora X -> Y
      (IR2) (Augmentation rule per FD):
      { X -> Y }   XZ -> YZ

      (IR3) (Transitive rule per FD) :
      { X -> Y, Y -> Z }  X -> Z
\end{lstlisting}
\subsection*{Spiegazione estesa}
[AGGIUNTA] Questa spiegazione collega il contenuto testuale della slide al contesto del corso. La slide enfatizza i concetti mostrati nell'estratto e il loro ruolo nella progettazione di basi di dati relazionali, con particolare attenzione a correttezza semantica, qualità dello schema e compromessi progettuali.
[AGGIUNTA] Fonte integrazione: riformulazione didattica del testo presente nella slide 47 del file Lezione 02 - Algoritmi per la progettazione di DB Relazionali e ulteriori dipendenze.pdf.
\subsection*{Derivazioni}
[AGGIUNTA] Nella slide compaiono espressioni simboliche/relazioni. Derivazione didattica generale:
\begin{align*}
X^+ &= X \\
X^+ &\leftarrow X^+ \cup Y \quad \text{se } (U \to V) \in F \text{ e } U \subseteq X^+ \text{ con } Y=V \\
\text{Stop} &\text{ quando non si aggiungono nuovi attributi in } X^+
\end{align*}
[AGGIUNTA] Questa procedura esplicita il metodo di chiusura usato nelle slide su dipendenze e chiavi.
\begin{itemize}[leftmargin=*]
\item Forma osservata nella slide: \texttt{-  Presi R = (A1, A2, ...,An) e X, Y, Z, W Í R:}
\item Forma osservata nella slide: \texttt{se X   Ê Y allora X -> Y}
\item Forma osservata nella slide: \texttt{( X -> Y )   XZ -> YZ}
\end{itemize}
\subsection*{Esempio numerico / esercizio}
[AGGIUNTA] Esempio: dato uno schema R(A,B,C) con dipendenze funzionali A->B e B->C, mostrare la chiusura A^+ e discutere se A è superchiave. Soluzione sintetica: A^+={A,B,C}, quindi A determina tutti gli attributi di R ed è superchiave.
\subsection*{Note su assunzioni o integrazioni}
[ASSUNZIONE] L'estrazione OCR può perdere formattazione o simboli della slide originale.
[AGGIUNTA] Le spiegazioni sono estensioni didattiche per rendere espliciti passaggi impliciti nel materiale proiettato.
\paragraph{Riferimento fonte} File: Lezione 02 - Algoritmi per la progettazione di DB Relazionali e ulteriori dipendenze.pdf. Numero slide: 47.
\section{Regole di inferenza per FD e MVD (2) --- Slide 48}
\subsection*{Estratto originale dalla slide 48}
\begin{lstlisting}
Regole di inferenza per FD e MVD (2)
   (IR4) (Complementation rule per MVD):
   { X   Y }  { X   (R- ( X   Y)) }

   (IR5) (Augmentation rule per MVD):
   se X   Y e W Ê Z allora WX   YZ

   (IR6) (Transitive rule per MVD):
   {X   Y, Y   Z }  X   (Z - Y)

   (IR7) (Replication rule da FD a MVD):
   {X -> Y}  X   Y

   IR7 is applicable with the additional implicit restriction that at
   most one value of Y is associated with each value of X:
      the set of values of Y determined by a value of X is restricted to
      being a singleton set with only one value
      In practice, this rarery happens.
\end{lstlisting}
\subsection*{Spiegazione estesa}
[AGGIUNTA] Questa spiegazione collega il contenuto testuale della slide al contesto del corso. La slide enfatizza i concetti mostrati nell'estratto e il loro ruolo nella progettazione di basi di dati relazionali, con particolare attenzione a correttezza semantica, qualità dello schema e compromessi progettuali.
[AGGIUNTA] Fonte integrazione: riformulazione didattica del testo presente nella slide 48 del file Lezione 02 - Algoritmi per la progettazione di DB Relazionali e ulteriori dipendenze.pdf.
\subsection*{Derivazioni}
[AGGIUNTA] Nella slide compaiono espressioni simboliche/relazioni. Derivazione didattica generale:
\begin{align*}
X^+ &= X \\
X^+ &\leftarrow X^+ \cup Y \quad \text{se } (U \to V) \in F \text{ e } U \subseteq X^+ \text{ con } Y=V \\
\text{Stop} &\text{ quando non si aggiungono nuovi attributi in } X^+
\end{align*}
[AGGIUNTA] Questa procedura esplicita il metodo di chiusura usato nelle slide su dipendenze e chiavi.
\begin{itemize}[leftmargin=*]
\item Forma osservata nella slide: \texttt{(X -> Y)  X   Y}
\end{itemize}
\subsection*{Esempio numerico / esercizio}
[AGGIUNTA] Esempio: dato uno schema R(A,B,C) con dipendenze funzionali A->B e B->C, mostrare la chiusura A^+ e discutere se A è superchiave. Soluzione sintetica: A^+={A,B,C}, quindi A determina tutti gli attributi di R ed è superchiave.
\subsection*{Note su assunzioni o integrazioni}
[ASSUNZIONE] L'estrazione OCR può perdere formattazione o simboli della slide originale.
[AGGIUNTA] Le spiegazioni sono estensioni didattiche per rendere espliciti passaggi impliciti nel materiale proiettato.
\paragraph{Riferimento fonte} File: Lezione 02 - Algoritmi per la progettazione di DB Relazionali e ulteriori dipendenze.pdf. Numero slide: 48.
\section{Regole di inferenza per FD e MVD (3) --- Slide 49}
\subsection*{Estratto originale dalla slide 49}
\begin{lstlisting}
Regole di inferenza per FD e MVD (3)
         (IR8) (Regola di unione per FD e MVD):
         se X -» Y e  W con le proprietà che
          1.   W   Y è vuoto
          2.   W->Z
          3.   Y ÊZ
         allora X -> Z


-    L'insieme delle regole di inferenza è corretto e
    completo:
         Dato un insieme F di FD e MVD specificate su
         R = {A1, A2, ...,An}, possiamo usare IR1 / IR8 per
         inferire l'insieme completo di tutte le dipendenze
         funzionali e multivalued.
\end{lstlisting}
\subsection*{Spiegazione estesa}
[AGGIUNTA] Questa spiegazione collega il contenuto testuale della slide al contesto del corso. La slide enfatizza i concetti mostrati nell'estratto e il loro ruolo nella progettazione di basi di dati relazionali, con particolare attenzione a correttezza semantica, qualità dello schema e compromessi progettuali.
[AGGIUNTA] Fonte integrazione: riformulazione didattica del testo presente nella slide 49 del file Lezione 02 - Algoritmi per la progettazione di DB Relazionali e ulteriori dipendenze.pdf.
\subsection*{Derivazioni}
[AGGIUNTA] Nella slide compaiono espressioni simboliche/relazioni. Derivazione didattica generale:
\begin{align*}
X^+ &= X \\
X^+ &\leftarrow X^+ \cup Y \quad \text{se } (U \to V) \in F \text{ e } U \subseteq X^+ \text{ con } Y=V \\
\text{Stop} &\text{ quando non si aggiungono nuovi attributi in } X^+
\end{align*}
[AGGIUNTA] Questa procedura esplicita il metodo di chiusura usato nelle slide su dipendenze e chiavi.
\begin{itemize}[leftmargin=*]
\item Forma osservata nella slide: \texttt{2.   W->Z}
\item Forma osservata nella slide: \texttt{allora X -> Z}
\item Forma osservata nella slide: \texttt{R = (A1, A2, ...,An), possiamo usare IR1 / IR8 per}
\end{itemize}
\subsection*{Esempio numerico / esercizio}
[AGGIUNTA] Esempio: dato uno schema R(A,B,C) con dipendenze funzionali A->B e B->C, mostrare la chiusura A^+ e discutere se A è superchiave. Soluzione sintetica: A^+={A,B,C}, quindi A determina tutti gli attributi di R ed è superchiave.
\subsection*{Note su assunzioni o integrazioni}
[ASSUNZIONE] L'estrazione OCR può perdere formattazione o simboli della slide originale.
[AGGIUNTA] Le spiegazioni sono estensioni didattiche per rendere espliciti passaggi impliciti nel materiale proiettato.
\paragraph{Riferimento fonte} File: Lezione 02 - Algoritmi per la progettazione di DB Relazionali e ulteriori dipendenze.pdf. Numero slide: 49.
\section{Quarta Forma Normale (4NF) --- Slide 50}
\subsection*{Estratto originale dalla slide 50}
\begin{lstlisting}
Quarta Forma Normale (4NF)
\end{lstlisting}
\subsection*{Spiegazione estesa}
[AGGIUNTA] Questa spiegazione collega il contenuto testuale della slide al contesto del corso. La slide enfatizza i concetti mostrati nell'estratto e il loro ruolo nella progettazione di basi di dati relazionali, con particolare attenzione a correttezza semantica, qualità dello schema e compromessi progettuali.
[AGGIUNTA] Fonte integrazione: riformulazione didattica del testo presente nella slide 50 del file Lezione 02 - Algoritmi per la progettazione di DB Relazionali e ulteriori dipendenze.pdf.
\subsection*{Derivazioni}
[AGGIUNTA] In questa slide non sono state rilevate equazioni esplicite nel testo estratto; non è quindi richiesta una derivazione algebrica puntuale.
\subsection*{Esempio numerico / esercizio}
[AGGIUNTA] Esempio: dato uno schema R(A,B,C) con dipendenze funzionali A->B e B->C, mostrare la chiusura A^+ e discutere se A è superchiave. Soluzione sintetica: A^+={A,B,C}, quindi A determina tutti gli attributi di R ed è superchiave.
\subsection*{Note su assunzioni o integrazioni}
[ASSUNZIONE] L'estrazione OCR può perdere formattazione o simboli della slide originale.
[AGGIUNTA] Le spiegazioni sono estensioni didattiche per rendere espliciti passaggi impliciti nel materiale proiettato.
\paragraph{Riferimento fonte} File: Lezione 02 - Algoritmi per la progettazione di DB Relazionali e ulteriori dipendenze.pdf. Numero slide: 50.
\section{Quarta Forma Normale --- Slide 51}
\subsection*{Estratto originale dalla slide 51}
\begin{lstlisting}
Quarta Forma Normale
-    Definizione:
    Uno schema di relazione R è detto essere in 4NF
    rispetto a un insieme di dipendenza F se, per
    ogni dipendenza multivalued non-trivial X   Y in
    F+, X è una superchiave per R.
\end{lstlisting}
\subsection*{Spiegazione estesa}
[AGGIUNTA] Questa spiegazione collega il contenuto testuale della slide al contesto del corso. La slide enfatizza i concetti mostrati nell'estratto e il loro ruolo nella progettazione di basi di dati relazionali, con particolare attenzione a correttezza semantica, qualità dello schema e compromessi progettuali.
[AGGIUNTA] Fonte integrazione: riformulazione didattica del testo presente nella slide 51 del file Lezione 02 - Algoritmi per la progettazione di DB Relazionali e ulteriori dipendenze.pdf.
\subsection*{Derivazioni}
[AGGIUNTA] In questa slide non sono state rilevate equazioni esplicite nel testo estratto; non è quindi richiesta una derivazione algebrica puntuale.
\subsection*{Esempio numerico / esercizio}
[AGGIUNTA] Esempio: dato uno schema R(A,B,C) con dipendenze funzionali A->B e B->C, mostrare la chiusura A^+ e discutere se A è superchiave. Soluzione sintetica: A^+={A,B,C}, quindi A determina tutti gli attributi di R ed è superchiave.
\subsection*{Note su assunzioni o integrazioni}
[ASSUNZIONE] L'estrazione OCR può perdere formattazione o simboli della slide originale.
[AGGIUNTA] Le spiegazioni sono estensioni didattiche per rendere espliciti passaggi impliciti nel materiale proiettato.
\paragraph{Riferimento fonte} File: Lezione 02 - Algoritmi per la progettazione di DB Relazionali e ulteriori dipendenze.pdf. Numero slide: 51.
\section{4NF: Esempio --- Slide 52}
\subsection*{Estratto originale dalla slide 52}
\begin{lstlisting}
4NF: Esempio
-    Nella relazione EMP risulta:
      ENAME   PNAME,

      ENAME   DNAME,

    ma ENAME non è una superchiave.

-    Ciò implica che EMP non è in 4NF.
\end{lstlisting}
\subsection*{Spiegazione estesa}
[AGGIUNTA] Questa spiegazione collega il contenuto testuale della slide al contesto del corso. La slide enfatizza i concetti mostrati nell'estratto e il loro ruolo nella progettazione di basi di dati relazionali, con particolare attenzione a correttezza semantica, qualità dello schema e compromessi progettuali.
[AGGIUNTA] Fonte integrazione: riformulazione didattica del testo presente nella slide 52 del file Lezione 02 - Algoritmi per la progettazione di DB Relazionali e ulteriori dipendenze.pdf.
\subsection*{Derivazioni}
[AGGIUNTA] In questa slide non sono state rilevate equazioni esplicite nel testo estratto; non è quindi richiesta una derivazione algebrica puntuale.
\subsection*{Esempio numerico / esercizio}
[AGGIUNTA] Esempio: dato uno schema R(A,B,C) con dipendenze funzionali A->B e B->C, mostrare la chiusura A^+ e discutere se A è superchiave. Soluzione sintetica: A^+={A,B,C}, quindi A determina tutti gli attributi di R ed è superchiave.
\subsection*{Note su assunzioni o integrazioni}
[ASSUNZIONE] L'estrazione OCR può perdere formattazione o simboli della slide originale.
[AGGIUNTA] Le spiegazioni sono estensioni didattiche per rendere espliciti passaggi impliciti nel materiale proiettato.
\paragraph{Riferimento fonte} File: Lezione 02 - Algoritmi per la progettazione di DB Relazionali e ulteriori dipendenze.pdf. Numero slide: 52.
\section{4NF: Esempio (2) --- Slide 53}
\subsection*{Estratto originale dalla slide 53}
\begin{lstlisting}
4NF: Esempio (2)




  -    EMP viene decomposta in EMP_PROJECTS e
      EMP_DEPENDENTS, che sono in 4NF:
        ENAME   PNAME è una MVD in EMP_PROJECTS

        ENAME   DNAME è una MVD banale in
       EMP_DEPENDENTS.
\end{lstlisting}
\subsection*{Spiegazione estesa}
[AGGIUNTA] Questa spiegazione collega il contenuto testuale della slide al contesto del corso. La slide enfatizza i concetti mostrati nell'estratto e il loro ruolo nella progettazione di basi di dati relazionali, con particolare attenzione a correttezza semantica, qualità dello schema e compromessi progettuali.
[AGGIUNTA] Fonte integrazione: riformulazione didattica del testo presente nella slide 53 del file Lezione 02 - Algoritmi per la progettazione di DB Relazionali e ulteriori dipendenze.pdf.
\subsection*{Derivazioni}
[AGGIUNTA] In questa slide non sono state rilevate equazioni esplicite nel testo estratto; non è quindi richiesta una derivazione algebrica puntuale.
\subsection*{Esempio numerico / esercizio}
[AGGIUNTA] Esempio: dato uno schema R(A,B,C) con dipendenze funzionali A->B e B->C, mostrare la chiusura A^+ e discutere se A è superchiave. Soluzione sintetica: A^+={A,B,C}, quindi A determina tutti gli attributi di R ed è superchiave.
\subsection*{Note su assunzioni o integrazioni}
[ASSUNZIONE] L'estrazione OCR può perdere formattazione o simboli della slide originale.
[AGGIUNTA] Le spiegazioni sono estensioni didattiche per rendere espliciti passaggi impliciti nel materiale proiettato.
\paragraph{Riferimento fonte} File: Lezione 02 - Algoritmi per la progettazione di DB Relazionali e ulteriori dipendenze.pdf. Numero slide: 53.
\section{Nota sulle MVD non-trivial --- Slide 54}
\subsection*{Estratto originale dalla slide 54}
\begin{lstlisting}
Nota sulle MVD non-trivial
-    Relazioni contenenti MVD non-trivial tendono ad
    essere relazioni "tutta chiave", nel senso che la
    chiave è formata da tutti gli attributi.
\end{lstlisting}
\subsection*{Spiegazione estesa}
[AGGIUNTA] Questa spiegazione collega il contenuto testuale della slide al contesto del corso. La slide enfatizza i concetti mostrati nell'estratto e il loro ruolo nella progettazione di basi di dati relazionali, con particolare attenzione a correttezza semantica, qualità dello schema e compromessi progettuali.
[AGGIUNTA] Fonte integrazione: riformulazione didattica del testo presente nella slide 54 del file Lezione 02 - Algoritmi per la progettazione di DB Relazionali e ulteriori dipendenze.pdf.
\subsection*{Derivazioni}
[AGGIUNTA] In questa slide non sono state rilevate equazioni esplicite nel testo estratto; non è quindi richiesta una derivazione algebrica puntuale.
\subsection*{Esempio numerico / esercizio}
[AGGIUNTA] Esempio: dato uno schema R(A,B,C) con dipendenze funzionali A->B e B->C, mostrare la chiusura A^+ e discutere se A è superchiave. Soluzione sintetica: A^+={A,B,C}, quindi A determina tutti gli attributi di R ed è superchiave.
\subsection*{Note su assunzioni o integrazioni}
[ASSUNZIONE] L'estrazione OCR può perdere formattazione o simboli della slide originale.
[AGGIUNTA] Le spiegazioni sono estensioni didattiche per rendere espliciti passaggi impliciti nel materiale proiettato.
\paragraph{Riferimento fonte} File: Lezione 02 - Algoritmi per la progettazione di DB Relazionali e ulteriori dipendenze.pdf. Numero slide: 54.
\section{Decomposizione LJ in Relazioni 4NF --- Slide 55}
\subsection*{Estratto originale dalla slide 55}
\begin{lstlisting}
Decomposizione LJ in Relazioni 4NF
-  Decomponendo uno schema R in R 1= (X   Y) e
  R2= (R - Y) in base a una MVD X   Y, la
  decomposizione gode della proprietà LJ
  (condizione necessaria e sufficiente).
-  Proprietà LJ1': Gli schemi di relazione R1e R2
  formano una decomposizione LJ di R se e solo
  se:
      (R1   R2)   (R1- R2)
      oppure, per simmetria se e solo se
      (R1   R2)   (R2- R1)
\end{lstlisting}
\subsection*{Spiegazione estesa}
[AGGIUNTA] Questa spiegazione collega il contenuto testuale della slide al contesto del corso. La slide enfatizza i concetti mostrati nell'estratto e il loro ruolo nella progettazione di basi di dati relazionali, con particolare attenzione a correttezza semantica, qualità dello schema e compromessi progettuali.
[AGGIUNTA] Fonte integrazione: riformulazione didattica del testo presente nella slide 55 del file Lezione 02 - Algoritmi per la progettazione di DB Relazionali e ulteriori dipendenze.pdf.
\subsection*{Derivazioni}
[AGGIUNTA] Nella slide compaiono espressioni simboliche/relazioni. Derivazione didattica generale:
\begin{align*}
X^+ &= X \\
X^+ &\leftarrow X^+ \cup Y \quad \text{se } (U \to V) \in F \text{ e } U \subseteq X^+ \text{ con } Y=V \\
\text{Stop} &\text{ quando non si aggiungono nuovi attributi in } X^+
\end{align*}
[AGGIUNTA] Questa procedura esplicita il metodo di chiusura usato nelle slide su dipendenze e chiavi.
\begin{itemize}[leftmargin=*]
\item Forma osservata nella slide: \texttt{-  Decomponendo uno schema R in R 1= (X   Y) e}
\item Forma osservata nella slide: \texttt{R2= (R - Y) in base a una MVD X   Y, la}
\end{itemize}
\subsection*{Esempio numerico / esercizio}
[AGGIUNTA] Esempio: dato uno schema R(A,B,C) con dipendenze funzionali A->B e B->C, mostrare la chiusura A^+ e discutere se A è superchiave. Soluzione sintetica: A^+={A,B,C}, quindi A determina tutti gli attributi di R ed è superchiave.
\subsection*{Note su assunzioni o integrazioni}
[ASSUNZIONE] L'estrazione OCR può perdere formattazione o simboli della slide originale.
[AGGIUNTA] Le spiegazioni sono estensioni didattiche per rendere espliciti passaggi impliciti nel materiale proiettato.
\paragraph{Riferimento fonte} File: Lezione 02 - Algoritmi per la progettazione di DB Relazionali e ulteriori dipendenze.pdf. Numero slide: 55.
\section{Decomposizione LJ in relazioni 4NF --- Slide 56}
\subsection*{Estratto originale dalla slide 56}
\begin{lstlisting}
Decomposizione LJ in relazioni 4NF
-     Algoritmo 5
     Decomposizione di un schema in relazioni 4NF
     con la proprietà LJ:
    1.   porre D = {R};
    2.   finché esiste uno schema di relazione Q in D che non è
         in 4NF
         {
             scegliere uno schema di relazione Q in D che non è in 4NF;
             trovare una MVD non banale X   Y in Q che viola 4NF;
             rimpiazzare Q in D con due schemi (Q - Y) e (X   Y);
         }
\end{lstlisting}
\subsection*{Spiegazione estesa}
[AGGIUNTA] Questa spiegazione collega il contenuto testuale della slide al contesto del corso. La slide enfatizza i concetti mostrati nell'estratto e il loro ruolo nella progettazione di basi di dati relazionali, con particolare attenzione a correttezza semantica, qualità dello schema e compromessi progettuali.
[AGGIUNTA] Fonte integrazione: riformulazione didattica del testo presente nella slide 56 del file Lezione 02 - Algoritmi per la progettazione di DB Relazionali e ulteriori dipendenze.pdf.
\subsection*{Derivazioni}
[AGGIUNTA] Nella slide compaiono espressioni simboliche/relazioni. Derivazione didattica generale:
\begin{align*}
X^+ &= X \\
X^+ &\leftarrow X^+ \cup Y \quad \text{se } (U \to V) \in F \text{ e } U \subseteq X^+ \text{ con } Y=V \\
\text{Stop} &\text{ quando non si aggiungono nuovi attributi in } X^+
\end{align*}
[AGGIUNTA] Questa procedura esplicita il metodo di chiusura usato nelle slide su dipendenze e chiavi.
\begin{itemize}[leftmargin=*]
\item Forma osservata nella slide: \texttt{1.   porre D = (R);}
\end{itemize}
\subsection*{Esempio numerico / esercizio}
[AGGIUNTA] Esempio: dato uno schema R(A,B,C) con dipendenze funzionali A->B e B->C, mostrare la chiusura A^+ e discutere se A è superchiave. Soluzione sintetica: A^+={A,B,C}, quindi A determina tutti gli attributi di R ed è superchiave.
\subsection*{Note su assunzioni o integrazioni}
[ASSUNZIONE] L'estrazione OCR può perdere formattazione o simboli della slide originale.
[AGGIUNTA] Le spiegazioni sono estensioni didattiche per rendere espliciti passaggi impliciti nel materiale proiettato.
\paragraph{Riferimento fonte} File: Lezione 02 - Algoritmi per la progettazione di DB Relazionali e ulteriori dipendenze.pdf. Numero slide: 56.
\section{Conclusione}
[AGGIUNTA] Punti chiave: (i) la qualità degli schemi relazionali si valuta tramite ridondanza, anomalie e proprietà formali; (ii) dipendenze funzionali e chiusure guidano analisi e trasformazioni; (iii) normalizzazione (3NF/BCNF) e algoritmi top-down/bottom-up aiutano a ottenere schemi robusti; (iv) la progettazione richiede equilibrio tra rigore formale e vincoli pratici.
\appendix
\section{Trascrizioni complete delle slide}
[AGGIUNTA] Le trascrizioni sono incluse nelle sezioni principali slide-per-slide; questa appendice funge da richiamo strutturale.
\section{Dati usati per i grafici}
\begin{table}[ht]
\centering
\caption{Dati sintetici usati nelle Figure \ref{fig:crescita-combinatoria}, \ref{fig:confronto-strategie}, \ref{fig:miglioramento-normalizzazione}}
\label{tab:dati-grafici-riepilogo}
\begin{tabular}{@{}lll@{}}
\toprule
Figura & Tipo dati & Motivazione \\
\midrule
\texttt{fig:crescita-combinatoria} & sintetici & visualizzare andamento esponenziale \\
\texttt{fig:confronto-strategie} & sintetici & confronto qualitativo tra approcci \\
\texttt{fig:miglioramento-normalizzazione} & sintetici & mostrare miglioramento progressivo \\
\bottomrule
\end{tabular}
\end{table}
\section{Mappa slide -> sezione}
\begin{longtable}{@{}p{0.60\textwidth}cc@{}}
\toprule
File sorgente & Slide & Sezione \\
\midrule
\endfirsthead
\toprule
File sorgente & Slide & Sezione \\
\midrule
\endhead
Lezione 00 - Presentazione del corso.pdf & 1 & CORSO DI LAUREA MAGISTRALE \\
Lezione 00 - Presentazione del corso.pdf & 2 & Caratteristiche generali \\
Lezione 00 - Presentazione del corso.pdf & 3 & Calendario del corso \\
Lezione 00 - Presentazione del corso.pdf & 4 & Laboratorio \\
Lezione 00 - Presentazione del corso.pdf & 5 & Descrizione del corso \\
Lezione 00 - Presentazione del corso.pdf & 6 & Obiettivi \\
Lezione 00 - Presentazione del corso.pdf & 7 & Prerequisiti \\
Lezione 00 - Presentazione del corso.pdf & 8 & Link \\
Lezione 00 - Presentazione del corso.pdf & 9 & Modalità d'esame \\
Lezione 00 - Presentazione del corso.pdf & 10 & Testi consigliati \\
Lezione 01 - Dipendenze funzionali e Normalizzazione di DB Relazionali.pdf & 1 & BASI DI DATI 2 \\
Lezione 01 - Dipendenze funzionali e Normalizzazione di DB Relazionali.pdf & 2 & Le fasi di progettazione di un Database \\
Lezione 01 - Dipendenze funzionali e Normalizzazione di DB Relazionali.pdf & 3 & Qualità di schemi relazionali \\
Lezione 01 - Dipendenze funzionali e Normalizzazione di DB Relazionali.pdf & 4 & Qualità di schemi relazionali (2) \\
Lezione 01 - Dipendenze funzionali e Normalizzazione di DB Relazionali.pdf & 5 & Misure informali di qualità \\
Lezione 01 - Dipendenze funzionali e Normalizzazione di DB Relazionali.pdf & 6 & Modello ER Company \\
Lezione 01 - Dipendenze funzionali e Normalizzazione di DB Relazionali.pdf & 7 & Outline \\
Lezione 01 - Dipendenze funzionali e Normalizzazione di DB Relazionali.pdf & 8 & 1) Semantica degli attributi di una \\
Lezione 01 - Dipendenze funzionali e Normalizzazione di DB Relazionali.pdf & 9 & 1) Semantica degli attributi di una relazione: \\
Lezione 01 - Dipendenze funzionali e Normalizzazione di DB Relazionali.pdf & 10 & 1) Semantica degli attributi di una relazione: \\
Lezione 01 - Dipendenze funzionali e Normalizzazione di DB Relazionali.pdf & 11 & 1) Semantica degli attributi di una relazione: \\
Lezione 01 - Dipendenze funzionali e Normalizzazione di DB Relazionali.pdf & 12 & 1) Semantica degli attributi di una relazione: \\
Lezione 01 - Dipendenze funzionali e Normalizzazione di DB Relazionali.pdf & 13 & 1) Semantica degli attributi di una relazione: \\
Lezione 01 - Dipendenze funzionali e Normalizzazione di DB Relazionali.pdf & 14 & Linea guida 1 \\
Lezione 01 - Dipendenze funzionali e Normalizzazione di DB Relazionali.pdf & 15 & Linea guida 1: Esempio \\
Lezione 01 - Dipendenze funzionali e Normalizzazione di DB Relazionali.pdf & 16 & Outline \\
Lezione 01 - Dipendenze funzionali e Normalizzazione di DB Relazionali.pdf & 17 & 2) Riduzione dei valori ridondanti nelle tuple \\
Lezione 01 - Dipendenze funzionali e Normalizzazione di DB Relazionali.pdf & 18 & 2) Riduzione dei valori ridondanti nelle \\
Lezione 01 - Dipendenze funzionali e Normalizzazione di DB Relazionali.pdf & 19 & 2) Riduzione dei valori ridondanti nelle \\
Lezione 01 - Dipendenze funzionali e Normalizzazione di DB Relazionali.pdf & 20 & 2) Riduzione dei valori ridondanti nelle \\
Lezione 01 - Dipendenze funzionali e Normalizzazione di DB Relazionali.pdf & 21 & Update Anomalies \\
Lezione 01 - Dipendenze funzionali e Normalizzazione di DB Relazionali.pdf & 22 & Insertion Anomalies \\
Lezione 01 - Dipendenze funzionali e Normalizzazione di DB Relazionali.pdf & 23 & Insertion Anomalies (2) \\
Lezione 01 - Dipendenze funzionali e Normalizzazione di DB Relazionali.pdf & 24 & Deletion Anomalies \\
Lezione 01 - Dipendenze funzionali e Normalizzazione di DB Relazionali.pdf & 25 & Modification Anomalies \\
Lezione 01 - Dipendenze funzionali e Normalizzazione di DB Relazionali.pdf & 26 & Linea guida 2 \\
Lezione 01 - Dipendenze funzionali e Normalizzazione di DB Relazionali.pdf & 27 & Outline \\
Lezione 01 - Dipendenze funzionali e Normalizzazione di DB Relazionali.pdf & 28 & 3) Riduzione dei valori null nelle tuple \\
Lezione 01 - Dipendenze funzionali e Normalizzazione di DB Relazionali.pdf & 29 & Interpretazioni di null \\
Lezione 01 - Dipendenze funzionali e Normalizzazione di DB Relazionali.pdf & 30 & Linea guida 3 \\
Lezione 01 - Dipendenze funzionali e Normalizzazione di DB Relazionali.pdf & 31 & Linea guida 3: Esempio \\
Lezione 01 - Dipendenze funzionali e Normalizzazione di DB Relazionali.pdf & 32 & Outline \\
Lezione 01 - Dipendenze funzionali e Normalizzazione di DB Relazionali.pdf & 33 & Tuple spurie \\
Lezione 01 - Dipendenze funzionali e Normalizzazione di DB Relazionali.pdf & 34 & Tuple spurie (2) \\
Lezione 01 - Dipendenze funzionali e Normalizzazione di DB Relazionali.pdf & 35 & Tuple spurie: Esempio \\
Lezione 01 - Dipendenze funzionali e Normalizzazione di DB Relazionali.pdf & 36 & Tuple spurie: Esempio (2) \\
Lezione 01 - Dipendenze funzionali e Normalizzazione di DB Relazionali.pdf & 37 & Motivo delle tuple spurie \\
Lezione 01 - Dipendenze funzionali e Normalizzazione di DB Relazionali.pdf & 38 & Linea guida 4 \\
Lezione 01 - Dipendenze funzionali e Normalizzazione di DB Relazionali.pdf & 39 & Sommario delle linee guida \\
Lezione 01 - Dipendenze funzionali e Normalizzazione di DB Relazionali.pdf & 40 & Dipendenze Funzionali \\
Lezione 01 - Dipendenze funzionali e Normalizzazione di DB Relazionali.pdf & 41 & Qualità di schemi relazionali \\
Lezione 01 - Dipendenze funzionali e Normalizzazione di DB Relazionali.pdf & 42 & Dipendenza Funzionale \\
Lezione 01 - Dipendenze funzionali e Normalizzazione di DB Relazionali.pdf & 43 & Dipendenza Funzionale (2) \\
Lezione 01 - Dipendenze funzionali e Normalizzazione di DB Relazionali.pdf & 44 & Dipendenza Funzionale (3) \\
Lezione 01 - Dipendenze funzionali e Normalizzazione di DB Relazionali.pdf & 45 & Dipendenza Funzionale (4) \\
Lezione 01 - Dipendenze funzionali e Normalizzazione di DB Relazionali.pdf & 46 & Dipendenza Funzionale: Esempio \\
Lezione 01 - Dipendenze funzionali e Normalizzazione di DB Relazionali.pdf & 47 & Dipendenza funzionale dalla semantica \\
Lezione 01 - Dipendenze funzionali e Normalizzazione di DB Relazionali.pdf & 48 & Dipendenza Funzionale: Esempio \\
Lezione 01 - Dipendenze funzionali e Normalizzazione di DB Relazionali.pdf & 49 & Regole di inferenza per le FD \\
Lezione 01 - Dipendenze funzionali e Normalizzazione di DB Relazionali.pdf & 50 & Regole di inferenza: Esempio \\
Lezione 01 - Dipendenze funzionali e Normalizzazione di DB Relazionali.pdf & 51 & Regole di inferenza per le FD (2) \\
Lezione 01 - Dipendenze funzionali e Normalizzazione di DB Relazionali.pdf & 52 & Regole di inferenza fondamentali \\
Lezione 01 - Dipendenze funzionali e Normalizzazione di DB Relazionali.pdf & 53 & Regole di inferenza fondamentali (2) \\
Lezione 01 - Dipendenze funzionali e Normalizzazione di DB Relazionali.pdf & 54 & Dimostrazione di IR1 \\
Lezione 01 - Dipendenze funzionali e Normalizzazione di DB Relazionali.pdf & 55 & Dimostrazione di IR2 \\
Lezione 01 - Dipendenze funzionali e Normalizzazione di DB Relazionali.pdf & 56 & Dimostrazione di IR3 \\
Lezione 01 - Dipendenze funzionali e Normalizzazione di DB Relazionali.pdf & 57 & Dimostrazione di IR4 \\
Lezione 01 - Dipendenze funzionali e Normalizzazione di DB Relazionali.pdf & 58 & Dimostrazione di IR5 \\
Lezione 01 - Dipendenze funzionali e Normalizzazione di DB Relazionali.pdf & 59 & Dimostrazione di IR6 \\
Lezione 01 - Dipendenze funzionali e Normalizzazione di DB Relazionali.pdf & 60 & Regole di inferenza di Armstrong \\
Lezione 01 - Dipendenze funzionali e Normalizzazione di DB Relazionali.pdf & 61 & Regole di inferenza nella progettazione dei db \\
Lezione 01 - Dipendenze funzionali e Normalizzazione di DB Relazionali.pdf & 62 & Algoritmo per il calcolo di X+ \\
Lezione 01 - Dipendenze funzionali e Normalizzazione di DB Relazionali.pdf & 63 & Algoritmo per il calcolo di X+: Esempio \\
Lezione 01 - Dipendenze funzionali e Normalizzazione di DB Relazionali.pdf & 64 & Equivalenza di insiemi di dipendenze \\
Lezione 01 - Dipendenze funzionali e Normalizzazione di DB Relazionali.pdf & 65 & Forme Normali \\
Lezione 01 - Dipendenze funzionali e Normalizzazione di DB Relazionali.pdf & 66 & Normalizzazione dei dati \\
Lezione 01 - Dipendenze funzionali e Normalizzazione di DB Relazionali.pdf & 67 & Forme normali \\
Lezione 01 - Dipendenze funzionali e Normalizzazione di DB Relazionali.pdf & 68 & Forme normali (2) \\
Lezione 01 - Dipendenze funzionali e Normalizzazione di DB Relazionali.pdf & 69 & Definizioni \\
Lezione 01 - Dipendenze funzionali e Normalizzazione di DB Relazionali.pdf & 70 & Definizioni (2) \\
Lezione 01 - Dipendenze funzionali e Normalizzazione di DB Relazionali.pdf & 71 & Definizioni: Esempio \\
Lezione 01 - Dipendenze funzionali e Normalizzazione di DB Relazionali.pdf & 72 & Esempio di attributi primi \\
Lezione 01 - Dipendenze funzionali e Normalizzazione di DB Relazionali.pdf & 73 & Prima forma normale (1NF) \\
Lezione 01 - Dipendenze funzionali e Normalizzazione di DB Relazionali.pdf & 74 & Prima forma normale: Esempio \\
Lezione 01 - Dipendenze funzionali e Normalizzazione di DB Relazionali.pdf & 75 & Prima forma normale: Esempio (2) \\
Lezione 01 - Dipendenze funzionali e Normalizzazione di DB Relazionali.pdf & 76 & Prima forma normale: Esempio (3) \\
Lezione 01 - Dipendenze funzionali e Normalizzazione di DB Relazionali.pdf & 77 & Prima forma normale: Esempio 2 \\
Lezione 01 - Dipendenze funzionali e Normalizzazione di DB Relazionali.pdf & 78 & Prima forma normale: Esempio 2 (2) \\
Lezione 01 - Dipendenze funzionali e Normalizzazione di DB Relazionali.pdf & 79 & Seconda Forma Normale (2NF) \\
Lezione 01 - Dipendenze funzionali e Normalizzazione di DB Relazionali.pdf & 80 & Esempio dipendenze piene \\
Lezione 01 - Dipendenze funzionali e Normalizzazione di DB Relazionali.pdf & 81 & Seconda Forma Normale (2NF) \\
Lezione 01 - Dipendenze funzionali e Normalizzazione di DB Relazionali.pdf & 82 & Seconda Forma Normale: Esempio \\
Lezione 01 - Dipendenze funzionali e Normalizzazione di DB Relazionali.pdf & 83 & Seconda Forma Normale \\
Lezione 01 - Dipendenze funzionali e Normalizzazione di DB Relazionali.pdf & 84 & Esempio di 2NF \\
Lezione 01 - Dipendenze funzionali e Normalizzazione di DB Relazionali.pdf & 85 & Terza Forma Normale (3NF) \\
Lezione 01 - Dipendenze funzionali e Normalizzazione di DB Relazionali.pdf & 86 & Terza Forma Normale: Esempio \\
Lezione 01 - Dipendenze funzionali e Normalizzazione di DB Relazionali.pdf & 87 & Terza Forma Normale: Esempio (2) \\
Lezione 01 - Dipendenze funzionali e Normalizzazione di DB Relazionali.pdf & 88 & Terza Forma Normale \\
Lezione 01 - Dipendenze funzionali e Normalizzazione di DB Relazionali.pdf & 89 & Sommario \\
Lezione 01 - Dipendenze funzionali e Normalizzazione di DB Relazionali.pdf & 90 & Slide 90 \\
Lezione 01 - Dipendenze funzionali e Normalizzazione di DB Relazionali.pdf & 91 & Definizione generale di Seconda e Terza \\
Lezione 01 - Dipendenze funzionali e Normalizzazione di DB Relazionali.pdf & 92 & Forma Normale di Boyce-Codd (BCNF) \\
Lezione 01 - Dipendenze funzionali e Normalizzazione di DB Relazionali.pdf & 93 & Forme Normali: Esempio \\
Lezione 01 - Dipendenze funzionali e Normalizzazione di DB Relazionali.pdf & 94 & Forme Normali: Esempio (2) \\
Lezione 01 - Dipendenze funzionali e Normalizzazione di DB Relazionali.pdf & 95 & Forme Normali: Esempio (3) \\
Lezione 01 - Dipendenze funzionali e Normalizzazione di DB Relazionali.pdf & 96 & Forme Normali: Esempio (4) \\
Lezione 01 - Dipendenze funzionali e Normalizzazione di DB Relazionali.pdf & 97 & Forme Normali: Esempio (5) \\
Lezione 01 - Dipendenze funzionali e Normalizzazione di DB Relazionali.pdf & 98 & Forme Normali: Esempio (6) \\
Lezione 02 - Algoritmi per la progettazione di DB Relazionali e ulteriori dipendenze.pdf & 1 & BASI DI DATI 2 \\
Lezione 02 - Algoritmi per la progettazione di DB Relazionali e ulteriori dipendenze.pdf & 2 & Top-Down o Bottom-Up? \\
Lezione 02 - Algoritmi per la progettazione di DB Relazionali e ulteriori dipendenze.pdf & 3 & Algoritmi per la progettazione di \\
Lezione 02 - Algoritmi per la progettazione di DB Relazionali e ulteriori dipendenze.pdf & 4 & Decomposizione delle relazioni \\
Lezione 02 - Algoritmi per la progettazione di DB Relazionali e ulteriori dipendenze.pdf & 5 & Conservazione degli attributi \\
Lezione 02 - Algoritmi per la progettazione di DB Relazionali e ulteriori dipendenze.pdf & 6 & Decomposizione delle relazioni (2) \\
Lezione 02 - Algoritmi per la progettazione di DB Relazionali e ulteriori dipendenze.pdf & 7 & Conservazione delle dipendenze \\
Lezione 02 - Algoritmi per la progettazione di DB Relazionali e ulteriori dipendenze.pdf & 8 & Conservazione delle dipendenze (2) \\
Lezione 02 - Algoritmi per la progettazione di DB Relazionali e ulteriori dipendenze.pdf & 9 & Proiezione \\
Lezione 02 - Algoritmi per la progettazione di DB Relazionali e ulteriori dipendenze.pdf & 10 & Decomposizione dependency-preserving \\
Lezione 02 - Algoritmi per la progettazione di DB Relazionali e ulteriori dipendenze.pdf & 11 & Decomposizione dependency-preserving: \\
Lezione 02 - Algoritmi per la progettazione di DB Relazionali e ulteriori dipendenze.pdf & 12 & Decomposizione dependency-preserving: \\
Lezione 02 - Algoritmi per la progettazione di DB Relazionali e ulteriori dipendenze.pdf & 13 & Proposizione 1 \\
Lezione 02 - Algoritmi per la progettazione di DB Relazionali e ulteriori dipendenze.pdf & 14 & Copertura minimale \\
Lezione 02 - Algoritmi per la progettazione di DB Relazionali e ulteriori dipendenze.pdf & 15 & Algoritmo di sintesi relazionale \\
Lezione 02 - Algoritmi per la progettazione di DB Relazionali e ulteriori dipendenze.pdf & 16 & Algoritmo per trovare una copertura \\
Lezione 02 - Algoritmi per la progettazione di DB Relazionali e ulteriori dipendenze.pdf & 17 & Algoritmo per trovare una copertura \\
Lezione 02 - Algoritmi per la progettazione di DB Relazionali e ulteriori dipendenze.pdf & 18 & Decomposizione e Join Lossless \\
Lezione 02 - Algoritmi per la progettazione di DB Relazionali e ulteriori dipendenze.pdf & 19 & Lossless Join \\
Lezione 02 - Algoritmi per la progettazione di DB Relazionali e ulteriori dipendenze.pdf & 20 & Test della proprietà di lossless join \\
Lezione 02 - Algoritmi per la progettazione di DB Relazionali e ulteriori dipendenze.pdf & 21 & Test della proprietà di lossless join (2) \\
Lezione 02 - Algoritmi per la progettazione di DB Relazionali e ulteriori dipendenze.pdf & 22 & Test della proprietà di lossless join (3) \\
Lezione 02 - Algoritmi per la progettazione di DB Relazionali e ulteriori dipendenze.pdf & 23 & Descrizione dell'algoritmo 2 \\
Lezione 02 - Algoritmi per la progettazione di DB Relazionali e ulteriori dipendenze.pdf & 24 & Descrizione dell'algoritmo 2 (2) \\
Lezione 02 - Algoritmi per la progettazione di DB Relazionali e ulteriori dipendenze.pdf & 25 & Descrizione dell'algoritmo 2 (3) \\
Lezione 02 - Algoritmi per la progettazione di DB Relazionali e ulteriori dipendenze.pdf & 26 & Algoritmo 2: Esempio \\
Lezione 02 - Algoritmi per la progettazione di DB Relazionali e ulteriori dipendenze.pdf & 27 & Trovare una decomposizione \\
Lezione 02 - Algoritmi per la progettazione di DB Relazionali e ulteriori dipendenze.pdf & 28 & Proprietà delle decomposizioni lossless join \\
Lezione 02 - Algoritmi per la progettazione di DB Relazionali e ulteriori dipendenze.pdf & 29 & Proprietà delle decomposizioni lossless join (2) \\
Lezione 02 - Algoritmi per la progettazione di DB Relazionali e ulteriori dipendenze.pdf & 30 & Decomposizione LJ in schemi in BCNF \\
Lezione 02 - Algoritmi per la progettazione di DB Relazionali e ulteriori dipendenze.pdf & 31 & Decomposizione LJ e dependency-preserving \\
Lezione 02 - Algoritmi per la progettazione di DB Relazionali e ulteriori dipendenze.pdf & 32 & Trovare una chiave K per uno schema di \\
Lezione 02 - Algoritmi per la progettazione di DB Relazionali e ulteriori dipendenze.pdf & 33 & Nota \\
Lezione 02 - Algoritmi per la progettazione di DB Relazionali e ulteriori dipendenze.pdf & 34 & Decomposizioni e valori null \\
Lezione 02 - Algoritmi per la progettazione di DB Relazionali e ulteriori dipendenze.pdf & 35 & Decomposizioni e valori null: Esempio \\
Lezione 02 - Algoritmi per la progettazione di DB Relazionali e ulteriori dipendenze.pdf & 36 & Decomposizioni e valori null: Esempio \\
Lezione 02 - Algoritmi per la progettazione di DB Relazionali e ulteriori dipendenze.pdf & 37 & Decomposizioni e valori null: Esempio \\
Lezione 02 - Algoritmi per la progettazione di DB Relazionali e ulteriori dipendenze.pdf & 38 & Decomposizioni e valori null (2) \\
Lezione 02 - Algoritmi per la progettazione di DB Relazionali e ulteriori dipendenze.pdf & 39 & Algoritmi di normalizzazione \\
Lezione 02 - Algoritmi per la progettazione di DB Relazionali e ulteriori dipendenze.pdf & 40 & Dipendenze Multivalore \\
Lezione 02 - Algoritmi per la progettazione di DB Relazionali e ulteriori dipendenze.pdf & 41 & Dipendenze Multivalore (MVD) \\
Lezione 02 - Algoritmi per la progettazione di DB Relazionali e ulteriori dipendenze.pdf & 42 & Dipendenze Multivalore: Esempio \\
Lezione 02 - Algoritmi per la progettazione di DB Relazionali e ulteriori dipendenze.pdf & 43 & Definizione formale di MVD: \\
Lezione 02 - Algoritmi per la progettazione di DB Relazionali e ulteriori dipendenze.pdf & 44 & Definizione formale di MVD (2) \\
Lezione 02 - Algoritmi per la progettazione di DB Relazionali e ulteriori dipendenze.pdf & 45 & MVD: Esempio \\
Lezione 02 - Algoritmi per la progettazione di DB Relazionali e ulteriori dipendenze.pdf & 46 & MVD banali \\
Lezione 02 - Algoritmi per la progettazione di DB Relazionali e ulteriori dipendenze.pdf & 47 & Regole di inferenza per FD e MVD \\
Lezione 02 - Algoritmi per la progettazione di DB Relazionali e ulteriori dipendenze.pdf & 48 & Regole di inferenza per FD e MVD (2) \\
Lezione 02 - Algoritmi per la progettazione di DB Relazionali e ulteriori dipendenze.pdf & 49 & Regole di inferenza per FD e MVD (3) \\
Lezione 02 - Algoritmi per la progettazione di DB Relazionali e ulteriori dipendenze.pdf & 50 & Quarta Forma Normale (4NF) \\
Lezione 02 - Algoritmi per la progettazione di DB Relazionali e ulteriori dipendenze.pdf & 51 & Quarta Forma Normale \\
Lezione 02 - Algoritmi per la progettazione di DB Relazionali e ulteriori dipendenze.pdf & 52 & 4NF: Esempio \\
Lezione 02 - Algoritmi per la progettazione di DB Relazionali e ulteriori dipendenze.pdf & 53 & 4NF: Esempio (2) \\
Lezione 02 - Algoritmi per la progettazione di DB Relazionali e ulteriori dipendenze.pdf & 54 & Nota sulle MVD non-trivial \\
Lezione 02 - Algoritmi per la progettazione di DB Relazionali e ulteriori dipendenze.pdf & 55 & Decomposizione LJ in Relazioni 4NF \\
Lezione 02 - Algoritmi per la progettazione di DB Relazionali e ulteriori dipendenze.pdf & 56 & Decomposizione LJ in relazioni 4NF \\
\bottomrule
\end{longtable}
\section{Elenco dei simboli e notazione}
\begin{tabular}{@{}ll@{}}
\toprule
Simbolo & Significato \\
\midrule
$R$ & schema di relazione \\
$A,B,C,\dots$ & attributi \\
$X \to Y$ & dipendenza funzionale \\
$X^+$ & chiusura di un insieme di attributi \\
3NF & terza forma normale \\
BCNF & Boyce-Codd Normal Form \\
\bottomrule
\end{tabular}
\section{Elenco contenuti aggiunti}
\begin{enumerate}[leftmargin=*]
\item [AGGIUNTA] Spiegazioni didattiche per ogni slide: per rendere espliciti passaggi impliciti.
\item [AGGIUNTA] Tre grafici vettoriali TikZ/PGFPlots: per chiarire trend e trade-off concettuali.
\item [AGGIUNTA] Esempio numerico/esercizio ricorrente: per fissare metodo operativo su chiusure e chiavi.
\item [AGGIUNTA] Tabelle dati sintetici: per trasparenza sui valori usati nei grafici non presenti nelle slide.
\end{enumerate}
\section{Bibliografia}
Non sono presenti riferimenti bibliografici espliciti nelle tre slide analizzate. [INCERTO] Se presenti come immagini non OCR, non sono stati estratti automaticamente.
\end{document}